\section{Splitting the critical set by gradedness}\label{sec:split}

Every maximizer is critical, so it suffices to control the critical set. The following inequality splits it in two.

\subsection{The compression inequality and its equality class}\label{sec:compression}

\begin{Lem}{7}[A bound at every critical point, and its equality class]
At every critical point,
\[
\theta\tr(W_-)+\tr(W_-^2)\;\le\;\|\Ric(W_-)\|^2+6\,\delta^2 ,
\]
with equality precisely when $\ran W_+\oplus\ran W_-$ is a $\mathbb Z_2$-graded subalgebra with even part $\ran W_+$, which is the condition defining a compact symmetric pair with $E_+$ the isotropy algebra \emph{when} $\ker W=0$, and is strictly weaker otherwise: at $W_{\CP^2}$ the graded condition holds while $\dim E_+=1$ is the even part of no involution of $\so(4)$.
\end{Lem}

\begin{proof}
Write $E_\pm$ for the positive and negative eigenspaces of $W$, $P_\pm$ for the orthogonal projections, and $W_\pm\succeq0$ for the spectral parts, so that
\[
W=W_+-W_-,\qquad |W|=W_++W_-,\qquad W_+W_-=0 .
\]

\emph{Step 1: compress criticality onto $E_-$.} Apply $P_-(\cdot)P_-$ to $Q(W)=\theta W$ and take the trace. On the right, $P_-WP_-=-W_-$, so the right side is $-\theta\tr(W_-)$. On the left, $P_-W^2P_-=W_-^2$, and therefore
\begin{equation}\label{eq:gc-compress}
\theta\tr(W_-)+\tr(W_-^2)\;=\;-\tr\big(P_-W^{\#}P_-\big).
\end{equation}
This is an identity at every critical point, and it is the only place criticality is used. It is also exactly where the hypothesis bites: an independent audit reformulated the whole lemma as the single exact identity $S-(T_1+2T_2+T_3)=-\tr\big(P_-(Q(W)-\theta W)P_-\big)$, whose right side vanishes precisely on the critical set. Off that set the margin runs to $-0.18$ under directed minimization at $n=6$, so criticality is load-bearing rather than decorative.

\emph{Step 2: expand $W^{\#}$ bilinearly.} Since $\#$ is symmetric and bilinear, $W=W_+-W_-$ gives
\[
W^{\#}=W_+^{\#}-2\,W_+\#W_-+W_-^{\#},
\]
so the right side of \eqref{eq:gc-compress} splits into three compressions, of which the middle one carries the sign that matters.

\emph{Step 3: $\#$-positivity.} The rank-one formula
\[
(uu^{\mathsf T})\#(vv^{\mathsf T})=\tfrac12\,[u,v][u,v]^{\mathsf T}
\]
(with the constant measured as exactly $\tfrac12$ in an independent rebuild) shows that $A\#B\succeq0$ whenever $A,B\succeq0$, since both are sums of rank-one terms with nonnegative coefficients. In particular $W_+\#W_-\succeq0$, hence
\[
\tr\big(P_-(W_+\#W_-)P_-\big)\;\ge\;0 .
\]
\textbf{This is the only inequality in the proof.} Dropping it converts \eqref{eq:gc-compress} into
\[
\theta\tr(W_-)+\tr(W_-^2)\;\le\;-\tr\big(P_-W_+^{\#}P_-\big)-\tr\big(P_-W_-^{\#}P_-\big).
\]

\emph{Step 4: the partial trace.} The two remaining compressions are partial traces of $\#$-squares of positive semidefinite operators, and collapse onto the scalar invariants of $W_-$ alone.

\emph{Step 5: the polarized scalar-curvature identity.} Lemma~15 gives, for every symmetric $R$ on $\Lambda^2\R^n$,
\[
\scal\big(R^2+R^{\#}\big)=\|\Ric(R)\|^2+6\|\Alt(R)\|^2 ,
\]
with the constant exactly $6$ in every dimension and no cross-term. Applied to $R=W_-$ this identifies the right-hand side of Step 4 as $\|\Ric(W_-)\|^2+6\delta^2$, which is the assertion. Each of the five steps was machine-verified at $10^{-14}$--$10^{-16}$ as it was written.

\emph{The equality case.} Equality holds exactly when the term dropped in Step 3 vanishes, i.e.\ when $\tr\big(P_-(W_+\#W_-)P_-\big)=0$. By the rank-one formula this is a sum of squared norms of brackets $\big\|P_-[x,y]\big\|^2$ with $x$ ranging over $E_+$ and $y$ over $E_-$, weighted by products of eigenvalues, all of which are strictly positive. So equality is equivalent to the vanishing of that one bracket condition, $[E_+,E_-]\subseteq E_-$.
\end{proof}

\begin{remarkstar}[Slack zero does not by itself mean graded]
\textbf{The equality case is not ``exactly the compact symmetric pairs''}: that reading is too strong, and two corrections separate it from the truth.

\emph{Gradedness is two conditions, not three.} The form $\langle[X,Y],Z\rangle$ is totally antisymmetric (residual exactly $0$), so of the three conditions usually written down,
\[
[E_+,E_+]\perp E_-,\qquad [E_-,E_-]\perp E_-,\qquad [E_+,E_-]\subseteq E_- ,
\]
the first and third are the same statement: $\|\Pi_{\mm}[\hh,\hh]\|^2\equiv\|\Pi_{\hh}[\hh,\mm]\|^2$ identically, measured as $16/16$, $24/24$, $32/32$. The same antisymmetry says $S=0$ counts \emph{triples} drawn from $\mm$, so it is $\binom{\dim\mm}3$ conditions, and $S\equiv0$ identically whenever $\dim\mm\le2$, there being no triples at all.

\emph{The implication still fails off the subalgebra family, but the standard witness is a padding artefact.} On the family of Lemma~23 the remaining condition holds by construction, since $E_+$ is a subalgebra there, so $S=0$ genuinely is gradedness, which is why Lemma~23 goes through. Off it the implication fails. $W_{\CP^2}$ was recorded here as the counterexample, but that reading is now corrected, since $W_{\CP^2}$ is \textbf{graded in its native dimension $4$} ($\sigma_1=1.3\times10^{-30}$, with $(\so(4),\uu(2))$ a symmetric pair), and its ungradedness at larger $n$ is manufactured by the padding. Genuine ungraded $S=0$ points exist all the same, and now have a clean source: for two-eigenvalue critical points the law $\sigma_1=3\sigma_3+\tfrac{6d}n(n^2-4d)$ gives \textbf{graded $\iff\dim\mm=n^2/4$}, so every $S=0$ point at any other $\dim\mm$ is ungraded. The triangle-free strongly regular graphs supply them explicitly: Clebsch at $n=16$ has $\dim\mm=40$, $S=0$ exactly, Ricci and Bianchi both zero, criticality residual $3\times10^{-18}$.

The inequality $S\ge0$ is untouched, but what is withdrawn is the reduction of Conjecture~A that read $S=0$ as a classification hypothesis. \S\ref{sec:outside} records what this opens up.
\end{remarkstar}

\begin{certificate}
The equality/strictness split reproduces across the whole catalogue: equality at the rival, at every two-block sphere product for any split (exact by hand), at $\CP^3$, at the $\su(m)$ family, at $\HP^2$ and at the Fox plateau (\emph{every} known critical point with value above $1$), and strict at three-block products, at the $SU(3)$ bi-invariant tensor, at the anti-graded $\delta^2=3/25$ point, and at the $\theta=0$ point. The associated fusion tables were computed in exact rationals: at the rival the Peirce spectrum of $L_u$ is $\{-1,-\tfrac12,-\tfrac16,0,\tfrac13,\tfrac12,1\}\cdot\theta$ with multiplicities $(2,16,72,155,36,18,1)$, and at the balanced sphere $\{-\tfrac13,0,\tfrac13,\tfrac12,1\}\cdot\theta$ with $(81,148,54,16,1)$. In both the Peirce cap $\theta/2$ is attained on exactly the orbit dimension, $18$ and $16$, which reconfirms the cap lemma of Lemma~13 independently.
\end{certificate}

Three by-products fall out of the same computation and are used later: $\|\Ric(W_-)\|^2+6\delta^2\ge\|W\|^2$ for every Weyl $W$ with no criticality hypothesis; $\lambda_{\max}(W_+)\le\theta$; and one-line proofs of the three block-positivity facts that the lifted chain of Lemma~18 had only observed numerically, through $X_{\mathrm{inv}}=4W_+\#W_--2\theta W_--2W_-^2$.

The hypotheses are criticality, Ricci-flatness and the first Bianchi identity. The proof uses bilinearity of $\#$, its rank-one positivity, a partial trace, the polarized identity of Lemma~15, and the convention $\scal=2\tr$. The value lies in the equality case: every known critical point with value above $1$ satisfies it (the sphere pairs, the rival, the Ricci-flat K\"ahler family, $\CP^3$, $\HP^2$), and \emph{equality is not confined to those}: it also holds at $S^2\times S^4$ with $\theta^2=15/16$ and at $\CP^4$ with $\theta^2=35/36$, both below $1$. So the graded case becomes a classification problem and the ungraded case an analytic one.

\subsection{What ``graded'' names, in three inequivalent senses}\label{sec:senses}

The word carries three meanings in this problem, and they do not nest. Each result below is stated in the sense its own proof uses.

\begin{itemize}
\item \textbf{(C) pair-graded.} $W$ is the $\ad(\pp)$-invariant Weyl operator of a symmetric pair. The invariant subspace is one-dimensional at every pair checked ($\so(k)\oplus\so(l)$ and $\uu(m)$ alike, at $n=4,6,8,10,12,13,14$), so such a $W$ is unique per pair.
\item \textbf{(A) sign-graded.} $\pp=E_+(W)$ and $\mm=E_-(W)$ span $\so(n)$, equivalently $\ker W=0$, and the bracket relations hold. This is what the defect screen tests.
\item \textbf{(B) slice-graded.} $W$ commutes with $\Lambda^2(\sigma)$ for some involution $\sigma$. The reading matters: $\sigma$ an orthogonal involution of $\R^n$ is strictly narrower than $\sigma$ an involutive automorphism of $\so(n)$, since $\Lambda^2(J)$ is an involution for a complex structure $J$ that is not one on $\R^n$.
\end{itemize}

Each of (C) and (A) implies (B), and \emph{neither implies the other}. Two witnesses separate them, and both are objects the programme already cares about: $W_{\CP^2}$ at $n=4$ is pair-graded (the $\ad(\uu(2))$-invariant subspace of $\Weyl_4$ is one-dimensional and $W_{\CP^2}$ generates it), while its kernel carries $3$ of $6$ dimensions, so it is not sign-graded; and the $n=8$ witness at $\theta=1.088948$ is sign-graded and not pair-graded. The failure of (C)$\Rightarrow$(A) is confined to $\CP^2$: the $\su(m)$ eigenvalue is $(m-2)/(2(m+1)(2m-1))$, zero exactly at $m=2$, and $\CP^3$, $\CP^4$, $\CP^5$ all have $\ker W=0$.

Two consequences are load-bearing below. \emph{The defect screen is vacuous where $\ker W\ne0$}, since it traces only over $E_+\oplus E_-$: padded $W_{\CP^2}$ reports $S=0$ while $63$ of $66$ dimensions sit in the kernel at $n=12$. And \emph{a slice grading is not unique} (the $n=8$ witness admits at least eight inequivalent $\so(3)\oplus\so(5)$ splittings, and every coordinate sphere product lies in every coordinate slice), so a quantity built from $\pp$ is a function of the pair $(W,\sigma)$ and never of $W$ alone. Masses such as $\|W|_{\pp}\|^2$ inherit this: across the witness's eight gradings it spans $[0.319349,\,0.808897]$, straddling the value the equivariant point carries.

\subsection{The graded half: settled on the balanced splits}\label{sec:gradedhalf}

\begin{Conj}{8}[The graded case as a finite check]
A graded critical point is pinned by its pair's dimensions:
\[
\theta=\frac{c_{\kk}q^2+2q^2+p^2}{(2q-p)\sqrt{pq(p+q)}},\qquad p=\dim\kk,\ q=\dim\mm,
\]
and a second identity, counting the pair's structure constants two ways, reduces $c_{\kk}$ to a tabulated Casimir ratio. So where the pinning holds the value follows from three integers and one table lookup, with no curvature operator built at any stage. Since the compact symmetric pairs are classified, the graded case is then a finite check. \textbf{The pinning is established on the balanced splits and refuted off them}, so the check as run covers that half. The enumeration also needs $n\ne8$, since at $n=8$ triality carries $S^3\times S^5$ to a graded critical point outside Cartan's list.
\end{Conj}

\begin{remarkstar}
\textbf{What the pinning needs, and where it fails.} The two-eigenvalue ansatz below is the pair's invariant Weyl operator only when the isotropy's simple ideals are interchangeable. $S^3\times S^5$ carries \emph{three} distinct $\Lambda^2$ eigenvalues, so no two-eigenvalue operator represents it, and $(\so(8),\so(3)\oplus\so(5))$ would need an irrational Casimir ratio. The formula is exact on balanced splits with $c_{\kk}=2(k-2)$, and fails at $S^3\times S^5$, $S^3\times S^7$, $S^4\times S^6$, $S^5\times S^7$, $S^3\times S^9$, $S^4\times S^8$, $S^2\times S^6$ and at $\CP^m$ for $m\ge4$. The census behind ``settles it'' screens on a single-Casimir isotropy, which rejects every unbalanced split before evaluating it. \textbf{What settles the graded half above the rival's value is Lemma~69, whose route is independent of this one.}
\end{remarkstar}

\begin{proof}
On a graded critical point on which $W$ is scalar on each of $E_\pm$, the trace-free unit-norm operator supported on $\hh$ is unique, namely $W=(q\Pi_{\kk}-p\Pi_{\mm})/\sqrt{pq(p+q)}$. Substituting into $Q(W)=\theta W$ and evaluating the structure-constant sums block by block leaves a single scalar equation, whose solution is the displayed formula with $C$ the Casimir eigenvalue of $\kk$ acting on $\mm$ in the $\so(n)$-induced normalisation. A second identity, counting the pair's structure constants two ways, expresses $C$ as a tabulated Casimir ratio, so no curvature operator need be built.
\end{proof}

The law derives two closed forms. For $(\su(m),\so(m))$, substituting $p=m(m-1)/2$, $q=(m-1)(m+2)/2$, $C=m$ gives $\theta^2=m(m+1)/((m+2)(m-1))$ identically, decreasing in $m$ with the rival as its $m=2$ member, so the rival's maximality inside that family has a structural reason. For $S^m\times S^m$, $\kk=\so(m)\oplus\so(m)$ and $\mm$ the $m^2$ mixed planes give $p=m(m-1)$, $q=m^2$, $C=2(m-1)$, returning $\theta_0(2m)$ identically for every $m$.

\begin{remarkstar}
Two limits on the reduction, both stated rather than hidden. Graded is broader than two-eigenvalue: $\HP^2$ is graded with $\dim E_+=13$, but $W$ is non-scalar on $E_+$ (the $\sype(1)$ and $\sype(2)$ ideals carry different eigenvalues), so whenever $\kk$ is non-simple the ansatz needs one scalar per simple ideal, which is what Lemma~10 supplies. And $\theta$ depends only on $(p,q,C)$ while $\delta^2$ does not: $\su(4)$ at $n=8$ and $S^3\times S^3$ at $n=6$ share $(p,q,C)=(6,9,4)$ and hence the identical $\theta$, yet carry $\delta^2=1/10$ and $0$. The defect is independent data, which is why the $(\theta,\delta^2)$ frontier is a staircase rather than a curve.
\end{remarkstar}

Sphere products are provably defect-free for every dimension and every split, since the projection onto $\mm$ is a sum over decomposable bivectors whose $4$-form parts vanish. \textbf{Which of them exceed $\thstar$ is settled in closed form}: from
\[
\theta\big(S^p\times S^q\big)^2=\frac{2(n-1)(p-1)(q-1)}{(n-2)\,pq},
\]
$\theta\ge\thstar$ holds exactly when $p,q\ge5$, which needs $n^2-7n-54\ge0$ and hence $n\ge12$, identically Lemma~69's own inequality, reached there by a different route. \textbf{Below that the inequality genuinely fails}: at $n=12$, $S^3\times S^9$ gives $1.14179$ and $S^4\times S^8$ gives $1.20156$, against $\thstar=1.224745$. Everything else either falls below or does not exist: $SO(2m)/U(m)$ and the complex Grassmannians are genuinely defective but sit below the rival; the Ricci-flat K\"ahler family touches the rival's value only at the rival itself; and $(\sype(2j),\sype(j)\oplus\sype(j))$ and $(\su(2k),\sype(k))$ carry \emph{no} graded critical point at all, the latter shown in exact arithmetic.

Two safeguards proved load-bearing: the criticality relation must be \emph{tested} against each pair's true Casimirs rather than solved for them, and the isotropy ideals must be split correctly. The versions that got either wrong produced, respectively, values above the rival for pairs carrying no critical point, and a false report that no complex Grassmannian carried one.

One subtlety had to be closed. The value depends only on the abstract pair, but the defect depends on how the isotropy algebra sits inside $\so(n)$: at $n=8$ the sphere product $S^3\times S^5$ and the quaternionic-K\"ahler operator on $\HP^2$ are isospectral at the identical value $1.115546702$, yet one has zero defect and the other does not. Two non-conjugate embeddings of a single pair require an outer automorphism, which $\so(N)$ possesses only at $N=8$ (triality), so the phenomenon is confined to one dimension, below the threshold and below the rival's value.

\begin{Lem}{10}[The reducible-isotropy families carry one critical point each]
Ricci-flatness together with the first Bianchi identity cuts each one-parameter-per-factor family down to a \emph{one-dimensional} space of Weyl operators, so the critical point in each is forced rather than found.
\end{Lem}

\begin{proof}
Ricci-flatness imposes two linear conditions on the one-parameter-per-factor family, and the first Bianchi identity imposes the rest, a constraint the count is wrong without. Together they cut each family down to a \emph{one-dimensional} space of Weyl operators. A one-dimensional subspace preserved by $Q$ forces $Q(W)=\theta W$, so ``solving criticality within the family'' is vacuous: the family contains exactly one unit Weyl operator, and the content is the verified $Q(W)=\theta W$ at $5\times10^{-16}$. A circle scan does find further stationary points, and every one of them violates Bianchi.
\end{proof}

\begin{center}
\begin{tabular}{@{}llll@{}}
\toprule
Class & $n$ & $\theta$ & graded? \\
\midrule
$\CP^m$ (Hermitian pairs), $\theta^2=\tfrac{(m+1)(2m-1)}{3m(m-1)}$ & $8,10,12,14$ & $0.986013,\ 0.948683,\ 0.924962,\ 0.908514$ & graded \\
$\HP^2$, $\sype(1)\oplus\sype(2)$ & $8$ & $1.115547$ & graded, $\delta^2\ne0$ \\
$\HP^3$, $\sype(1)\oplus\sype(3)$ & $12$ & $0.852944$ & ungraded, $S=0.6734$ \\
$\mathrm{Gr}_{\mathbb C}(2,2)$ & $8$ & $0.720082$ & ungraded, $S=0.5714$ \\
$\mathrm{Gr}_{\mathbb C}(2,3)$ & $12$ & $0.596068$ & ungraded, $S=0.9631$ \\
\bottomrule
\end{tabular}
\end{center}

All sit far below $\thstar$. The $\CP^m$ closed form is Fox's: his Lemma~10.1 constructs the idempotents $S^r(\alpha)$ from an orthogonal complex structure on an $r$-dimensional subspace and his (10.4) gives $|S^r|^2=6r(r-2)/((r+2)(r-1))$, whence $\theta(S^r)^2=2(r+2)(r-1)/(3r(r-2))$, reproducing the measured values exactly at $r=2m$ and supplying the proof the measurement lacked.

That $\HP^2$ is \emph{graded} mattered more than it looks. Had it been ungraded, the recorded ungraded ceiling would have jumped from $0.900828$ to $1.115547$, cutting the gap to $\thstar$ from $0.324$ to $0.109$ and turning the ungraded half from a loose problem into a tight one. It is graded, and all three genuinely ungraded points here sit below $0.900828$.

\paragraph{A caution about closing this by citation.} Petersen--Stanfield \cite{ps2025} compute the curvature operator of every compact irreducible symmetric space, showing the eigenvalues lie in $[0,\tfrac12]$ with one, two or three distinct nonzero values: one per simple factor of the isotropy algebra, the centre carrying $\tfrac12$, and only the complex Grassmannians \emph{able} to carry three. Their tables are what Lemma~10 is checked against, and the check reproduces their Theorem~A case by case from criticality rather than from root data. But closing the graded enumeration \emph{from} their tables is unsound: graded critical points here are indexed by symmetric \emph{pairs} $(\hh,\kk)$ with $\hh\subseteq\so(n)$, whereas they classify symmetric \emph{spaces}, and a space's Weyl operator is graded only when its isotropy is a symmetric subalgebra of $\so(n)$. Their Theorem~A is also an inclusion, not an equivalence: $\mathrm{Gr}_{\mathbb C}(2,2)$ above carries two distinct values, not three.

\paragraph{The class closes without the enumeration.} What the subsections above establish family by family is superseded in scope by \S\ref{sec:symmetric}: a single dimension count covers every compact irreducible symmetric space of rank $\ge2$ together with $\mathbb{OP}^2$, and a partition argument covers the reducible ones, so no per-class check remains owed. The enumeration is retained because it is what the closed forms and the three-eigenvalue structure were read off, and because it agrees with \cite{ps2025} case by case.

\paragraph{Gradedness is strictly weaker than symmetry.} Jensen's second homogeneous Einstein metric on $S^3\times S^3$ (non-symmetric, with a negative curvature eigenvalue that Theorem~A forbids for symmetric spaces) was the natural candidate for a high-value ungraded point. Rebuilt here, its spectrum reproduces theirs exactly as a validation gate, and its Weyl part has $\theta=\sqrt{6/5}$, criticality residual $3.0\times10^{-16}$, $\delta^2=5/48$ and $S=0$, so it is \emph{graded}. So candidate ungraded witnesses cannot be found simply by picking non-symmetric homogeneous spaces.

\begin{Lem}{69}[Above $\thstar$ the pair-graded points are sphere products with both factors at least five]
For $n\ge12$, every pair-graded critical point with $\theta\ge\thstar$ is a two-factor sphere product $S^p\times S^q$ with $p,q\ge5$, and each of those is Morse--Bott nondegenerate.
\end{Lem}

\begin{proof}
A pair-graded point is the $\ad(\pp)$-invariant Weyl operator of its pair, and that operator is unique up to scale. The one-dimensionality is an invariant count rather than an appeal to the classification: at every $\Zt$-grading of $\so(n)$,
\[
\dim\operatorname{End}^{\mathrm{sym}}_K(\Lambda^2)-\operatorname{rank}(\Alt)-\operatorname{rank}(\Ric)=1,
\]
running $3-0-2$ at $p\ge5$, $4-1-2$ at $p=4$ and $3-1-1$ at $\uu(m)$, so it is uniform over the gradings rather than checked case by case, and $p=1$ carries no pair-graded point at all. At an $\so(p)\oplus\so(q)$ pair the invariant is the sphere product's. At a $\uu(m)$ pair it is $\CP^m$'s, and those are excluded by the threshold, $\theta(\CP^m)$ running $1.054093$, $0.986013$, $0.948683$, $0.924962$ at $m=3,\dots,6$ and decreasing, so every $\uu(m)$ pair with $n=2m\ge12$ falls far below $\thstar$.

Which splits survive is one inequality. With $f(m)=\tfrac m{8(m-1)}=\tfrac18\bigl(1+\tfrac1{m-1}\bigr)$, the condition $\|W\|^2\le\tfrac16$ is $f(p)+f(q)-\tfrac n{8(n-1)}\le\tfrac16$, and the $\tfrac18$ cancels to leave
\[
\frac1{p-1}+\frac1{q-1}-\frac1{n-1}\;\le\;\frac13 .
\]
At $p=2$ and $p=3$ the left side already exceeds $\tfrac13$. At $p=4$ it reads $\tfrac1{q-1}\le\tfrac1{n-1}$, that is $q\ge n$, impossible at every $n$. At $p=5$ it is $\tfrac5{(n-6)(n-1)}\le\tfrac1{12}$, that is $n^2-7n-54\ge0$, that is $n\ge12$. Hence $p,q\ge5$, which sits strictly inside the range where $S^p\times S^q$ is Morse--Bott nondegenerate.

Three or more factors are excluded by the same budget: with $k$ factors $\|W\|^2>\tfrac18\bigl(k-1-\tfrac1{n-1}\bigr)$, which already exceeds $\tfrac16$ at $k=3$. Every factor is a sphere by the quantitative $\|W_a\|^2\ge\tfrac16$ rather than by non-vanishing, which breaks no budget on its own. And $n\ge12$ enters as the dimension at which an admissible split first exists, rather than as a hypothesis the argument consumes.

\end{proof}

The enumeration this rests on carries an exception. The $\Zt$-gradings of $\so(n)$ are $\so(p)\oplus\so(q)$ and $\uu(n/2)$ for $n\ne8$. At $n=8$, where $\Out(\so(8))=S_3$, triality adds $\spin(7)$ of dimension $21$ and $\sype(1)\oplus\sype(2)$ of dimension $13$, both acting irreducibly on $\R^8$ and so carried to no reducible $\so(p)\oplus\so(q)$ by any element of $O(8)$. The second is $\HP^2=\tau(S^3\times S^5)$ at $\theta=1.115547$. Neither reaches Lemma~69, which is stated for $n\ge12$ and above $\thstar$.

The margin the lemma leaves is smaller than the two-block gap suggests. The second eigenvalue of $L_e$ below $\tfrac12$ is $(M-2)/(2(M-1))$ with $M=\max(p,q)$, so the gap is $1/(2(M-1))$, which is $\tfrac1{12}$ at $S^5\times S^7$, decaying like $1/n$. Nondegeneracy is unaffected, but an argument that spends the margin is.

\begin{remarkstar}
Certifying the hypothesis by the defect is unsound off $\ker W=0$. Padded $W_{\CP^2}$ has $\Delta=0$ vacuously and is $\ad(\pp)$-equivariant for no symmetric subalgebra: at $n=8$ its isotropy has dimension $10$ while every $\Zt$-grading subalgebra of $\so(8)$ has dimension $12$, $13$, $16$ or $21$. The hypothesis is certified by the commutant, not by $\Delta$. At $n=12$ the dimension count alone does not decide it, $\so(6)\oplus\so(6)$ having dimension $30$ against an isotropy of $32$, and structure is needed there.
\end{remarkstar}

\subsection{Conjecture A on the compact symmetric spaces, proved}\label{sec:symmetric}

The graded critical points are the compact symmetric pairs, so the class carries Conjecture~A as a statement in its own right. It holds, with its equality case, at every $n\ge12$, and the argument is a dimension count against a quadratic quantity rather than a case analysis over Cartan's list.

\begin{Lem}{46}[At a locally symmetric Einstein space the Weyl part is critical, and $\theta$ is quadratic]
Let $R$ be the curvature operator of a locally symmetric Einstein space, $\Ric=\lambda g$, and set $W=R-cI$ with $c=\lambda/(n-1)$. Then
\[
Q(W)=\lambda W,\qquad\text{so}\qquad \theta\Big(\tfrac W{\|W\|}\Big)=\frac\lambda{\|W\|} .
\]
In Petersen--Stanfield's normalisation \cite{ps2025}, where $\lambda=\tfrac12$ for every compact irreducible symmetric space, this reads
\[
\theta^2=\frac1{4\big(\|R\|^2-(\tr R)^2/N\big)},\qquad N=\tbinom n2 .
\]
\end{Lem}

\begin{proof}
Four lines. $\nabla R=0$ gives $\Delta R=0$, and the metric shrinks homothetically under Ricci flow, so $Q(R)=\lambda R$. Next $I^{\#}=(n-2)I$ gives $Q(I)=(n-1)I$. The trilinear form $B(X,Y,Z)=\langle X\circ Y,Z\rangle$ is totally symmetric, which collapses $P(W)$ and $\lambda\|W\|^2$ to the same expression, and $R\#I=\lambda I-R$ gives $R\circ I=R\#I+R=\lambda I$. Combining, $Q(W)=Q(R)-2c\,R\circ I+c^2Q(I)=\lambda(R-cI)=\lambda W$.
\end{proof}

\begin{certificate}
Two reviewers reached the identity independently by different routes on identical briefs, and a third corroborated it exactly on four values. One route goes through the bridge identity $\lambda_i=1/(c_i+2)$ between block eigenvalues and Casimir data, and locates the residual: it vanishes precisely on $\tr R=n/4$, so the Einstein hypothesis is load-bearing rather than incidental.
\end{certificate}

\textbf{What Lemma~46 removes is the error-prone half of the computation.} $\theta$ becomes a quadratic functional of published eigenvalue tables (no structure constants, no $S$-blocks, no Casimir dictionary, no dual Coxeter numbers, no Dynkin indices), and the closed forms follow by substitution.

\begin{center}
\begin{tabular}{@{}llll@{}}
\toprule
\textbf{Class} & \textbf{$\theta^2$} & \textbf{Supremum} & \\
\midrule
$\CP^m$ & $\dfrac{(m+1)(2m-1)}{3m(m-1)}$ & $\thstar$ at $m=2$ & decreasing to $\tfrac23$ \\[6pt]
$\HP^m$ & $\dfrac{(m+2)^2(4m-1)}{9m(m-1)(2m+1)}$ & $\tfrac{56}{45}$ at $m=2$ & decreasing to $\tfrac29$ \\[6pt]
$\mathrm{Gr}_{\mathbb C}(p,q)$ & $\dfrac{(p+q)^2(2pq-1)}{pq\,(4p^2q^2-p^2-q^2-2)}$ & $\tfrac{14}{27}$ at $(2,2)$ & \\[6pt]
Type II, $G\times G/\Delta G$ & $\dfrac{4(n-1)}{n(n-3)}$, $n=\dim\mathfrak g$ & $\tfrac7{10}$ at $\su(3)$ & \\
\bottomrule
\end{tabular}
\end{center}

\noindent The Type II row is the whole class of compact simple Lie groups with their bi-invariant metrics, closed in one line: the adjoint representation has Killing constant $2$ for every $\mathfrak g$, so $\lambda=\tfrac14$ universally and $\theta$ depends on nothing but $\dim\mathfrak g$. Its $\su(3)$ value $\sqrt{7/10}=0.836660$ reproduces the report's independently computed $SU(3)$ bi-invariant point by a route needing only that dimension. The one degenerate member is $\su(2)$ at $\dim3$, which is $S^3$: constant curvature, $W=0$, no $\theta$.

\begin{Lem}{47}[The dimension lemma, and the irreducible theorem]
Normalise to Einstein constant $\tfrac12$, so $\tr R=n/4$ and the centre constant $c=\tr R/N=\tfrac1{2(n-1)}$ depends on $n$ alone. Then $\|W\|^2=\sum_id_i\lambda_i^2-\tfrac n{8(n-1)}$ over the isotropy blocks, and Cauchy--Schwarz with weights $d_i$ gives
\[
\dim\hh\;\le\;\frac{3n^2(n-1)}{2(7n-4)}\quad\Longrightarrow\quad\|W\|^2\ge\tfrac16 .
\]
The lemma covers every compact irreducible symmetric space of rank $\ge2$, and the rank-one $\mathbb{OP}^2$. Hence
\[
\boxed{\ \|W\|^2\ge\tfrac16\ \text{on every compact irreducible symmetric space, with equality exactly at }\CP^2,\ }
\]
equivalently $\theta\le\thstar$, attained only there.
\end{Lem}

\begin{proof}
The block eigenvalues are \cite{ps2025} Proposition 2.2: $\lambda_i$ with multiplicity $d_i=\dim$ of each isotropy block, $\tfrac12$ on the centre, and $0$ on the remaining $N-\dim\hh$ directions of $\Lambda^2$. Cauchy--Schwarz gives $\sum_id_i\lambda_i^2\ge(n/4)^2/\dim\hh$, and substituting the universal $c$ turns $\|W\|^2\ge\tfrac16$ into the stated dimension bound.

Each classical family reduces to a polynomial positivity statement, proved by an all-nonnegative-coefficient certificate after shifting to its first admissible rank. The twelve exceptional spaces are checked directly and mostly by an order of magnitude, $E_8/SO(16)$ having $\dim\hh=120$ against a bound of $3499$. The residue is exactly the compact rank-one spaces, and every one is already settled: $S^n$ and $\R\mathrm P^n$ have $W=0$, so $\theta$ is undefined and the conjecture's hypothesis excludes them; $\CP^m$ and $\HP^m$ have the closed forms above, with minima $\tfrac16$ at $\CP^2$ and $\tfrac{45}{224}$ at $\HP^2$; $\mathbb{OP}^2$ falls under the lemma.
\end{proof}

\begin{certificate}
Exact rational arithmetic throughout. The certificates carry \emph{strictly positive} constant terms ($616$, $48$, $7200$, $114$, $3696$, $2192$, $236$, $84$), which is what rules out a second equality case and earns ``equality exactly at $\CP^2$''. $\CP^2$ itself is confirmed from isotropy data alone: $\hh=\uu(1)\oplus\su(2)$, $\lambda=\tfrac12,\tfrac16$ with multiplicities $1,3$, giving $\tr R=1=n/4$, $c=\tfrac16$ and $\|W\|^2=(\tfrac13)^2+2(\tfrac16)^2=\tfrac16$ exactly. Two reviewers, independently and on identical briefs, rebuilt seven family closed forms from the Killing forms at residual $0$ and confirmed every claim.
\end{certificate}

\begin{Lem}{48}[The reducible case, and Conjecture A on the whole symmetric class]
The Einstein normalisation forces every factor of a product to carry the same constant $\tfrac12$, so no flat factors occur. For $M=M_1\times\dots\times M_k$ the curvature operator is block-diagonal on $\Lambda^2$ and zero on the mixed terms, so factorwise
\[
\|W\|^2=\sum_a\|W_a\|^2+B,\qquad B=\sum_af(n_a)-\frac n{8(n-1)},\quad f(m)=\frac m{8(m-1)} .
\]
$B$ depends on the partition alone. $f$ is decreasing and convex, so Jensen minimises $\sum_af(n_a)$ at the balanced split for fixed $k$, and $k f(n/k)=\tfrac{kn}{8(n-k)}$ increases in $k$. Each $\|W_a\|^2\ge0$ vanishes exactly when that factor has constant curvature. Hence the reducible minimum sits at $S^{n/2}\times S^{n/2}$, with $\|W\|^2=\tfrac{n^2}{8(n-1)(n-2)}=1/(4\theta_0(n)^2)$. Combining with Lemma~47:
\[
\boxed{\ \text{for every }n\ge12,\ \text{every compact \emph{Einstein} symmetric space of dim }n\text{ has }\theta\le\theta_0(n),\ }
\]
with equality exactly at $S^{n/2}\times S^{n/2}$ for even $n$. For odd $n$ the best split falls strictly below: at $n=13$, $\theta^2=\tfrac{120}{77}$ against $\theta_0^2=\tfrac{264}{169}$.
\end{Lem}

\paragraph{The two halves meet with no gap, and the polynomial that makes them meet is the threshold's.} The irreducible bound clears the reducible minimum precisely when
\[
\tfrac16\ \ge\ \frac{n^2}{8(n-1)(n-2)}\iff n^2-12n+8\ge0 ,
\]
which is Lemma~2's equation. The conjecture's own dimension hypothesis is therefore recovered rather than assumed: $n=12$ is the smallest integer at which the sphere-product value clears $\CP^2$'s $\thstar$, with margin exactly $\tfrac1{36}$, since $\theta_0(12)^2=\tfrac{55}{36}$ against $\tfrac{54}{36}$.

\begin{center}\itshape
Corollary. For every $n\ge12$, no compact irreducible symmetric space (nor any padding of one into $\Weyl_n$, padding leaving $\theta$ unchanged) attains the maximum of $P$. Each has $\theta\le\thstar<\theta_0(n)$, strictly.
\end{center}

\paragraph{The dimension lemma fails precisely on the maximizer, and that is structural.} The isotropy of the balanced sphere product is $\so(n/2)\oplus\so(n/2)\subset\so(n)$, of dimension $n(n-2)/4$, and it exceeds Lemma~47's threshold exactly when
\[
\frac{n(n-2)}4>\frac{3n^2(n-1)}{2(7n-4)}\iff n(n^2-12n+8)>0\iff n\ge12 .
\]
So the lemma is exactly sharp enough to eliminate every competitor and exactly too weak to reach the winner, which is why the class closes only in combination with the partition argument. It is also the fourth independent route to $n^2-12n+8$, after the value comparison of Lemma~2, the cone criterion of Lemma~44, and the meeting of the two halves above.

\paragraph{The threshold $\|W\|^2\ge\tfrac16$ is $\CP^2$'s, not the conjecture's, and the difference is decisive.} Reading Conjecture~A on the symmetric class as ``$\|W\|^2\ge\tfrac16$'' is false: $S^5\times S^7$ has $\|W\|^2=\tfrac{175}{1056}$, short of $\tfrac16$ by exactly $\tfrac1{1056}$, and $S^6\times S^6$ has $\tfrac9{55}$, both below $\tfrac16$ while \emph{satisfying} Conjecture~A. The correct quadratic form is $\|W\|^2\ge\tfrac{n^2}{8(n-1)(n-2)}=1/(4\theta_0(n)^2)$, and $1/(4\|W\|^2)=\tfrac{264}{175}$ reproduces the recorded $\theta=1.228239158$ at $S^5\times S^7$.

\paragraph{A second refutation on the same line, about which isotropies are large.} The subalgebras of $\so(n)$ exceeding the dimension bound are not $\so(n)$ and $\uu(n/2)$ alone. The sharpest is $\so(n/2)\oplus\so(n/2)$ above. At low dimension there is also $\sype(1)\oplus\sype(2)\subset\so(8)$, the isotropy of $\HP^2$, at $\dim13$ against a bound of $12.92$, next to $\dim\uu(4)=16$ and $\dim\so(8)=28$. It is the only quaternionic escape: $\HP^m$ sits inside the bound for every $m\ge3$.

\paragraph{What this closes, and what it leaves.} Attack (a) of \S\ref{sec:families} is closed \emph{on the symmetric class}: $\HP^2$ and the three-eigenvalue points, which it still owed as separate family caps, fall under Lemma~47, and no per-class programme over Cartan's list remains. \textbf{What it does not exhaust is $\{S=0\}$.} Gradedness is strictly broader than symmetry (Jensen's second homogeneous Einstein metric on $S^3\times S^3$ is non-symmetric and has $S=0$, at $\theta=\sqrt{6/5}$), so the residual of \S\ref{sec:outside} keeps a non-symmetric part that a dimension count over isotropy algebras cannot see. Beyond that lies the ungraded set, whose ceiling on record stands at $0.900828$ against $\thstar=1.224745$, a gap this section does not touch. \S\ref{sec:fenceposts} is where the reason lives.

\subsection{The ungraded half}\label{sec:ungraded}

This half is open, and the honest framing is that it is a case of a problem the source literature poses and leaves open: Fox writes that ``it would be interesting to describe completely the $G$-orbits of idempotents in subalgebras of $(\mathrm{MC}(V^{*}),*)$'', and gives such a description only for $\dim V^{*}=3,4$. Partial results are therefore the realistic target rather than a full enumeration.

\begin{Obs}{12}[The ungraded half, measured]
The strict side of Lemma~7 is measured by the slack $S$ of Definition~2. On a structured class, under two bracket conditions that hold there, the slack has a closed form,
\[
S=\tr\big(P_{E_-}W_-^{\#}P_{E_-}\big)=\tfrac12\sum_{i,j}\mu_i\mu_j\big\|P_{E_-}[v_i,v_j]\big\|^2,\qquad Wv_i=-\mu_iv_i \text{ on } E_-,
\]
so on that class the slack is exactly the bracket mass $E_-$ fails to expel, weighted by the negative spectrum. The identity is \emph{conditional} on those bracket conditions: an audit found generic critical points violating them, with ratio $11.5$ rather than $\tfrac12$. The inequality $S\ge0$ itself is untouched.
\end{Obs}

\begin{certificate}
The decomposition $S=T_1+2T_2+T_3$ underlying the closed form holds at $21$ structured points, at both saved witnesses, and at generic Levenberg--Marquardt points at $n=6$ \emph{and} at the untested odd dimension $n=7$, with worst deviation $5.7\times10^{-9}$, and that only at a witness whose own criticality residual is $4\times10^{-8}$. Two upgrades came out of the audit: the two supporting lemmas are identities of the \emph{whole} Weyl space rather than of critical points (verified at non-critical points to $3.5\times10^{-11}$), and the decomposition is equivalent to the single exact identity $S-(T_1+2T_2+T_3)=-\tr\big(P_-(Q(W)-\theta W)P_-\big)$, whose right side vanishes precisely on the critical set. Safe to treat as settled.
\end{certificate}

\paragraph{The carrier inequality: valid, Fox-sharp, and unable to separate.} A companion bound assembled from the same identities (sharp at the rival in closed form, with the crossover constant emerging as $\sqrt2\sqrt{1-\tfrac14}$ from its Ricci shape) was for a while the single statement the crux reduced to. An independent audit rebuilt it from definitions, confirmed the master identity at $\le2\times10^{-14}$ across roughly $380$ critical points spanning $n=6,\dots,24$, confirmed the Fox equality exactly in all four closed-form constants, and then \textbf{refuted the separation claim} $S>0\Rightarrow\mathrm{RHS}<\thstar$ with the third-simplest sphere product there is. The balanced family $(S^m)^3$ is zero-defect, ungraded, exactly critical, and has
\[
\mathrm{RHS}=\theta\,\frac{4m-1}{3m-1},\qquad \theta^2=\frac{(3m-1)(m-1)}{3m^2},
\]
reaching $1.234230>\thstar$ at $S^8\times S^8\times S^8$ ($n=24$) and tending to $4/3$. The witness is exact to twelve digits with criticality residual $2\times10^{-16}$. The bound keeps its place as a valid Fox-sharp fence post and loses its role as the finish, so any repair must be $n$-aware or value-aware. One regularity survives and is worth carrying: across all data, every point whose bound reaches $\thstar$ has $\theta\le0.916$.

\paragraph{The multi-block family, closed in form.} For $S^{n_1}\times\dots\times S^{n_k}$, verified symbolically for $k=2,3,4$ with symbolic block sizes and against the full $Q$ at residual $10^{-16}$,
\[
\theta^2=\frac{2(n-1)^2}{(k-1)n^2+\sum_a\frac{(n-n_a)^2}{n_a-1}},\qquad
S=\frac{3\theta^2\,e_3(n_1,\dots,n_k)}{(n-1)^2},
\]
with the vertex system solved by hand rather than by sampling:
\[
w_{ab}=-\frac{\theta}{n-1}\ (a\ne b),\qquad w_{aa}=\frac{\theta\,(n-n_a)}{(n-1)(n_a-1)} .
\]
Two consequences follow immediately: $\theta^2<2/(k-1)$ strictly, so $k\ge3$ forces $\theta<1$ for every partition and every $n$; and \emph{graded $\iff k\le2$}, since $e_3$ vanishes exactly then. The closed form reproduces every recorded number and supersedes the sampled enumeration, whose seeds had missed the uniform $(6,6,6,6)$ point and the $(12,12)$ coarsening at $n=24$. It also settles that $\theta$ is globally maximised at balanced blocks, by strict convexity of $1/(n_a-1)$.

\paragraph{Why the surviving lead closes nothing.} No counterexample exists among roughly $1500$ verified critical points, and the value $1$ is approached strictly from below and never attained. But three separate objections stand.
\begin{enumerate}[leftmargin=1.6em]
\item \emph{Gradedness and $\theta$ are perfectly confounded on this family.} Ungraded means $k\ge3$, and $\theta$ is a function of $k$ alone through $\theta^2<2/(k-1)$, with the constant $1$ being exactly $\sqrt{2/(k-1)}$ at $k=3$. So the evidence shows a \textbf{block-count} statement, not a statement about gradedness.
\item \emph{The trigger does no filtering.} The right-hand side used had dropped a factor $\sqrt{1-e^2}$. Corrected, the fraction of the family with $\mathrm{RHS}\ge\thstar$ falls from $100\%$ to $0.6\%$.
\item \emph{The complementary branch is an ungranted grant.} At every witness $\delta^2=J_4=0$, so $\mathrm{RHS}-\theta=J_5/r^2$ identically, and the statement being tested is $J_5\ge0$.
\end{enumerate}

\paragraph{$J_5\ge0$, proven on the multi-block family, where it is not needed.} Substituting the closed forms above gives, for every partition,
\[
\frac{J_5}{r^2}=\frac{\theta\,\big(4n\,e_3+8\,e_4\big)}{(n-1)\sum_a n_a(n-n_a)^2}\;\ge\;0,
\]
with equality exactly when $k\le2$. On balanced blocks this is $\theta\,m/(3m-1)>0$ for every $m$. That supersedes both the balanced-only proof and a sweep over $147{,}811$ ungraded partitions. \textbf{But the chain invokes $J_5\ge0$ only above $\theta\ge(n-3)\sqrt{2/((n-1)(n-2))}$, and $\theta^2<2/(k-1)$ caps this family below that floor}, so the proof lands off the load-bearing region. What is owed is the grant where $\theta$ is large, on families this one does not reach.

\paragraph{Why the ungraded ceiling has resisted, stated as a fact about method.} Two populations exist and they do not meet.
\begin{itemize}[nosep,leftmargin=1.4em]
\item \emph{Structured constructions} (the named families, the sphere products, $(S^m)^3$) supply everything on file above $\theta=0.512$, and enumerating all two-block sphere products and the $\CP^m$ family across $n=8,\dots,24$ gives $61$ critical points, \textbf{all graded}. So the structured population contains no ungraded members that matter.
\item \emph{Random-start search} at $n=6$ returns $155$ converged critical points, \textbf{all ungraded}, occupying $\theta\in[0.007,0.512]$, entirely below the range where the question lives. At $n=8,10,12$ random-start Levenberg--Marquardt converges nothing at all ($0$ of $80$, $0$ of $60$, $0$ of $40$), a budget failure in dimension $300$.
\end{itemize}
Continuation between them is impossible for a structural reason already established here: critical points are isolated modulo their orbits, so there is no path within $\Crit(P)$ to ascend along, and a padded seed returns its own $\theta=0.390894929$ unchanged to nine digits. Ungraded critical points do exist at $n=8,10,12$ constructively (zero-padding from $\Lambda^2(\R^6)$ preserves Ricci-flatness and first Bianchi, so a padded critical point is an exact critical point of the larger problem, residual $6\times10^{-16}$), so their absence from sweeps is a property of the search, not of the geometry.

The ceiling is $\theta\le0.9067$, against a target of $\thstar=1.224745$: a gap of about $0.32$, which makes this a \emph{loose} problem where any separating constant suffices, and which stands on $n=6$ and $n=8$ with no method currently able to test it above $n=8$.

\subsection{A population where gradedness varies freely, and Conjecture A on it}\label{sec:subalgebra}

The obstruction recorded above is that on the multi-block products, ungraded $\iff k\ge3$ while $\theta$ is a function of $k$ alone, so every observation there is consistent with a statement about \emph{block count} and can say nothing about gradedness. What was needed was a population where the two vary independently. Indexing critical points by a \emph{subalgebra} rather than by a partition supplies one, since block count is not a parameter of the construction at all.

\emph{Literature first, and it paid.} The pairs $(\so(n),\hh)$ with $\hh$ and its complement both irreducible are classified: Wolf's Theorem 11.1 and its four tables \cite{wolf1968}, with a legible modern transcription in \cite{lauret2025}. Reading the tables before deriving anything is what showed that the orthogonal-ambient non-symmetric entries all have $\dim\hh=O(n)$ against $N\sim n^2/2$, which is what made the dimension question worth asking in the first place.

For a subalgebra $\hh\subseteq\so(n)=\Lambda^2\R^n$ with $\mm=\hh^\perp$, $r=\dim\hh$ and $N=\binom n2$, set
\[
W_{\hh}=c\Big(\Pi_{\hh}-\tfrac rN\Id\Big),\qquad c=\sqrt{\tfrac{N}{r(N-r)}},
\]
the unit-norm Casimir operator of $\hh$. Write $S_{\hh\hh\hh}=\sum_{a,b,c\in\hh}\langle[\varphi_a,\varphi_b],\varphi_c\rangle^2$ over an orthonormal basis, and likewise $S_{\hh\mm\mm}$ and $S_{\mm\mm\mm}$.

\begin{Lem}{23}[Conjecture A on the subalgebra family]
Let $0<r<N$ with $\hh\not\cong\mm$, and suppose $Q(W_{\hh})=\theta W_{\hh}$. Then
\[
\theta=\frac2n\sqrt{\frac{rN}{N-r}},\qquad
S_{\mm\mm\mm}=(N-r)\Big[2(n-2)-\frac{4(n-1)r}{N}\Big],
\]
and consequently
\[
r\le\frac{n(n-2)}4,\qquad\text{hence}\qquad \theta\le\theta_0(n)=\frac{\sqrt{2(n-1)(n-2)}}n,
\]
with equality in either place exactly when $W_{\hh}$ is $\mathbb Z_2$-graded, that is when $(\so(n),\hh)$ is a symmetric pair, and among those only at the balanced split. So \textbf{the maximum of $P$ over this family is $\theta_0(n)$, attained exactly at the Weyl operator of $S^{n/2}\times S^{n/2}$.}
\end{Lem}

\begin{proof}
Write $\kappa=2(n-2)$ for the Killing constant of $\so(n)$ in this normalisation and $\sigma=1+\kappa/2=n-1$. Three identities and one non-negativity do all the work.

\emph{Three identities.} $\sum_{b,c}c_{abc}^2=\kappa$ for every $a$; and, because $[\hh,\hh]\subseteq\hh$,
\[
S_{\hh\hh\hh}+S_{\hh\mm\mm}=\kappa r,\qquad 2S_{\hh\mm\mm}+S_{\mm\mm\mm}=\kappa(N-r).
\]

\emph{The eigen-equation splits.} $W_{\hh}$ has exactly two eigenvalues, $u=c(1-r/N)>0$ on $\hh$ and $v=-cr/N<0$ on $\mm$, so $Q(W_{\hh})=\theta W_{\hh}$ is two scalar equations. Eliminating $S_{\hh\hh\hh}$ and $S_{\mm\mm\mm}$ gives $S_{\hh\mm\mm}=2\sigma r(N-r)/N$ and $\theta=\sigma ur/(N-r)$, which is the stated closed form, and back-substitution gives the stated $S_{\mm\mm\mm}$.

\emph{The non-negativity.} $S_{\mm\mm\mm}$ is a sum of squares, so $2(n-2)-4(n-1)r/N\ge0$, i.e.\ $r\le n(n-2)/4$, and $\theta$ is increasing in $r$, so $\theta\le\theta_0(n)$.

\emph{The equality case.} Equality in either place means $S_{\mm\mm\mm}=0$, i.e.\ $[\mm,\mm]\subseteq\hh$, i.e.\ $W_{\hh}$ is $\mathbb Z_2$-graded, i.e.\ $(\so(n),\hh)$ is a symmetric pair. Among symmetric pairs only the balanced split reaches the cap: $\binom p2+\binom q2=n(n-2)/4$ with $p+q=n$ forces $(p-q)^2=0$, and $\uu(m)\subset\so(2m)$ has $r=m^2\ne m(m-1)$. Equality needs $n$ even: for odd $n$ the cap is not an integer and the inequality is strict.
\end{proof}

\paragraph{Criticality is free, under three hypotheses.} With $K=N_{O(n)}(\hh)$ the \textbf{full normaliser}: (F1) $\hh$ and $\mm$ both $K$-irreducible; (F2) $(\mathrm{Sym}^2\R^n)^K=\R\,\Id$; (F3) $(\Lambda^4\R^n)^K=0$. Then (F2) makes $W_{\hh}$ Ricci-flat, (F3) makes it Bianchi-free, $Q$ preserves the Weyl space, and the $K$-invariant traceless Weyl operators are $\R W_{\hh}$, so $Q(W_{\hh})=\theta W_{\hh}$ follows rather than being assumed. Taking $K$ connected \emph{loses the equality case}: $\so(4)+\so(4)\subset\so(8)$ then fails (F1), each $\so(4)$ splitting into self-dual and anti-self-dual parts, and fails (F3) through the two block volume forms. The $O(4)\times O(4)$ reflections and the block swap restore both.

\paragraph{All three hypotheses are load-bearing, and what they exclude is the Berger list.} The subalgebras exceeding the cap are $\mathfrak g_2\subset\so(7)$ ($r=14$ against $8.75$) and $\mathfrak{spin}(7)\subset\so(8)$ ($r=21$ against $12$), each satisfying (F1) and (F2) and failing only (F3); $\so(n-1)\subset\so(n)$, failing only (F2); and $\uu(m)$, $\su(m)$, $\sype(m)+\sype(1)$ and the unbalanced splits, failing (F1). Every one of them carries a parallel $4$-form, so \textbf{(F3) says exactly that $\hh$ is not a special-holonomy algebra}, and the theorem restates as a Lie-theoretic bound:

\begin{quote}
\emph{An irreducible non-special-holonomy subalgebra of $\so(n)$ has dimension at most $n(n-2)/4$, with equality only for the balanced split.}
\end{quote}

No statement of this was located in the literature or by search.

\paragraph{What it does not reach, stated plainly.} The slack on the family is $S=\tfrac{r}{nN}\big[4(N-r)-n^2\big]$, non-negative under the same condition, so \textbf{the cap is the graded compression inequality of Lemma~7 restricted here}: a sum-of-squares re-derivation and a cross-check rather than independent leverage. The family holds no near-extremal member other than the extremum: $\theta/\theta_0$ runs $1.000$ at the balanced products and then $0.845,\ 0.730,\ 0.671,\ 0.655,\ 0.612,\ 0.598,\ 0.507,\ 0.471$. And the maximum is unique only \emph{within} the family: $W_{\CP^3}$ at $n=6$ is critical at residual $5.1\times10^{-16}$ with $\theta=\theta_0(6)$ exactly, tying $S^3\times S^3$ from outside it. The unbalanced products, $\CP^m$, $\HP^m$, the complex Grassmannians and every generic ungraded point sit outside the family.

\paragraph{Why this is the strongest partial result on file.} It is Conjecture A, with its equality case, \emph{proven outright} on an infinite population, and proven at every $n$, with no dimension threshold, no criticality assumption (criticality follows from the hypotheses), and no numerical ingredient. It is also the population the ungraded half asked for: gradedness there is the symmetric-pair condition and no block count enters, so the separation is not the confounding that makes the multi-block evidence unreadable. What it is not is a step that reaches the frontier: the family is thin, and the ungraded points that make the conjecture hard lie outside it.


\paragraph{The family does not exhaust the graded set, and no threshold repairs that.} What Lemma~23 proves on the subalgebra family cannot be extended to (H) by raising a value threshold, because the graded critical set is strictly wider than the family above every threshold the question can carry.

\begin{Lem}{58}[The graded set is strictly wider than the subalgebra family]
A Casimir $W_{\hh}=c\big(\Pi_{\hh}-\tfrac rN\Id\big)$ with $0<r<N$ has exactly \emph{two} eigenvalues, $c(1-\tfrac rN)$ on $\hh$ and $-c\tfrac rN$ on its complement. Hence no graded critical point with three or more distinct eigenvalues lies in the family, and the following do not:
\begin{center}
\begin{tabular}{@{}llp{0.22\linewidth}p{0.40\linewidth}@{}}
\toprule
Witness & $n$ & $\theta$ & Why it is not a Casimir \\
\midrule
$\CP^3$ & $6$ & $\theta_0(6)=1.054092553$ & three eigenvalues \\
$(\su(3),\so(3))$ & $6$ & $\sqrt{6/5}=1.095445115$ & three eigenvalues \\
$S^5\times S^7$ & $12$ & $1.2282390$ & isotropy $\so(5)\oplus\so(7)$ has $r=31$, and the family's value there is $1.2742878$ \\
\bottomrule
\end{tabular}
\end{center}
Since $S^5\times S^7$ sits \emph{above} $\thstar$, restricting the family by a value threshold does not recover it.
\end{Lem}

\begin{proof}
Two eigenvalues is immediate from $\Pi_{\hh}^2=\Pi_{\hh}$. For the third row, the family's value depends on $\hh$ through $\dim\hh$ alone, while a sphere product's depends on the \emph{partition}, so the two agree only where the partition is balanced, and $S^5\times S^7$ is not.
\end{proof}

\noindent\textbf{What this costs.} (H) cannot be closed by classifying $\hh$, at any threshold. Lemma~23 remains what it is, a proof on an infinite population, and stops being a step towards (H). \emph{What (H) still needs is a route that does not classify $\hh$, and none is named}, so the obligation moves to (G), which has a destination and no route. \S\ref{sec:closedroutes} records this among the closed routes.

\subsection{Second-order optimality: the one constraint no earlier attack imposed}\label{sec:secondorder}

Every inequality available here constrains the unknown maximizer only through criticality, $Q(W)=\theta W$, which is first-order information. A maximizer satisfies more than that.

\begin{Lem}{24}[Second-order optimality, as an eigenvalue bound on the multiplication operator]
Let $W$ be a critical point with $\theta=P(W)>0$, put $e=W/\theta$, and let $L_e=e\circ(\cdot\,)$ for the product $X\circ Y=\tfrac12\big(Q(X+Y)-Q(X)-Q(Y)\big)$. The Hessian of $P$ restricted to the unit sphere of $\Weyl_n$ is $6\theta\big(L_e-\tfrac12\Id\big)$. Hence if $W$ is a maximizer,
\[
\lambda_{\max}\big(L_e|_{W^\perp}\big)\;\le\;\tfrac12 .
\]
\end{Lem}

\begin{proof}
$P$ is a cubic form, so $\nabla P=3Q(W)$ and the second differential is $6\theta L_e$. The constraint $\|W\|^2=1$ contributes $-3\theta\,\Id$ through its multiplier, leaving $6\theta(L_e-\tfrac12\Id)$ on the tangent space. Since $e$ is an idempotent, $\langle e,L_ee\rangle=1$, and the orbit directions are exact $\tfrac12$-eigenvectors (differentiating $Q(W_t)=\theta W_t$ along the orbit gives $e\circ\dot W=\tfrac12\dot W$), so $\tfrac12$ is attained at every critical point and the bound is sharp.
\end{proof}

\paragraph{Scope.} The hypothesis $\theta>0$ is load-bearing: at $\uu(2)\subset\so(4)$ one has $\theta=-0.866$ and the inequality reverses. Verified through the geodesic second derivative to $10^{-10}$ and audited independently.

\paragraph{Prior art: this is not new.} The statement is Proposition 6.6 of Krasnov--Tkachev \cite{kt2018}: in any Euclidean metrised algebra an extremal idempotent satisfies $\sigma(L_c|_{c^\perp})\subseteq(-\infty,\tfrac12]$ with $1$ simple, by the same argument. It appears again as Proposition 2.5 of Tkachev's survey \cite{tkachev2019}. It was derived here independently and before the literature was found, so the derivation is retained above as confirmation, but \textbf{the credit belongs to those authors}. Their framework also identifies the whole setting: $\Weyl_n$ with the cubic $P$ is an algebra in their sense, its critical points are that algebra's idempotents, and maximizing $P$ is minimizing $\|c\|$ over them.

\begin{Obs}{25}[Every ungraded critical point on file fails that bound, except on the $W_{\CP^2}$ orbit]
Across $155$ converged critical points at $n=6$ (all ungraded, all with $\theta>0$) together with every structured ungraded point on record, $\lambda_{\max}(L_e|_{W^\perp})>\tfrac12$, except on the $SO(n)$-orbit of $W_{\CP^2}$ and its paddings, where it is exactly $\tfrac12$.
\end{Obs}

That exception is forced rather than incidental: $W_{\CP^2}$ realises $M(4)=\thstar$, so it is the maximiser at $n=4$ and Lemma~24 gives $\lambda_{\max}\le\tfrac12$ there. It is the same exception Conjecture~A carries at $8\le n\le11$. And $\lambda_{\max}(L_e|_{W^\perp})\ge\tfrac12$ holds at every critical point with $\theta\ne0$ by equivariance, so ``fails the bound'' means strictly, and the content of the observation is that equality is attained only on that orbit.

The margin that observation reports is not uniform, and where it is largest the locus is not the one the conjecture is about. On the \emph{diagonal} locus the minimum of $\lambda_{\max}(L_e)$ over $\Delta>0$ critical points fits a closed form over a census,
\[
\min\lambda_{\max}(L_e)\;=\;1+\frac1{2(n-3)(n-4)},
\]
exact across seven consecutive dimensions from a census of $571$ points ($1.011905$ at $n=10$, $1.006944$ at $n=12$, $1.003205$ at $n=16$), with every minimiser connected. \textbf{So the diagonal cushion converges to $1$}: the stronger reading $\Delta>0\Rightarrow\lambda_{\max}(L_e)\ge1$ survives on the censused diagonal points, with no slack for a lossy step. Off the diagonal locus the margin over $\tfrac12$ is $1.1027\times$, held by a point at $\theta=1.080092$, $\Delta=3.494141$ and criticality residual $7.9\times10^{-16}$. Four of the six sub-$0.56$ candidates at $n=8$ are paddings, separated by the frame-independent kernel-bivector rank, whereas a coordinate test passes all six.


\paragraph{The bound is checkable block by block, and it excludes a whole family.} At a coordinate-diagonal critical point $L_e$ is block-diagonal over the characters of the diagonal sign group: one block on the diagonal locus, $\binom n4$ blocks of dimension $2$ indexed by four-subsets, and $\binom n2$ of dimension $n-3$ indexed by pairs. The four-set block of $\{i,j,k,l\}$ has spectrum $\pm\tfrac1{2\theta}\sqrt\Sigma$ in the opposite-pair sums $A=M_{ij}+M_{kl}$, $B=M_{ik}+M_{jl}$, $C=M_{il}+M_{jk}$ of the sectional matrix, with $\Sigma=A^2+B^2+C^2-AB-BC-CA$, so Lemma~24 is one inequality per four-subset.

\begin{Lem}{126}[The four-set threshold at a sphere product]
At $S^p\times S^q$ with $n=p+q$ the nonzero four-set blocks are the $(2,2)$-splits, with the common value
\[
\mu_{p,q}\;=\;\frac{n-2}{2(p-1)(q-1)},
\]
and $\mu_{p,q}=\tfrac12$ exactly when $\tfrac1{p-1}+\tfrac1{q-1}=1$, whose only solution in integers is $p=q=3$. Consequently $\mu_{2,q}=\tfrac q{2(q-1)}>\tfrac12$ at every $q$, so \textbf{no $S^2\times S^q$ is a local maximum.}
\end{Lem}

\begin{proof}
The block form gives $\mu_{p,q}$. Setting it to $\tfrac12$ and writing $u=p-1$, $v=q-1$ gives $u+v=uv$, i.e.\ $\tfrac1u+\tfrac1v=1$, and $u=1$ forces $\tfrac1v=0$, leaving $u=v=2$. At $p=2$ the value exceeds $\tfrac12$, breaching Lemma~24.
\end{proof}

\paragraph{Where this reaches and the representation-theoretic route does not.} Lemma~27 obtains sphere-product nondegeneracy from a summand decomposition, and that decomposition over-counts at $p=2$ ($86$ against $\dim\Weyl_6=84$ at $S^2\times S^4$), so the exclusion above rests on the block structure rather than on the summands. \emph{The agreement elsewhere is exact}: the block accounting returns $\dim E(\tfrac12)=18$ at $S^3\times S^3$ against a $9$-dimensional orbit, the surplus being the nine $(2,2)$-split four-set blocks sitting at exactly $\tfrac12$. Those nine are a degeneracy of the second-order test alone: the third-order obstruction is $-0.5929271$ on the whole zero cone, so no branch of critical points emanates there.

\begin{remarkstar}
\textbf{The observation is a statement about the points on file, and it does not extend to all ungraded critical points.} Padded $W_{\CP^2}$ is ungraded and critical with $\lambda_{\max}=\tfrac12$ exactly, so the bound is attained rather than exceeded there. That is forced rather than accidental: the orbit direction is an exact $\tfrac12$-eigenspace at every critical point (\S\ref{sec:shape}), and \emph{the criterion detects non-maximality rather than gradedness}, which is the reading under which the range above is evidence.
\end{remarkstar}

On the subalgebra family, where gradedness varies with no block count present, the values are exact rationals:

\begin{center}
\begin{tabular}{@{}lllll@{}}
\toprule
point & $n$ & $\theta$ & slack $S$ & $\lambda_{\max}(L_e|_{W^\perp})$ \\
\midrule
$\so(m)\oplus\so(m)$ & $8,10,12$ & $\theta_0(n)$ & $0$ & $\tfrac12$ \\
$S^3\times S^5$ & $8$ & $1.115547$ & $0$ & $\tfrac12$ \\
$\su(3)$ adjoint & $8$ & $0.836660$ & $0.571$ & $\tfrac34$ \\
$\so(3)\subset\so(5)$ & $5$ & $0.828079$ & $0.180$ & $\tfrac34$ \\
$\so(5)$ adjoint & $10$ & $0.717137$ & $0.889$ & $\tfrac56$ \\
multi-block $k\ge3$, four points & $9$--$12$ & $0.638$--$0.829$ & $0.73$--$1.09$ & $1$ \\
\bottomrule
\end{tabular}
\end{center}

The value $1$ across the multi-block products records that the idempotent is not primitive.

\paragraph{What it does not say, and the correction matters.} \textbf{This is not a characterisation of gradedness.} Two \emph{graded} families violate the bound as well: $S^2\times S^{n-2}$ at $\tfrac12+\tfrac1{2(n-3)}$, and $\uu(m)\subset\so(2m)$ at $\tfrac{m+2}{2(m+1)}$, the latter including $W_{\CP^3}$, which is graded, sits at $\theta_0(6)$ exactly, and yet fails to be a local maximizer. So the criterion detects \emph{non-maximality}, and any attempt to make it an equivalence with $S>0$ is ruled out.

\paragraph{The characterisation built on it is refuted in both directions, and the refutation is what leaves the useful statement behind.} The conjecture (that equality in Lemma~24 picks out the orbits, up to triality at $n=8$, of two-block sphere products with $p,q\ge3$ together with $\CP^2$) was tested and failed at both ends, two reviewers agreeing independently on identical briefs.

\begin{center}
\begin{tabular}{@{}p{0.30\linewidth}p{0.62\linewidth}@{}}
\toprule
\textbf{Direction} & \textbf{Witness} \\
\midrule
Saturation does not imply the list & Padded $\HP^2$ saturates at exactly $\tfrac12$ at $n=8,10,12$, criticality residual $1.4\times10^{-15}$. The record's ``padded $\HP^2$ violates at $1.0172$'' came from Weyl-projecting the \emph{full} curvature tensor in dimension $10$, which is a different operator \\
\addlinespace
The list does not imply saturation & Padded \emph{sphere products} (the conjecture's own family) violate, at exactly $\tfrac12+\tfrac1{2(n-3)}$: $0.6$ at $n=8$, $0.571428$ at $n=10$, $0.555556$ at $n=12$ \\
\bottomrule
\end{tabular}
\end{center}

Padding acts oppositely on the two families, so no reading of ``family member'' saves both directions at once. \textbf{The triality identification survives, and on stronger evidence than before}: $\tau$ was built explicitly as the spin representation on $\Delta_+=\mathbb O$, exactly orthogonal, and the $O(8)$-invariant Casimir of the top eigenspace separates the orbits: $(-2,-2,-2,0^5)$ for $S^3\times S^5$ against $(-\tfrac34)^{\times8}$ for both its triality image and $\HP^2$. So $\HP^2$ and $S^3\times S^5$ are isospectral, one diagonal and one not, and \textbf{that mechanism exists only at $n=8$}, where the diagonal maximum $\sqrt{21/16}=1.145644$ is below $\thstar$ anyway.

\paragraph{What Lemma~24 is instead: an exclusion device, and it is worth more than the characterisation was.} One direction of the bound needs no classification at all: a critical point that violates it is not a global maximizer, full stop. Applied across the catalogue, on both reviewers' independent implementations, it excludes $\CP^m$ for $m\ge3$ (at $\tfrac{m+2}{2(m+1)}$ exactly), $S^2\times S^{n-2}$, the $C_5$ blow-up (at $1.20601133$), every product with three or more blocks, $\HP^3$, $\mathrm{Gr}_{\mathbb C}(2,2)$, $\mathrm{Gr}_{\mathbb C}(2,3)$, five real Grassmannians, $SO(8)/U(4)$, $Sp(3)/U(3)$, $SO(6)/U(3)$, $Sp(2)/U(2)$, $SU(4)/SO(4)$, $SU(3)/SO(3)$, $\CP^5$, the $\su(3)$ bi-invariant, and thirteen of the fourteen $8$-cycle circulants.

\textbf{The saturating set across everything tested is: two-block sphere products with $p,q\ge3$, $\CP^2$, $\HP^2$, and paddings of the last two.} \emph{Owed}: a proof that the exclusion is exhaustive, rather than a list of the points it has been run on. \emph{Falsifier}: a critical point outside that set saturating with $1$ simple. $\HP^2$ was exactly that, so the falsifier has already fired once.

\begin{remarkstar}[Excess is a theorem, and the $E(\tfrac12)$ multiplicity refines nothing]
Differentiating $Q(W_t)=\theta W_t$ along an orbit gives $e\circ\dot W=\tfrac12\dot W$, so \emph{every orbit tangent lies in $E(\tfrac12)$} and the excess $\dim E(\tfrac12)-\dim(\text{orbit})$ is non-negative for structural reasons. Adding ``and the excess vanishes'' to the discriminator separates nothing: excess $0$ holds at the saturating $S^4\times S^4$ \emph{and} at the violating $\CP^3$ and $\CP^4$. The two non-zero rows are the already-known exceptions: $S^3\times S^3$, the degenerate split Lemma~27 excludes, and the $\circ$-decomposable $[3,3,3]$. So the second-order discriminator is the eigenvalue bound alone.
\end{remarkstar}

\begin{remarkstar}[A standing filter, withdrawn]
A tempting filter on candidate bounds reads: \emph{a bound tight at $W_{\CP^2}$ has no room left, so it cannot separate anything from $\thstar$}. The inference is unsound, and the target statement itself is the counterexample: ``$\theta\le\thstar$ on the band'' is tight at $W_{\CP^2}$ and separates the band by definition, so taken literally the filter rejects the theorem the programme is proving. What survives is the opposite reading: $W_{\CP^2}$ is the \emph{named equality model} and the target is a gap statement, and tightness at an equality model is what a sharp proof looks like.
\end{remarkstar}

\paragraph{Why it is nevertheless the live route.} Conjecture~A needs only one direction: that an ungraded critical point with $\theta\ge\theta_0(n)$ cannot be a maximizer. Three candidate constructions of the violating direction are already eliminated (rank-two directions are capped at $0$, $K$-invariant directions are multiples of $W$, and the slack's own subspace $\mathrm{Sym}^2(\mm)$ is strictly negative), leaving a collective direction in a nontrivial $K$-module, which at the $\su(3)$ point is exactly the $d_{abc}$ family of the symmetric invariant cubic.

\subsection{The eigenspace-pair test, in closed form}\label{sec:pairtest}

The collective direction is sought first among \emph{pair directions}. For two eigenvectors $\psi_\alpha,\psi_\beta$ of $e$ on $\Lambda^2$, put $X=\pi_{\Weyl_n}(\psi_\alpha\odot\psi_\beta)$. Lemma~24 then gives one inequality per pair, $\langle X,e\circ X\rangle\le\tfrac12\lVert X\rVert^2$ at a maximizer. The target on this line is that a maximizer has exactly two eigenvalues, since every critical point with two eigenvalues has $\theta\le\theta_0(n)$ by the nonnegativity of the bracket tensor alone, so the pair test is the instrument for excluding three or more.

\begin{Lem}{234}[The projected pair quotient in closed form]
Let $n\ge5$, $e\in\Weyl_n$ arbitrary, $u,v\in\Lambda^2$ orthonormal with matrices $U,V$, $Y_0=u\odot v$ and $X=\pi_{\Weyl_n}Y_0$. With $h=\mathrm{Ric}(Y_0)=-\tfrac12(UV+VU)$ and $\omega=(Y_0)_{[ijkl]}$,
\[
\langle X,e\circ X\rangle=\langle Y_0,e\circ Y_0\rangle-\frac{1}{2(n-2)}\sum e_{ijkl}h_{ik}h_{jl}-\frac38\sum e_{ijkl}\,\omega_{ijab}\,\omega_{klab},
\qquad
\lVert X\rVert^2=\tfrac12-\frac{\lvert h\rvert^2}{n-2}-\lVert\omega\rVert^2_{\mathrm{op}} .
\]
\end{Lem}

\begin{proof}
Since $e\circ\Weyl_n\subseteq\Weyl_n$, the complement $N=Y_0-X$ satisfies $\langle X,e\circ X\rangle=\langle Y_0,e\circ Y_0\rangle-\langle N,e\circ N\rangle$, and $N=h\owedge g/(n-2)+\omega$. The cross term between the Ricci and the totally skew parts is an $O(n)$-equivariant map whose only candidate is killed by the Bianchi identity. The Ricci term is read off a diagonal $h$ through the diagonal $R^\#$ formula and extended by equivariance, giving $\tfrac1{2(n-2)}$. The space of $O(n)$-equivariant maps $S^2\Lambda^4\to\Weyl_n$ is at most one-dimensional by the first fundamental theorem, and evaluating at a volume form $\omega=s\,\mathrm{vol}_B$ gives $\tfrac38$. At $n=5$ both sides vanish.
\end{proof}

\emph{The count in the last step is taken under $O(n)$}: under $SO(8)$ the tensor $\omega_{ijab}(\star\omega)_{klab}$ is a second invariant. Two reviewers audited the lemma independently.

\begin{Lem}{233}[The pair quotient at a two-block sphere product]
At $S^p\times S^q$, every pair direction built from two coordinate planes lies in a sum of the summands of Lemma~27's decomposition, so its projected quotient is a combination of Lemma~27's eigenvalues: $\mu_{p,q}=\frac{n-2}{2(p-1)(q-1)}$ for disjoint planes one in each block, $0$ for disjoint planes in one block, and $\frac{q}{2(p-1)(n-3)}$ for two planes of the $p$-block sharing an index.
\end{Lem}

All three values sit below $\tfrac12$ except $\mu_{2,q}$, which is Lemma~126's $S^2\times S^q$ saddle. \textbf{The pair inequalities are strictly weaker than second-order optimality.} At a critical point with ten distinct eigenvalues at $n=8$ and $\theta=1.088948$, $\lambda_{\max}(L_e|_{W^\perp})=0.608885$, so the point is a saddle, while the largest pair quotient is $0.492123$ and every pair inequality holds. The unstable direction lies outside every eigenspace-pair span, so the pair test can exclude a critical point with three or more eigenvalues only in combination with the first-order rows, never alone.

\begin{Lem}{235}[The quintic corrections cancel from two combinations, at three values]
Let $W$ be a unit maximizer whose $e$ has three eigenvalues $\varepsilon_1<\varepsilon_2<\varepsilon_3$, with gaps $p=\varepsilon_2-\varepsilon_1$ and $q=\varepsilon_3-\varepsilon_2$, and let $S_{ab}\le0$ be the pair inequality summed over the eigenspaces $E_a$, $E_b$. Then
\[
p\,S_{12}+(p+q)\,S_{13}\le0,\qquad q\,S_{23}+(p+q)\,S_{13}\le0
\]
are exact inequalities in the class data (eigenvalues, multiplicities, the bracket totals $\tau_{abc}$ and one Gram parameter $G_{12}$). Every nonnegative combination of $S_{12},S_{13},S_{23}$ free of the two quintic moments of Lemma~234's corrections is a nonnegative combination of these two.
\end{Lem}

\begin{Lem}{236}[The gap-weighted pair sum is spectral]
For every unit $W\in\Weyl_n$ with $\theta=P(W)\ne0$, writing $s_{\alpha\beta}=\langle X,e\circ X\rangle-\tfrac12\lVert X\rVert^2$ at $X=\pi_{\Weyl_n}(\psi_\alpha\odot\psi_\beta)$,
\[
\sum_{\alpha,\beta}(\varepsilon_\alpha-\varepsilon_\beta)^2s_{\alpha\beta}=\frac{n+4}{6}\Big[(n-1)\tr e^3-(n-3)\tr e^2\Big].
\]
Consequently a unit maximizer satisfies $(n-1)\tr W^3\le(n-3)\theta$.
\end{Lem}

\textbf{The consequence is inert above $\sqrt{3/2}$}: Okumura's bound on $\tr W^3$ makes it automatic whenever $\theta\ge\sqrt{3/2}$ at $n\ge11$, and the maximal torus attains equality at $n=6,8$. So the only part of the pair test that is free of $G_{12}$ says nothing in the range the conjecture is about, and what the test can still say at three values sits in Lemma~235's two inequalities.

\begin{Lem}{237}[The Ky Fan bracket cap]
For a unit $X_0$ and orthonormal $X_1,\dots,X_r$ in $\so(n)$, $\sum_j\lVert[X_0,X_j]\rVert^2\le\kappa_r$, the largest sum of $r$ top eigenvalues of $-\mathrm{ad}_X^2$ over unit $X$, with $\kappa_1,\kappa_2=2,4$ for $n\ge4$ and $\kappa_4,\kappa_6=6,8$ for $n\ge6$. Along the spectral ray of a relaxation witness at $n=12$ that satisfies every linear class constraint with the pair cap (\S\ref{sec:pairfences}), the class system with these caps gives $\theta\le\frac{83\sqrt{21}}{294}\approx1.29372$ exactly, at every $n$.
\end{Lem}

\begin{proof}
Ky Fan's maximum principle applied to $-\mathrm{ad}_{X_0}^2$, whose root-slot spectrum is the one in the proof of Lemma~31, each slot of multiplicity two. The ray ceiling is an exact dual certificate on the class rows together with the cap $3P_{00}+3P_{02}\le8$.
\end{proof}

\begin{Lem}{238}[A seven-frame's bracket total, and the isotropic bound $42$]
Let $X_1,\dots,X_7$ be orthonormal in $\so(n)$, $Q_i=-X_i^2$, $P=\sum_iQ_i$ with $d=\operatorname{rank}P$, and $T=\sum_{i\ne j}\lVert[X_i,X_j]\rVert^2$. With $r=\min(7,d-1)$,
\[
T\le\Big(1-\frac1r\Big)\big(196+2\tr P^2\big)-168,
\]
and when $P$ is a multiple of the identity on its range, $T\le42$. The bound $42$ is attained at $d=8$ by the octonion frame and at $d=7$ by the complement of $\mathfrak g_2$ in $\so(7)$.
\end{Lem}

The anisotropic remainder contains a frame at $42$ as well, the pencil frame $X_i=u\wedge w_i$ through one common vector, where the inequality above loses. Lemma~238's identity and its bound at $r=7$ are machine-checked in Lean. Lemmas~233--238 each carry two independent audits. \textbf{The class-level pair test at three values cannot close the question} (\S\ref{sec:pairfences}): Lemma~237 moves the relaxation witness without removing it.

\subsection{Beyond the pair test: the three-value class system}\label{sec:classsystem}

A two-value spectrum is settled outright, and that is what makes the number of eigenvalues a route to Problem~P.

\begin{Lem}{105}[Two eigenvalues suffice]
Let $W$ be a unit $\theta$-critical point whose spectrum takes exactly two values, $u>0$ with multiplicity $r$ and $-w<0$ with multiplicity $N-r$. Then
\[
\theta\ \le\ \min\Big\{(n-1)w,\ \frac{(n-2)u^2-2uw-w^2}{w}\Big\}\ \le\ \theta_0(n),
\]
with equality in the second inequality exactly at $r=n(n-2)/4$, which forces $u/w=n/(n-2)$ and $n$ even, the balanced sphere product's spectrum. At odd $n$ the inequality is strict.
\end{Lem}

Two reviewers rebuilt the reduced system independently, and the lemma, its equality case and the odd-$n$ strictness are machine-checked in Lean at the operator. So Problem~P follows from one statement: \emph{a global maximiser has a two-value spectrum}. The rest of this subsection and the next test that statement at three values.

At three eigenvalues the pair test becomes linear constraints on \emph{class data}: the eigenvalues and multiplicities, the bracket totals $\tau_{abc}$ between eigenspaces, and isotropy parameters $G_{ab}$. The results below show that at $n=12$ no constraint at that level excludes a three-valued critical point above $\sqrt{3/2}$, and then turn to single eigenspaces and to pairs of them.

\begin{Lem}{240}[The seven-frame lemma]
For $X_0,\dots,X_6$ orthonormal in $\so(n)$, any $n$, with $\langle X,Y\rangle=-\tfrac12\tr XY$,
\[
T=\sum_{i\ne j}\lVert[X_i,X_j]\rVert^2\le54.064 .
\]
\end{Lem}

The proof splits on the top eigenspace of $-\mathrm{ad}_{X_i}^2$ at an element of small row deficit, which is always a triangle of three planes or a star through one vector.

\begin{Lem}{239}[The seven-frame lemma, triangle side]
If some $X_i$ has row deficit $8-\sum_{j\ne i}\lVert[X_i,X_j]\rVert^2<0.277$ and the six largest eigenvalues of $-\mathrm{ad}_{X_i}^2$ are the positive slots of three of its planes, then $T\le53.05$.
\end{Lem}

The star side is a separate count. Two reviewers audited each side independently, and the star count was repaired and re-audited with margin $0.0106$. The best frames known total $42$ (Lemma~238), so Lemma~240 is not sharp. Added to the capped class system it puts the ray ceiling of \S\ref{sec:pairfences} below $\sqrt{3/2}$, and the next two results show that this does not close the class system.

\begin{Lem}{241}[An exact point of the capped class system off the ray]
At $n=12$ the spectrum $e=(-\tfrac15,0,\tfrac15)$ with multiplicities $(7,52,7)$, so $\theta=5/\sqrt{14}\approx1.336$, carries class values that satisfy, in exact rational arithmetic, the first-order rows, the pair caps, the Ky Fan caps of Lemma~237, the seven-frame cap of Lemma~240 for every composition, the Gram constraints, the per-pair inequalities and Lemma~235's window. Hence no argument using only these class-level constraints excludes a three-valued critical point above $\sqrt{3/2}$ at $n=12$.
\end{Lem}

\begin{Lem}{242}[An exact point under every valid seven-frame cap]
At $n=12$ the spectrum $e=(-\tfrac3{20},0,\tfrac3{20})$ with multiplicities $(12,42,12)$, so $\theta\approx1.361$, satisfies the same constraints with every seven-frame class total at most $\tfrac{1172}{33}\approx35.5$. Since orthonormal seven-frames attain $42$, no valid sharpening of Lemma~240 excludes it.
\end{Lem}

Both are points of the relaxation, not Weyl tensors, and each rests on one audit. Lemma~242's point is not realisable, and the constraint that excludes it is a subspace cap rather than a frame cap:

\begin{Lem}{243}[The Gaussian subspace cap]
Let $X_1,\dots,X_k$ be orthonormal in $\so(n)$, $P=\sum_i(-X_i^2)$ and $T=\sum_{i\ne j}\lVert[X_i,X_j]\rVert^2$. Then
\[
T\le\Big(1-\frac1r\Big)\big(4k^2+2\tr P^2\big)-4k(k-1),\qquad r=\min(k,n-1).
\]
Applied to a union of eigenspaces it is a constraint linear in the class values $(\tau,G)$.
\end{Lem}

It excludes Lemma~242's point and keeps Lemma~241's, and it is machine-checked in Lean. With every class constraint, this one included, a bounded screen at $n=12$ still finds $190$ of the $2080$ multiplicity triples with a feasible direction just above $\theta_0(12)$.

\paragraph{Single eigenspaces.} Criticality also holds vector by vector. For a unit $x$ in an eigenspace $E_c$, write $b_{ab}(x)$ for its bracket weight between $E_a$ and $E_b$.

\begin{Lem}{244}[The per-vector rows, and a ceiling on the symmetric family]
At a critical point with three eigenvalues, every unit $x\in E_c$ satisfies
\[
\sum_{a,b}b_{ab}(x)=n-2,\qquad\sum_{a,b}\lambda_b\,b_{ab}(x)=-\lambda_c,\qquad\sum_{a,b}\lambda_a\lambda_b\,b_{ab}(x)=\lambda_c(\theta-\lambda_c),
\]
with Ky Fan caps on the weights. On the symmetric family, eigenvalues $(-a,0,a)$ with $a=1/\sqrt{2p}$ and multiplicities $(p,N-2p,p)$, every unit $x$ in the negative eigenspace gives $\theta\le\big(\min(\kappa_{p-1}(x),\kappa_p(x)-2,n-4)-1\big)/\sqrt{2p}$.
\end{Lem}

At $n=12$ this excludes every symmetric point above $\theta_0(12)$ with $p\ge17$, and it bars equal-angle elements from the negative eigenspace at $p=12$.

\begin{Lem}{245}[Isotropy, capped direction by direction]
Let $W$ be a Weyl tensor with three eigenvalues $\lambda_0<\lambda_1<\lambda_2$ and put $A_c=\sum_{\psi\in E_c}\psi\psi^{\mathsf T}$ on $\mathbb R^n$. For every unit $x$ the weights $a_c=\langle A_cx,x\rangle$ satisfy $a_0+a_1+a_2=n-1$, $\sum_c\lambda_ca_c=0$ and $a_c\ge0$, so each $a_c$ is at most the largest value $h_c$ of that coordinate on the resulting segment. In particular, if $E_c\supseteq e\wedge U$ with $U\perp e$, then $\dim U\le h_c$.
\end{Lem}

The bound is sharp at $S^5\times S^7$, where it is attained by stars in all three eigenspaces. On the symmetric family, $h_0=h_2=\tfrac{n-1}2$, so a star eigenspace must be small, but both single-eigenspace results together leave a window of stars open for large $n$ (post F77, \S\ref{sec:pairfences}). Closing it needs two eigenspaces at once:

\begin{Lem}{246}[Isotropy of two eigenspaces together]
With $A_c$ as in Lemma~245, the $A_c$ commute, and for any two classes $a\ne b$, the eigenvalues $\beta_1,\dots,\beta_n$ of $A_a+A_b$ satisfy
\[
\beta_j\le\sum_{i\ne j}\min(1,\beta_i)\qquad\text{for every }j .
\]
Consequently, if an extreme eigenspace is a star $E_0=e\wedge U$, then $\lambda_1<0$, and $\lambda_1>0$ for a star $E_2$. So on the symmetric family, and whenever $\lambda_1=0$, no extreme eigenspace is a star, at every $n$.
\end{Lem}

This removes every configuration in post F77's window. What remains owed on this line is the reach beyond stars, to near-decomposable eigenspaces, which the per-vector caps admit. Lemmas~243 to 246 each carry two independent audits.

\paragraph{The pair operator, traced over eigenspaces.} A second use of the per-vector rows traces them over whole eigenspaces, giving nine linear rows in the traces $t_{abc}=\tr(P_aB_{bc})$ of the bracket operators, and tests the pair operator on $\pi_{\Weyl_n}(E_0\odot E_2)$ against Lemma~24's bound.

\begin{Lem}{247}[The pair operator's bracket part on the symmetric family]
On the symmetric family, eigenvalues $(-a,0,a)$ with $\dim E_0=\dim E_2=p$, $a=1/\sqrt{2p}$ and $\kappa=\theta\sqrt{2p}$, put $D=\tfrac12(t_{222}-t_{002}+p)$. The nine traced rows give
\[
16D=(n-5-\kappa)\,p-4t_{000}-t_{011}-4t_{122},
\]
and the bracket part $B_{\mathrm{lead}}$ of the pair operator has $\lambda_{\max}(B_{\mathrm{lead}})\ge T/p$ for a trace $T$ with $2T/p-\tfrac12=-D/(p\kappa)$. So wherever $\theta\sqrt{2p}>n-5$, the bracket part exceeds Lemma~24's bound $\tfrac12$ by at least $(\kappa-(n-5))/(16\kappa)$.
\end{Lem}

\textbf{The hypothesis is never met at a critical point} (post F78): the rows also give $4t_{000}+4t_{001}+t_{011}=(n-5-\kappa)p$, so every point of the rows with nonnegative traces has $\kappa\le n-5$. The certificate extends to general three-value spectra, where on the slice $\lambda_1=0$ it can fire only when $\lvert\lambda_0\rvert>\lambda_2$, and there the window it leaves is nonempty.

\begin{Lem}{248}[The operator ceiling]
Let $W$ be a unit critical point with eigenvalues $\lambda_{\min}=-a<0<b=\lambda_{\max}$, any number of distinct values. Then
\[
a\,(\theta+a+2b)\ \le\ (n-2)\,b^2 ,
\]
and more generally $\lambda_c(\theta-\lambda_c+2b)+(n-2)b^2\ge0$ for every eigenvalue $\lambda_c$. At $a=b$ this is $\theta\le a(n-5)$.
\end{Lem}

It is a closed-form ceiling on $\theta$ from the two extreme eigenvalues alone, and it is the ceiling behind post F78. Traced over the extreme eigenspace it is a nonnegative combination of the rows and the traces, so it adds no exclusion to a class-level test that already carries both. Lemmas~247 and 248 each carry two independent audits and are machine-checked in Lean with the traced rows as hypotheses.

\paragraph{The symmetric family at $n=12$.} On the symmetric family (eigenvalues $(-a,0,a)$, $a=1/\sqrt{2p}$, $\kappa=\theta\sqrt{2p}$, $m=\dim E_1$) the per-vector rows sharpen through the spectrum of $-\ad_x^2$, and at $n=12$ they cut the multiplicity to a window.

\begin{Lem}{249}[The per-vector ceiling through the smallest eigenvalues of $-\ad_x^2$]
Every unit $x\in E_0$ gives $\kappa\le n-5-\tfrac12\mu_m(x)=\tfrac12\kappa_{2p}(x)-3$, with $\mu_m(x)$ the sum of the $m$ smallest eigenvalues of $-\ad_x^2$ on $\so(n)$ and $\kappa_{2p}(x)$ the sum of the $2p$ largest.
\end{Lem}

It dominates both terms of Lemma~244's ceiling and is tight at post F77's star, so it closes no window by itself. What it does close is the top of $p$ at $n=12$, one value at a time.

\begin{Lem}{250}[No critical point above $\theta_0(12)$ with $p=16$]
At $n=12$ every symmetric critical point with $p=16$ has $\theta<\theta_0(12)$.
\end{Lem}

\begin{Lem}{251}[No critical point above $\theta_0(12)$ with $p=15$]
At $n=12$ every symmetric critical point with $p=15$ has $\theta<\theta_0(12)$.
\end{Lem}

\begin{Lem}{252}[The coupled fourth moment, and $p=14$]
At every symmetric three-value critical point, averaging over the unit sphere of $E_0$,
\[
(p+2)\,\mathbb E\tr x^4+2\,\mathbb ER_1=2n-\tfrac3p\lVert\Omega_0\rVert^2,\qquad\text{so}\qquad\mathbb E\big[(p+2)\tr x^4+\mu_m(x)\big]\le2n .
\]
At $n=12$ every symmetric critical point with $p=14$ has $\theta<\theta_0(12)$.
\end{Lem}

\begin{Lem}{253}[$\kappa\le p-2$ at every $n$, and the window at $n=12$]
On the symmetric family at every $n$, $\kappa\le p-2$. At $n=12$ a symmetric three-value critical point above $\theta_0(12)$ has $8\le p\le13$.
\end{Lem}

Each argument pairs Lemma~249 at one element with a sphere average of $\tr x^4$, which is small, and a connectivity argument that forces every element of $E_0$ into one type, where $\tr x^4$ is large. Each lemma carries two independent audits, and each scalar step is machine-checked in Lean. Two steps rest on exact enumerations: the $p=14$ case checks $1{,}700{,}820$ edge structures by integer elimination, and the $p=7$ case every seven-subset of $D_6$'s roots. \emph{On the open window $8\le p\le13$ every admissible element is near-decomposable}, so what remains there is a rigidity statement for a $p$-dimensional space of near-decomposable $2$-forms, and the next subsection reduces the whole family to a single norm bound instead.

\subsection{The symmetric family reduces to statement \texorpdfstring{H$_\#$}{H\#}}\label{sec:statementH}

The per-vector line above works one element of $E_0$ at a time. Summed over the whole of $V=E_0\oplus E_2$ it becomes a statement about the subspace, and then about Hamilton's $\#$ on $\mathrm{Sym}^2\so(n)$ alone, with the Weyl tensor gone. Throughout, $F(A,B,C)=\langle A\#B,C\rangle$ on $\mathrm{Sym}^2\so(n)$ with the Frobenius norm, $N=\dim\so(n)$, and $\mathcal M(V)=\sum_{u,v,w}\langle[u,v],w\rangle^2$ over an orthonormal basis of $V$, the internal bracket mass.

\begin{Lem}{255}[The symmetric family reduces to an internal-mass cap]
Let $W$ be a unit critical point with spectrum $(-a,0,a)$, $a=1/\sqrt{2p}$, and $d=2p=\dim V$. Then $\theta\le\mathcal M(V)/(2d^{3/2})$, and $\mathcal M(V)\le8d^{3/2}$, so $\theta\le4$.
\end{Lem}

So the family's two-value question at $n$ follows from the \emph{internal-mass cap}: every subspace $V\subseteq\so(n)$ of dimension $d$ has $\mathcal M(V)<2\theta_0(n)\,d^{3/2}$. On subalgebras the cap holds with room.

\begin{Lem}{256}[The internal-mass cap on subalgebras]
For $n\ge9$ and every Lie subalgebra $\hh\subseteq\so(n)$ of dimension $d$, $\mathcal M(\hh)/d^{3/2}\le2(n-2)/\sqrt N$, the value of $\so(n)$ itself. Hence a symmetric three-value critical point whose $E_0\oplus E_2$ is closed under the bracket has $\theta\le\sqrt2(n-2)/\sqrt{n(n-1)}<\theta_0(n)$.
\end{Lem}

\begin{Lem}{257}[The cap follows from a quartic bound]
For every subspace $V$ of dimension $d$ with projector $\Pi$, $\mathcal M(V)\le2\sqrt d\,\lVert\Pi\#\Pi\rVert_F$, with equality on simple subalgebras. So the symmetric family at $n\ge12$ follows from the \emph{quartic bound} $\lVert\Pi\#\Pi\rVert_F\le(n-2)d/\sqrt N$ for every rank-$d$ projector.
\end{Lem}

\begin{Lem}{258}[The quartic bound on coordinate subspaces]
For every graph on $n\ge3$ vertices with $e$ edges, $\sum_{i<k}\operatorname{codeg}(i,k)^2\le(n-2)^2e^2/N$, with equality among graphs with an edge exactly at $K_n$. So the quartic bound holds on every subspace spanned by coordinate planes $e_i\wedge e_j$.
\end{Lem}

The quartic bound is false at $n\le8$, where $\Lambda^2_+\subset\so(n)$ violates it, and three independent Grassmannian screens at $n=9$ to $12$ find no violator. The three statements so far are one statement:

\begin{Lem}{259}[Statement H$_\#$]
Let $\kappa_n=\sup F(X,X,X)/\lVert X\rVert^3$ on $\mathrm{Sym}^2\so(n)$. \emph{(1)} The constant is also the norm of $\#$ as a bilinear map,
\[
\kappa_n=\sup_{X,Y\ne0}\frac{\lVert X\#Y\rVert}{\lVert X\rVert\lVert Y\rVert},
\]
and the quartic bound, the internal-mass cap and the symmetric family (where $\theta=F(W,W,W)$) each follow from \textbf{statement H$_\#$},
\[
\kappa_n=\frac{n-2}{\sqrt N}\qquad(n\ge9),
\]
which is false for $4\le n\le8$. \emph{(2)} Every symmetric orthogonal automorphism $\sigma$ has $\sigma\#\sigma=(n-2)\sigma$, so $I$ and every involutive $\mathrm{Ad}_g$ attain $(n-2)/\sqrt N$, and $I$ is a strict local maximum for $n\ge7$. \emph{(3)} Every $\#$-idempotent $C$ ($C\#C=C$) of trace $s$ satisfies H$_\#$, in the form $\lVert C\rVert^2\ge N/(n-2)^2$, unless $-\tfrac n{2(n-2)}<s<\tfrac{(n-1)(n-4)}{2(n-2)}$.
\end{Lem}

H$_\#$ is the sharp norm of $\#$ on symmetric operators, and its equality set is strictly larger than $\{I\}$, which is why no lossy step can prove it. Lemmas~255 to 259 each carry two independent audits, and the scalar cores of Lemmas~255 to 258 are machine-checked in Lean.

\paragraph{H$_\#$ as a bound on idempotents.} At a maximiser $X$ of $F$ on the unit sphere, $C=X/F(X)$ is a $\#$-idempotent with $\lVert C\rVert=1/\kappa_n$, so H$_\#$ is a lower bound on the norm of idempotents.

\begin{Lem}{260}[H$_\#$ is a norm bound on positive-definite idempotents]
$\kappa_n$ is attained on the positive semidefinite cone, and the range of a positive semidefinite idempotent is a semisimple subalgebra $\hh$, on which it is a positive-definite idempotent of $\hh$'s own $\#$. So H$_\#$ for $n\ge9$ is the bound $\lVert C\rVert^2\ge N/(n-2)^2$ on positive-definite idempotents of compact semisimple subalgebras of $\so(n)$. Every such idempotent satisfies $\tr(CK)\le d$, $K=\bigoplus_ic_i\Pi_i$ over the simple ideals ($\Pi_i\#\Pi_i=c_i\Pi_i$), with equality exactly at $\bigoplus_i\Pi_i/c_i$.
\end{Lem}

The idempotents $\bigoplus_i\Pi_i/c_i$ meet the bound by Lemma~256, so H$_\#$ would follow from uniqueness (U): \emph{these are the only positive-definite idempotents}. U holds on two families and fails in general.

\begin{Lem}{261}[U on the coordinate-diagonal family]
If $C=\sum_{i<k}a_{ik}\,e_{ik}e_{ik}^{\mathsf T}$ on $\so(n)$ has every $a_{ik}>0$ and $C\#C=C$, then $C=I/(n-2)$.
\end{Lem}

\begin{Lem}{262}[U on symmetric-pair families]
For a symmetric pair $\mathfrak g=\kk\oplus\pp$ of a compact simple $\mathfrak g$ ($I\#I=cI$) whose $\kk$ has one induced Casimir value, the idempotents of the form $aP_\kk+bP_\pp$ are $0$, $P_\kk/c_\kk$, $I/c$ and $\sigma/c$.
\end{Lem}

\begin{Lem}{263}[U is false]
On an isotropy-irreducible two-component family $aP_\kk+bP_\pp$ a second positive-definite idempotent exists exactly when $\mathrm{Cas}_\kk(\pp)<\tfrac12\mathrm{Cas}_\kk(\kk)$. At $SU(3)\subset G_2$, $C=\tfrac9{28}P_{\su(3)}+\tfrac3{28}P_\pp$ is a positive-definite idempotent on $\mathfrak g_2$ other than $I/4$, of squared norm $\tfrac{351}{392}>\tfrac78$.
\end{Lem}

\begin{Lem}{267}[U fails on $\so(n)$ itself]
For every $n\ge7$, under $SU(2)_+\times SU(2)_-\times O(n-4)\subset SO(n)$ the idempotent equations on $a_+P_++a_-P_-+a_sP_s+bP_V$ have a solution with $a_+\ne a_-$, a positive-definite idempotent on $\so(n)$ other than $I/(n-2)$. It is heavier than $I/(n-2)$, and unstable.
\end{Lem}

Both counterexamples are heavier than the identity's multiple, and both are unstable, which is the property a maximiser cannot have.

\paragraph{Stability replaces uniqueness.} Call an idempotent $C$ \emph{stable} when $\langle C\#H,H\rangle\le\tfrac12\lVert H\rVert^2$ for every $H\perp C$, the second-order condition at a maximum of $F$.

\begin{Lem}{264}[H$_\#$ follows from U$_{\mathrm s}$]
The idempotent of a maximiser is stable. On a compact semisimple $\hh$ with two or more simple ideals no positive-definite idempotent is stable. Hence H$_\#$ at $n\ge9$ follows from \textbf{U$_{\mathrm s}$}: on a compact simple $\hh$, every stable positive-definite idempotent equals $I/c$. On every two-component ansatz whose ambient algebra is $\so(n)$, $n\ge7$, U$_{\mathrm s}$ holds.
\end{Lem}

\begin{Lem}{265}[Where $I/c$ is stable]
On a compact simple $\hh$, $I\#\cdot$ acts on each irreducible $V\subseteq\mathrm{Sym}^2\hh$ by $c-\mathrm{Cas}_V/4$, so $I/c$ is stable exactly when no component's Casimir falls below the adjoint's. This holds strictly on $\so(n)$ ($n\ge7$), $\mathfrak g_2$, $\mathfrak f_4$, $\mathfrak e_6$, $\mathfrak e_7$, $\mathfrak e_8$ and $\su(2)$, with equality on $\su(m)$, $m\ge3$, and fails on $\sype(m)$, $m\ge2$.
\end{Lem}

\begin{Lem}{266}[Two constraints on the maximiser]
$F$ vanishes on the orbit tangents of every idempotent, so a third-order condition at a maximum binds only off the orbit. A maximiser supported on a simple $\hh\subsetneq\so(n)$ has $\tr C\le d/\mathrm{Cas}_\hh(V)$ for every $\hh$-irreducible $V\subseteq\hh^\perp$.
\end{Lem}

\begin{Lem}{270}[H$_\#$ needs only the light side]
On a compact simple $\hh$ with $I\#I=cI$, $d=\dim\hh$, let $C\ne I/c$ be a positive-definite idempotent with $\tau=\tr C$ and $q=\lVert C\rVert^2$. The projection of $I$ onto $C^\perp$ destabilises $C$ exactly when $q(2c\tau-d)>\tau^2$, which forces $q>d/c^2$. Hence H$_\#$ at $n\ge9$ follows from \textbf{U$_{\mathrm s}'$}: \emph{every stable positive-definite idempotent on a compact simple $\hh$ has $\lVert C\rVert^2\ge d/c^2$}, which asks nothing of idempotents heavier than $I/c$.
\end{Lem}

\paragraph{The torus families.} Fix a maximal torus of $SO(2m)$ and write a torus-invariant operator as $C=A\oplus\bigoplus_\alpha\mu_\alpha I_{R_\alpha}$, with $A$ on the Cartan and a scalar $\mu_\alpha$ on each root plane. With $U_{ij}=\mu_{\varepsilon_i-\varepsilon_j}$, $V_{ij}=\mu_{\varepsilon_i+\varepsilon_j}$, $P=U+V$, $Q=U-V$ and $d_i=\sum_l(U_{il}^2+V_{il}^2)$, every family above becomes matrix inequalities in $(P,Q)$.

\begin{Lem}{268}[The torus-invariant idempotents of $\so(2m)$]
$C$ is a positive-definite idempotent exactly when $A=2\sum_\alpha\mu_\alpha^2u_\alpha u_\alpha^{\mathsf T}$, $u_\alpha=\alpha/\sqrt2$, and $\mu>0$ solves $\mu_\alpha=4\mu_\alpha^3+\mu_\alpha\sum_{\beta\sim\alpha}\mu_\beta^2+\sum_{\beta\pm\gamma=\pm\alpha}\mu_\beta\mu_\gamma$.
\end{Lem}

\begin{Lem}{271}[The torus system in $P$ and $Q$]
Every torus-invariant $C$ has $\lVert C\rVert^2=\lVert A\rVert^2+\tfrac12\lVert P\rVert^2+\tfrac12\lVert Q\rVert^2$ and
\[
F(C,C,C)=3\sum_iA_{ii}d_i-3\sum_{i\ne j}A_{ij}P_{ij}Q_{ij}+\tfrac12\tr P^3+\tfrac12\tr Q^3 .
\]
The idempotent system is $P_{ij}(1-d_i-d_j-2Q_{ij}^2)=(P^2)_{ij}$ with its mirror under $P\leftrightarrow Q$, and $Q=0$ only at $I/c$.
\end{Lem}

Lemmas~269 and 272 are two further constraints on these solutions: a sufficient test for instability, and a Perron-root bound on the Cartan weights. The one that carries H$_\#$ is this:

\begin{Lem}{273}[H$_\#$ on the torus slice from one matrix inequality]
Put $b^2=\tfrac12(\lVert P\rVert^2+\lVert Q\rVert^2)$, $T=\tfrac12(\tr P^3+\tr Q^3)$, $g^2=\sum_id_i^2+\sum_{i\ne j}P_{ij}^2Q_{ij}^2$ and $\lambda_m=\sqrt{2m}\,(m-2)/(m-1)^{3/2}$. For $m\ge5$,
\[
\textbf{(J1)}\qquad T+\lambda_m\,g^2\le\lambda_m\quad\text{whenever }b=1
\]
implies H$_\#$ on the torus slice: $F(X)\le\kappa\lVert X\rVert^3$ for every torus-invariant $X\in\mathrm{Sym}^2\so(2m)$, $\kappa=2(m-1)/\sqrt{m(2m-1)}$. At $m=4$ H$_\#$ itself fails on the slice, at the torus-invariant $\Lambda^2_+$ point.
\end{Lem}

\paragraph{J1 is proved.} J1 is a cubic-quartic inequality in two symmetric zero-diagonal matrices, with equality at the uniform point and at the $\Lambda^2_+$ edge. Its proof reads each row of $U^2$ as a random walk with transition weights $U_{ij}^2/D_i$, $D_i=\sum_jU_{ij}^2$.

\begin{Lem}{274}[J1 at $Q=0$]
For real symmetric zero-diagonal $W\ne0$, $m\ge3$, with $D_i=\sum_jW_{ij}^2$ and $S=\lVert W\rVert^2$,
\[
\tr W^3\le K_m\,\frac{S^2-\sum_iD_i^2}{\sqrt S},\qquad K_m=\frac{\sqrt m\,(m-2)}{(m-1)^{3/2}},
\]
with equality at $J-I$. It contains the traceless cubic bound, and it is J1 on the face $Q=0$.
\end{Lem}

\begin{Lem}{275}[J1 on the face $P=Q$]
For real symmetric zero-diagonal $U$, $m\ge4$, with $x_{ij}=U_{ij}^2$, $\tr U^3\sqrt S\le\lambda_m\sum_{i\ne j}(D_iD_j-x_{ij}^2)$.
\end{Lem}

\begin{Lem}{276}[The decoupled walk inequality]
Let $m\ge7$, $\kappa=\lambda_m^2$, and let $w$ be a symmetric nonnegative $m\times m$ matrix with zero diagonal and total mass $1$, with row sums $\pi_i$, $q_{ij}=w_{ij}/\pi_i$ and $s=\sum_i\pi_i^2$. Then for every choice of $\gamma_{ij}\in[0,1]$, one per ordered pair,
\[
\sum_{i,j}\Big[w_{ij}(1-q_{ij})\big(1+\gamma_{ij}(1-2q_{ij})\big)+\kappa w_{ij}^2(1-\gamma_{ij}^2)\Big]\le\kappa(1-s).
\]
\end{Lem}

\begin{Lem}{277}[J1 for $m\ge7$, and H$_\#$ on the torus slice of $\so(2m)$]
J1 holds for every $m\ge7$. Hence $F(C)\le\kappa\lVert C\rVert^3$ for every torus-invariant $C\in\mathrm{Sym}^2\so(2m)$, $2m\ge14$, and every torus-invariant positive-definite idempotent there has $\lVert C\rVert\ge1/\kappa$.
\end{Lem}

The odd case runs the same chain one size up. Bordering the root matrices by the short-root weights $w_i/\sqrt2$ embeds the slice of $\so(2m+1)$ in that of $\so(2m+2)$.

\begin{Lem}{280}[The cubic form on the torus slice of $\so(2m+1)$]
With $w_i=\mu_{\varepsilon_i}$ on the short root planes,
\[
\lVert C\rVert^2=\lVert A\rVert^2+\lVert U\rVert^2+\lVert V\rVert^2+2\lVert w\rVert^2,\quad F(C)=3\sum_iA_{ii}(d_i+w_i^2)-3\sum_{i\ne j}A_{ij}P_{ij}Q_{ij}+T+3\,w^{\mathsf T}Pw .
\]
\end{Lem}

\begin{Lem}{281}[H$_\#$ on the torus slice of $\so(2m+1)$]
For $m\ge7$ every torus-invariant $C\in\mathrm{Sym}^2\so(2m+1)$ has $F(C)\le\frac{2m-1}{\sqrt{m(2m+1)}}\lVert C\rVert^3$, which is H$_\#$ at odd $n\ge15$.
\end{Lem}

Its inputs are the odd walk inequality (Lemma~278) and the odd J1 (Lemma~279), each the even statement with one distinguished vertex added. Smaller algebras then embed:

\begin{Lem}{282}[Every $\so(k)$ torus slice obeys the ambient bound]
Let $M\ge7$ and $\kappa_{2M}=2(M-1)/\sqrt{M(2M-1)}$. Every torus-invariant $C$ on $\so(2m)$ with $m\le M$, and on $\so(2m+1)$ with $m+1\le M$, has $F(C)\le\kappa_{2M}\lVert C\rVert^3$.
\end{Lem}

\begin{Lem}{283}[The $\su(k)$ torus family sits in $\so(2k)$'s slice]
An operator on $\su(k)$ invariant under $SU(k)$'s torus, extended by zero to $\so(2k)$, is invariant under $U(k)$'s torus, a maximal torus of $SO(2k)$, and keeps its cubic form. So for $k\le M$, $M\ge7$, it obeys $F(C)\le\kappa_{2M}\lVert C\rVert^3$.
\end{Lem}

So \textbf{the torus families of the standardly embedded $\so(k)$ and $\su(k)$ ask nothing more of H$_\#$ at any $n\ge14$}, since $\kappa_n$ increases in $n$ and H$_\#$ at $n$ asks of a proper subalgebra only the ambient constant. Lemmas~260 to 276 and 278 to 283 each carry two independent audits. Lemma~277 is a Lean proof whose statement's fidelity a reviewer confirmed. Lemmas~276 to 283 are proved in Lean on the standard three axioms, Lemma~283 at the level of $\so(2k)$ with the embedding step argued in prose.

\paragraph{Trace floors.} Where the torus count fails, a lower bound on the trace does the work, since pairing $C\#C=C$ with $I$ bounds the norm below by the trace.

\begin{Lem}{287}[An absolute trace floor]
Every nonzero positive semidefinite $\#$-idempotent $C$ on a compact simple $\hh$, in the metric where $I\#I=cI$, has $\tr C\ge h^\vee/c$ and $\lVert C\rVert^2\ge h^\vee/c^2$. In an embedding $\hh\subset\so(n)$ of Dynkin index $j$, $\lVert C\rVert^2\ge j^2/h^\vee$, so where $j^2\ge\tfrac34h^\vee$ and $n\ge9$, H$_\#$ asks nothing of $\hh$, on or off the torus. This covers $\mathfrak f_4$, $\mathfrak e_6$, $\mathfrak e_7$ and $\mathfrak e_8$ in every embedding, $\mathfrak g_2$ in every embedding other than $7$ plus trivial summands, and every $\hh$ in its adjoint embedding.
\end{Lem}

\begin{Lem}{286}[$\mathfrak g_2$'s torus family]
On $\mathfrak g_2$ every $\#$-idempotent with $\tr C\ge\tfrac{12}7$ has $\lVert C\rVert^2\ge\tfrac{36}{49}$, which is H$_\#$'s requirement at every $n\ge9$, and every positive-definite torus-invariant idempotent of $\mathfrak g_2$ has $\tr C\ge2$. So H$_\#$ asks nothing more of $\mathfrak g_2$'s torus family, and the same row argument does the same for $\sype(2)$'s and $\sype(3)$'s.
\end{Lem}

\begin{Lem}{285}[$\sype(k)$'s equal-weight sector]
On $\sype(k)$'s torus family, every positive solution of the idempotent equations whose short weights on $\varepsilon_i-\varepsilon_j$ and $\varepsilon_i+\varepsilon_j$ agree is $I/c$, at every $k$.
\end{Lem}

$I/c$ is unstable on $\sype(k)$ (Lemma~265), so U$_{\mathrm s}'$ holds on $\sype(k)$'s torus family wherever every positive solution has equal short weights. That is proved near $I/c$ and open globally. Lemmas~285 to 287 each carry two independent audits, and their row arguments are machine-checked in Lean.

What H$_\#$ still needs is U$_{\mathrm s}'$ on the classical algebras and on $\mathfrak g_2\subset\so(7)$ at index one, where the trace floor falls short. Three cases are open: the torus families of $\sype(k)$ for $k\ge4$, off the equal-weight sector; an idempotent invariant under no maximal torus of its own algebra; and $n=12,13$, where the $\so(12)$ slice's walk inequality at $m=6$ rests on one machine certificate that no numbered result covers.

\subsection{The mixing gap in the maximal frame}\label{sec:mixinggap}

A second route to Problem~P bypasses the eigenvalue count. In an orthonormal frame maximising the diagonal mass $\lVert\pi_{\mathcal D}W\rVert^2$, with $A=\pi_{\mathcal D}W$ the diagonal part and $\rho$ the critical value, the \emph{mixing gap} is the value dichotomy
\[
\textbf{(MG)}\qquad\rho\le\max\Big\{\sqrt{3/2},\ \frac{P(A)}{\lVert A\rVert^3}\Big\}.
\]
With Lemma~57, which gives $P(A)\le\theta_0(n)\lVert A\rVert^3$ at even $n$, (MG) is Conjecture~A at every even $n\ge12$ with no further input. (MG) and the route are due to an outside author, who reduces (MG) to two conjectures of their own, a closed cyclic trichotomy and a quantitative statement in the nonclosed case, and closes neither. Their exact variance formula is checked here on this report's objects:

\begin{Lem}{284}[The mixing-gap variance identity]
Let $W\in\Weyl_n$ be a unit critical point, $Q(W)=\rho W$ with $\rho\ne0$, and fix any orthonormal frame. Write $W=A+B$ with $A=\pi_{\mathcal D}W$, $t=\lVert A\rVert$, $b=\lVert B\rVert$, $0<t<1$, and split $B=\sum_UB_U$ over the nontrivial sign characters of the frame, $p_U=\lVert B_U\rVert^2$. Put $q_\perp=Q(A)-\tfrac{P(A)}{t^2}A$, write $Q(A,B_U)=\lambda_UB_U+Y_U$ with $Y_U\perp B_U$, and let $\bar\lambda=\sum_Up_U\lambda_U/b^2$. With
\[
z=\frac{A/t-tW}b,\qquad\mathcal U=\tfrac12I+\tfrac12W\otimes W-\rho^{-1}L_W,\qquad\eta=\langle z,\mathcal Uz\rangle,\qquad C=(\mathcal U-\eta)z,
\]
\[
\rho^2t^2b^2\lVert C\rVert^2=\lVert q_\perp\rVert^2+\sum_U\lVert Y_U\rVert^2+\sum_Up_U(\lambda_U-\bar\lambda)^2,\qquad\bar\lambda=\rho t^2\big(\eta+\tfrac12\big).
\]
\end{Lem}

Neither the choice of frame nor global maximality enters. Two reviewers confirmed it independently, and it is proved in Lean over any real inner product space, with the character facts as hypotheses. Its scalar part, $\eta\ge0$ at $\rho>0$, is the frame-split inequality $P(A)\le\tfrac\rho2t^2(1+t^2)$. \textbf{What it adds is the split of $\lVert C\rVert^2$ into three independent nonnegative defects}, so a proof of (MG) must use all three together: a bound on $\lVert q_\perp\rVert$ alone, which separates $P(A)$ from the coupling, is exactly the shape fence post F67 closes. \textbf{The route is blocked on a band of moderate diagonal mass.} Padded $W_{\CP^2}$ lies in it as a local maximum below $\theta_0(n)$ (post F73), so an argument there must use the value $\theta\ge\theta_0(n)$, and the balanced product rotated by $\tfrac\pi4$ across its matched planes sits at diagonal mass $t^2=\tfrac1n$ with value $\theta_0(n)$ (post F81), so only the maximality of the frame, at fourth order across several planes at once, separates small diagonal mass from the maximiser.

\subsection{Conjecture~L16 as a stability constant, and what is settled at its two equality points}\label{sec:stability}

Conjecture~L16 is a quantitative-stability statement, and reading it as one localises what it owes. Writing $\delta^2(W)=\|\mathrm{Alt}(W_-)\|^2$ for the defect and $\kappa_n=9(\sqrt{3/2}-\theta_0(n))<0$ for $n\ge12$, the conjecture asserts $P\le\theta_0(n)+\kappa_n\delta^2$ on the unit sphere, so its constant is
\[
  \kappa^*_n\;=\;\sup_{\delta(W)\neq0}\ \frac{P(W)-\theta_0(n)}{\delta(W)^2}\ \le\ \kappa_n .
\]
The supremum is the largest of three scenario constants: the limit on approach to the balanced orbit $SO(n)\cdot W^*$, the limit at the $\CP^2$-type competitor, and the value at an interior critical point of the ratio. The third \emph{is} the conjecture. The first two are local statements at the points where L16 is already an equality, and both are now settled.

This is the shape of the Bianchi--Egnell constant read backwards for a maximum: the local constant bounds the global one and need not equal it.

\begin{Lem}{133}[The balanced-orbit threshold, in closed form, \emph{audit pending}]
Let $n=2m\ge12$ and let $W^*$ be the unit-norm Weyl tensor of $S^m\times S^m$. Over directions $H\in T_{W^*}S$ with $D(\mathrm{Alt}\,W_-)_{W^*}[H]\neq0$,
\[
  \kappa_n^{\mathrm{Hess}}\;=\;\sup\ \frac{\tfrac12\mathrm{Hess}(P|_S)_{W^*}[H,H]}{\|D(\mathrm{Alt}\,W_-)_{W^*}[H]\|^2}\;=\;\frac{27}{2}\,\theta_0(n)\Big(\frac{2}{n-2}-\frac12\Big),
\]
attained on Lemma~26's direction in every $(2,2)$-split four-subset, and $\kappa_n^{\mathrm{Hess}}\le\kappa_n$ strictly for every $n\ge12$, the comparison reducing to $\theta_0(n)\big(\tfrac3{n-2}+\tfrac14\big)\le\sqrt{3/2}$. So Conjecture~L16 is not refuted at its own conjectured maximizer.
\end{Lem}

The denominator is supported on the $(2,2)$-split four-set blocks alone: $\mathrm{Alt}$ carries no diagonal or pair content, the $(4,0)$-splits are annihilated by the Daleckii--Krein multiplier and the $(3,1)$-splits by the first Bianchi identity, and on what remains Lemma~121's block gives the numerator. The Daleckii--Krein coefficient multiplies precisely the entry $\mathrm{Alt}$ annihilates, and the weight relating signed coordinates to the Frobenius norm cancels between numerator and denominator, so no normalisation convention enters. At $n=6$ the formula returns~$0$, which is where the four-set blocks meet $\tfrac12$ at $S^3\times S^3$.

\begin{Lem}{140}[The defect's first-order response is a $\tfrac12$-contraction, \emph{audit owed}]
Let $H\in\Weyl_n$ with $\|H\|=1$, let $g_1$ be the first-order coefficient of $\mathrm{Alt}(W_-)$ along $V+tH$ at any base point $V$, and let $W_-=\tfrac12(|W|-W)$. Then $\|g_1\|^2\le\tfrac14$.
\end{Lem}

The proof bounds $g_1$ without computing it. Since $H$ satisfies the first Bianchi identity, $\mathrm{Alt}(H)=0$, so $g_1=\mathrm{Alt}(Z+\mu H)$ for every scalar $\mu$. The Daleckii--Krein multiplier of $x\mapsto x_-$ has range of length one, so centring at $\mu=\tfrac12$ halves every entry, and on a kernel block the exact identity $D_-+\tfrac12D=\tfrac12|D|$ gives the same factor. \textbf{The constant is therefore independent of the sign convention for $W_-$}, both conventions having a multiplier range of length one, which matters, because a collision between those conventions is what refuted four earlier lemmas on this route.

\begin{Lem}{141}[The penalised second variation at the balanced orbit, with the sharp constant, \emph{audit owed}]
At $W^*$ with $n=2m\ge12$, the form $Q(H)=\tfrac12\mathrm{Hess}(P|_S)_{W^*}[H,H]-\kappa_n\|D(\mathrm{Alt}\,W_-)[H]\|^2$ satisfies, on $T_{W^*}S\cap T_{W^*}\mathcal O^\perp$,
\[
  Q(H)\;\le\;-c_n\|H\|^2,\qquad c_n\;=\;\frac{3\,\theta_0(n)}{n-2},
\]
and the constant is attained, on the within-factor pair block.
\end{Lem}

$Q$ is additive over the sign-character blocks and the defect response vanishes off the $(2,2)$ four-sets, so the pair block binds. On the four-sets the block ratio is monotone in the share of a direction's weight carried by the defect-active part, and its maximum sits at Lemma~133's own bound. \emph{The intersection with $T_{W^*}S$ is load-bearing: $H=W^*$ is orthogonal to the orbit and gives $Q=+3\theta_0$.}

\begin{Lem}{143}[A dimension-free-in-order bound on the cubic form]
For every $n$ and every $W\in\Weyl_n$ with $\|W\|=1$,
\[
\lvert P(W)\rvert\;\le\;1+\sqrt{n-2}.
\]
\end{Lem}

\begin{proof}
$P$ is odd and the sphere compact with $\nabla P=3Q$, so $\sup\lvert P\rvert=\sup P$ is attained at a critical point, where $\theta\le\lambda_{\max}p+B_5$ with $p=\sum_{\lambda>0}\lambda^2$ and $B_5=\tfrac12(\sum_\alpha\langle T_\alpha\mu,\mu\rangle^2)^{1/2}$. The first term is at most $1$. For the second, the pair cap $\sum_\gamma T_{\alpha\beta\gamma}=\|[v_\alpha,v_\beta]\|^2\le2$ with $\mu_\beta\mu_\gamma\le\tfrac12(\mu_\beta^2+\mu_\gamma^2)$ gives $\langle T_\alpha\mu,\mu\rangle\in[0,2]$, so $\sum_\alpha\langle T_\alpha\mu,\mu\rangle^2\le2\sum_\alpha\langle T_\alpha\mu,\mu\rangle\le4(n-2)$ by the Killing row mass, whence $B_5\le\sqrt{n-2}$.
\end{proof}

\emph{This improves the $\Theta(n)$ bounds above to $O(\sqrt n)$, and $O(\sqrt n)$ is the order that matters below. It does not improve on the marginal linear programme of \S\ref{sec:fenceposts}, whose feasible set it contains and which measures $\approx0.80\sqrt{n-2}$. What it adds is a closed form at every $n$.}

\begin{Lem}{145}[The tube about the balanced orbit]
At $W^*=S^{m}\times S^{m}$, $n=2m$, Conjecture~L16 holds on a tube about $SO(n)\cdot W^*$ of radius $\Omega(n^{-3/2})$.
\end{Lem}

\begin{proof}[Structure of the argument]
The spectrum of $W^*$ is two-valued with no kernel, so the gap at zero is $g_n=\lvert\lambda_-\rvert=\tfrac2n\sqrt{(n-2)/(2(n-1))}=\Theta(1/n)$ and $W\mapsto W_-$ is real-analytic on the ball of radius $g_n$. The Riesz projection onto the negative spectral part is holomorphic on a complex ball of radius $\rho<g_n/2$, and keeping the Frobenius norm on the perturbation alone ($\|ABC\|_2\le\|A\|_\infty\|B\|_2\|C\|_\infty$) gives a Cauchy supremum $M=O(g_n)$, hence $\|A^{(k)}\|=O(n^{k-1})$. The leading remainder term of $(\delta^2)'''=2\langle A''',A\rangle+6\langle A'',A'\rangle$ vanishes at $A(0)=\Alt(W^*_-)=0$, so $C_\delta=O(n)$. With Lemma~141's attained margin $c_n=3\theta_0(n)/(n-2)$ the radius is $\inf_H c(H)/C(H)=\Omega(n^{-3/2})$. \emph{The estimate is conservative in its analyticity radius, and Lemma~147 gives the range it lies in. The exponent is not claimed sharp.}
\end{proof}

\emph{The infimum over directions is what a tube radius is, and it is not the direction-uniform quotient $c_n/\sup_H C(H)$, which is $\Theta(n^{-2})$. Every ingredient is uniform in $H$, and the geodesic segments from the orbit fill the tube, so nothing is extended from rays.}

\begin{Lem}{148}[The defect vanishes identically on the margin summand]
Let $m\ge3$ and $n=2m$. At $W^*$, let $H$ lie in the joint $\lambda_m$-eigenspace (the Weyl image of \emph{both} copies of $S^2_0(\R^m)$ together, of dimension $(m+2)(m-1)$), with the diagonalising frame chosen after $H$. Then $W^*$ and $H$ are simultaneously diagonal on coordinate $2$-planes, hence so are $W(t)=W^*+tH$ and $W(t)_-$, and therefore
\[
\Alt\big(W(t)_-\big)=0\qquad\text{for every }t .
\]
This is an exact algebraic identity on the affine subspace $W^*+\Sigma$, using no spectral gap, no analyticity and no $\ker W^*=0$: $\Alt$ lands in $\Lambda^4$, whose components need four distinct indices, and a coordinate-diagonal operator has none. The branch that vanishes there is the second alone: at $w=0$ the first branch is present and equals $c_p/S\le 3/(n-2)$, and it is what caps the route.
\end{Lem}

\emph{So the summands on which Lemma~27 places the off-orbit margin maximum contribute no defect at all, and the tube radius is decided on the mixed directions: those with a component outside them, whose analyticity radius is $\Theta(g_n)$ rather than $\Omega(n^{-1/2})$. \textbf{Lemma~149 prices exactly those.}}

\begin{Lem}{147}[The radius tracks the cubic form's supremum, and the route is capped]
Write $\sup\lvert P\rvert=\Theta(n^\beta)$. Then $r_n\gtrsim n^{-1-\beta}$, and $\beta$ ranges over the closed interval $[0,\tfrac12]$: $\beta\ge0$ because $\sup\lvert P\rvert\ge P(W^*)=\theta_0(n)\uparrow\sqrt2$, and $\beta\le\tfrac12$ by Lemma~143. So the radius lies in $[n^{-3/2},n^{-1}]$, and $\Theta(n^{-1})$ caps this route.
\end{Lem}

\emph{At $\beta=\tfrac12$ the value $n^{-3/2}$ is confirmed and is pinned by the $P$-side branch alone. At $\beta=0$ the value $n^{-1}$ rests on Cauchy constants that are not pinned, and $\Theta(n^{-1})=\Theta(g_n)$ puts the tube at the order of the analyticity radius itself, which is where those constants decide it. The cap is independent of all of that: the $P$-side branch bounds itself, $c_p/\sup\lvert P\rvert\le3/(n-2)$.}

\begin{Lem}{149}[No direction is both low-margin and high-operator-norm]
Let $n=2m$ with $m\ge6$, and split $H=H_s+H_r$ with $H_s$ in the joint $\lambda_m$-eigenspace (\emph{both} copies of $S^2_0(\R^m)$) and $w=\|H_r\|$. Then $\|H\|_{op}\le a_0n^{-1/2}\sqrt{1-w^2}+w$ and the margin satisfies $c(H)\ge c_p+\rho^{\mathrm{eff}}_nw^2$ with $\rho^{\mathrm{eff}}_n=3\theta_0\tfrac{n-8}{2(n-2)}-\tfrac29|\kappa_n|$, so $\|A''\|=O(n^{1/2}+nw)$ and the defect branch of the radius is $\Omega(n^{-1})$ \emph{uniformly in} $w\in[0,1]$. Hence the bound of Lemma~147 holds over every direction of $T_{W^*}S\cap T_{W^*}\mathcal O^\perp$.
\end{Lem}

\begin{proof}
Writing that branch as $g(w)=\sqrt{c+w^2}/(A+Bw)$ with $c=\kappa/n$, $A=n^{1/2}$, $B=n$, one has $\operatorname{sign}g'(w)=\operatorname{sign}(wA-Bc)$, so $g$ has a single interior minimum at $w^*=Bc/A$, where $g(w^*)=\sqrt c/\sqrt{A^2+B^2c}=\tfrac1n\sqrt{\kappa/(1+\kappa)}$. The endpoints give $\sqrt\kappa/n$ and $\asymp1/n$.
\end{proof}

\emph{The content is that a direction of margin $\Theta(1/n)$ is forced to have $w\lesssim n^{-1/2}$, hence $\|A''\|=O(n^{1/2})$: low margin and high operator norm cannot occur together. The range $m\ge5$ is needed: the smallest off-summand margin, $3\theta_0(m-3)/(2(m-1))$, vanishes at $m=3$ and ties with $c_p$ at $m=4$.}

\paragraph{What the tube does not do.} At either exponent it contains no competitor and covers a vanishing fraction of the sphere, and at $\Theta(n^{-1})$ it is of the order of the analyticity radius itself. \emph{So the cap is a limit on this method rather than on the problem, and the statement it yields is local in a strong sense.}


\paragraph{What the two local scenarios leave.} At the balanced orbit the tube of Lemma~145 is unconditional: Lemma~143 supplies the bound on $\sup\lvert P\rvert$ that the remainder needs, and $O(\sqrt n)$ is what makes the tube unconditional rather than contingent on an unproved constant. \emph{An absolute bound $\sup\lvert P\rvert=O(1)$ would move the radius to $\Theta(n^{-1})$ by Lemma~147, and that is the end of this route rather than a step along it. What decides whether it is reachable is the mixed directions of Lemma~148: the summand carrying the margin contributes no defect at all, so the remainder is governed off it, and no estimate on the record covers those directions. The radius above is therefore a lower bound and not an equality.} At the competitor the corresponding statement reduces, through Lemma~140, to a single bound on the defect's second-order growth. Four formulations of that bound are closed (\S\ref{sec:fenceposts}), and what defeats them is a structural feature of the competitor rather than of the problem: it has a kernel of dimension $\binom n2-3$, so the zero-eigenvalue locus passes through it, while the balanced product has no kernel at all and $W\mapsto W_-$ is real-analytic on a ball about it.

\subsection{Nondegeneracy at the sphere products, proved}\label{sec:nondeg}

Observation~3's nondegeneracy claim is a scan: the eigenvalue that threatens to reach $\tfrac12$ is measured split by split, so it covers no range of $n$ at once. Two lemmas close that, and they run in opposite directions: one exhibits the dangerous eigenvector, the other bounds everything else.

\begin{Lem}{26}[The dangerous eigenvalue of the two-block spectrum, in closed form]
Let $n=p+q$ with $p,q\ge2$, let $\pi$ be a coordinate $2$-plane inside the $p$-block and $\rho$ one inside the $q$-block, and let $X=\phi_\pi\odot\phi_\rho$ with its $\Lambda^4$ component removed. Then $X\in\Weyl_n$ and, for the sphere product's $e$,
\[
e\circ X=\mu X,\qquad \mu=\frac{n-2}{2(p-1)(q-1)},\qquad \mu\ge\tfrac12\iff(p-2)(q-2)\le1 .
\]
So the family reaches $\tfrac12$ exactly at $p=2$, where it gives the excess $\tfrac1{2(n-3)}$, and at $(p,q)=(3,3)$, where it is an equality. For $p,q\ge3$ other than $(3,3)$ the gap is $1/(2\max(p,q)-2)$.
\end{Lem}

\begin{proof}
$X$ is supported on the three pairings of four indices, with coefficients $\tfrac23,\tfrac13,-\tfrac13$ and $\|X\|^2=\tfrac43$. In $T_{ab}=\sum_{m,r,t}c_{amr}\varepsilon_rc_{brt}X_{tm}$ the plane $r$ is determined by $a$ and $m$ as the third side of a triangle, so no sum over the ambient index set survives and the verification is \emph{finite}. Enumerating every plane type over a universe of the four core indices plus two spare indices in each block (enough, since a pair of planes spans at most four indices) gives $e\circ X=\mu X$ identically in $p$ and $q$.

Holding the three block values $\varepsilon_{PP},\varepsilon_{QQ},\varepsilon_{PQ}$ as free symbols, the same computation yields $e\circ X=\lambda X$ for \emph{any} three-block diagonal operator, with $\lambda=\tfrac12(\varepsilon_{PP}+\varepsilon_{QQ})-\varepsilon_{PQ}$. Criticality enters only through $Q(W_{p,q})=\tfrac{n-1}{pq}W_{p,q}$, which fixes the three values and collapses $\lambda$ to $\mu$.
\end{proof}

\paragraph{What it does not reach, and this is the point.} An exhibited eigenvector bounds $\lambda_{\max}$, and one multiplicity, from \emph{below}. The statements nondegeneracy needs (that $(3,3)$ is the \emph{unique} split reaching $\tfrac12$, and that $\dim E(\tfrac12)=pq$) are upper bounds on the rest of the spectrum, which no eigenvector supplies. That gap is what the next lemma closes, by bringing in the representation theory.

\begin{Lem}{27}[Sphere products are Morse--Bott nondegenerate]
Let $p,q\ge3$ with $(p,q)\ne(3,3)$ and $n=p+q$. Then $\lambda_{\max}(L_e|_{W^\perp})=\tfrac12$ at $W_{p,q}$, attained exactly on the orbit tangent, so $\dim E(\tfrac12)=pq=\dim(\text{orbit})$ and $S^p\times S^q$ is Morse--Bott nondegenerate.
\end{Lem}

\begin{proof}
Nienhaus's Proposition 4.6 and the display following it decompose
\[
\Weyl_n=\langle W\rangle\oplus\Weyl_p\oplus\Weyl_q\oplus(\R^p\!\otimes\!\R^q)\oplus S^2_0(\R^p)\oplus S^2_0(\R^q)\oplus(S^2_0\!\otimes\!S^2_0)\oplus(\Lambda^2\!\otimes\!\Lambda^2)\oplus(X_p\!\otimes\!\R^q)\oplus(X_q\!\otimes\!\R^p),
\]
with $X_m$ the $O(m)$-irreducible of Young label $[2,1]$ and $\dim X_m=m(m^2-4)/3$. The ten dimensions sum to $\dim\Weyl_n$ identically in $p,q$.

The Young labels are pairwise distinct for every $p,q\ge3$, so the decomposition is \textbf{multiplicity-free}. Since $e$ is $O(p)\times O(q)$-invariant, $L_e$ is equivariant, and since $L_e$ is symmetric, the $\mathbb C$ and $\mathbb H$ Schur cases cannot occur. So $L_e$ acts as a \emph{scalar} on each summand, and each summand's value \emph{is} an eigenvalue rather than a lower bound for one. That is the step the earlier scan could not supply.

The ten values: $1$ on $\langle W\rangle$; $\tfrac12$ on $\R^p\otimes\R^q$, the orbit tangent; $\lambda_m=(m-2)/(2(m-1))<\tfrac12$ on $S^2_0(\R^m)$; $-\mu$ and $+\mu$ on the two tensor summands, by Lemma~26; and $0$ on the remaining four. Then
\[
\tfrac12-\mu=\frac{(p-2)(q-2)-1}{2(p-1)(q-1)}>0\qquad\text{for }p,q\ge3,\ (p,q)\ne(3,3),
\]
so $\tfrac12$ is attained on the orbit summand alone.
\end{proof}

\paragraph{Both former grants are discharged.} $L_e$ vanishes on $\Weyl_p$ and $\Weyl_q$ by two independent proofs. On $X_p\otimes\R^q$ it vanishes by an index computation \emph{uniform in $p,q,n$}: writing the bilinear sharp as $\mathrm{Sym}_{\text{pair}}\Alt_{ij}H$ and specialising $X=I_p$, Ricci-flatness enters exactly once, to force the partial trace to vanish, and a six-case exhaustion leaves one live case in which the first Bianchi identity is used twice to give $-V/2$. The companion summand follows. The multiplicity-freeness at $p,q\le4$, asserted rather than proved in the source, is supplied by exact $O(m)$-commutant computation at $m=3,4$: every $\mathrm{End}=1$ and every cross-$\mathrm{Hom}=0$, including $\R^3$ against $\Lambda^2(\R^3)$ and $S^2_0$ against $X_3$, which share dimensions. The case $p,q\ge5$ is proved in the source.

\paragraph{The excluded case is a theorem, not an observation.} At $(3,3)$, $\mu=\tfrac12$ exactly and $\dim E(\tfrac12)=18$ against orbit dimension $9$. This is an \textbf{eigenvalue coincidence, not representation mixing}: under $O(3)$, $\Lambda^2(\R^3)\cong\R^3\otimes\det\not\cong\R^3$, measured as $\dim\mathrm{Comm}_{O(3)^2}=2$ against $\dim\mathrm{Comm}_{SO(3)^2}=4$, so Schur permits no mixing and the two summands carry $\tfrac12$ separately. A symbolic solve over all $45$ summand pairs finds it the only such coincidence with $p,q\ge3$, \emph{over all $p,q$} rather than over a tested range.

\begin{certificate}
Audited by four independent reviewers on identical briefs, agreeing on every verdict. Two proved the $\Weyl_p$ vanishing independently, and both refuted the mechanism first proposed for the $(3,3)$ exception. The closed form of Lemma~26 was first found numerically by two further reviewers independently of each other, then proved.
\end{certificate}

The graded classification has one other family, and it behaves better than the sphere products.

\begin{Lem}{28}[The $\CP^m$ family is nondegenerate at every $m\ge3$]
For every $m\ge3$ the $\uu(m)$-invariant line in $\Weyl_{2m}$ is spanned by a critical $W_{\CP^m}$ with $\theta^2=\frac{(m+1)(2m-1)}{3m(m-1)}$, and
\[
\dim E(\tfrac12)=\dim(\text{orbit})=m(m-1),\qquad
\lambda_{\max}\big(L_e|_{W^\perp}\big)=\frac{m+2}{2(m+1)}=\tfrac12+\frac1{2(m+1)} .
\]
So $\CP^m$ is Morse--Bott nondegenerate and, since $\lambda_{\max}>\tfrac12$, is \emph{not} a maximiser.
\end{Lem}

\begin{proof}
$\Weyl_{2m}\otimes\mathbb C$ is $U(m)$-multiplicity-free, so Schur fixes one eigenvalue per summand exactly as in Lemma~27. What replaces the case-by-case evaluation is a closed form. Writing $e$ as a combination of the three $U(m)$-isotypic projectors of $\Lambda^2\R^{2m}$ (with coefficients \emph{rational} in $m$, because the surd in $\theta$ cancels)
\[
e=-\frac1{2m-1}P_{\mm}+\frac{m-2}{(m+1)(2m-1)}P_{\su(m)}+\frac{2(m-1)}{2m-1}P_\omega ,
\]
and using $\sharp(vv^{\top},X)=-\tfrac12D(v)XD(v)$, every block value collapses onto the $U(1)$ charge, so no mixing angle enters and
\[
\lambda_S=1-\frac{c_S+q_S^2}{4(m+1)},\qquad c_S=\sum_i\nu_i(\nu_i+m+1-2i),\quad q_S=\sum_i\nu_i,
\]
with $\nu$ the summand's highest weight as a $GL(m)$ weight. Nondegeneracy is then arithmetic: $\lambda_S=\tfrac12$ exactly when $c_S+q_S^2=2(m+1)$, and across the summands that quantity takes the values $0$, $2m$, $2(m+1)$, $4(m-1)$, $4m$, $4(m+1)$, $4(m+4)$. It holds identically on the orbit. Among the rest only $4(m-1)$ can meet it, at $m=3$ alone, precisely where that summand's dimension $\tfrac{m^2(m+1)(m-3)}4$ vanishes. So $E(\tfrac12)$ \emph{is} the orbit at every $m\ge3$, with no exception. This is cleaner than the sphere products, where $(3,3)$ is genuine.
\end{proof}

\begin{certificate}
The critical row was reached three independent ways: the charge collapse above, a hand assembly $\frac{-(m+4)+3(m+1)+5(2m-1)}{3(m+1)(2m-1)}=\frac2{m+1}$, and explicit highest-weight vectors of weight $\varepsilon_1+\varepsilon_2-\varepsilon_{m-1}-\varepsilon_m$, exact in two primes at $m=3,\dots,9$. Measured confirmation at $m=3,\dots,7$, with $\theta^2$ exact to $m=12$. Audited by two reviewers on identical briefs.
\end{certificate}

\paragraph{The spectrum is derived rather than fitted, and the derivation found something a fit could not.} Writing $c=6\lambda-3$ for the constrained Hessian against the $L_e$ eigenvalue (an identity, since $P$ is cubic), Morse--Bott nondegeneracy reads \emph{$c=0$ on the orbit tangent and nowhere else}. The five families are
\[
c_{A_m}=\frac{-3}{m-1}\ (m=p,q),\qquad c_B=-3+\frac3{p-1}+\frac3{q-1},\qquad c_C=-3,\qquad c_D=-3-\frac3{p-1}-\frac3{q-1},
\]
and they follow from a Casimir master formula, $\lambda=1-\tfrac{\mathrm{Cas}_p}{4(p-1)}-\tfrac{\mathrm{Cas}_q}{4(q-1)}$, obtained by splitting the bilinear form block by block: the three block constants collapse to a common $1-\tfrac1{n-1}$, so \textbf{every dependence on which block a vector sits in cancels}. That derives the standing assertion of \S\ref{sec:landscape} (that each Casimir block carries exactly one Hessian eigenvalue) as a cancellation rather than an observation. Only $c_B$ can vanish, at $\tfrac1{p-1}+\tfrac1{q-1}=1$, i.e.\ $(p-2)(q-2)=1$, whose only solution with $p,q\ge3$ is $p=q=3$.

\begin{remarkstar}[What a fitted family list cannot see, and it was nearly fatal]
Two further families in the ambient $S^2\Lambda^2V$ also have vanishing $c$: $\Lambda^4V_1$ at $c=3(7-p)/(p-1)$, zero at $p=7$, and $\Lambda^3V_1\otimes V_2$ at $c=3(4-p)/(p-1)$, zero at $p=4$. \textbf{Both are killed by the first Bianchi identity}, each occurring once in $S^2\Lambda^2V$ and once in $\Lambda^4V$. Had either survived, $S^7\times S^q$ and $S^4\times S^q$ would be degenerate and Lemma~27 false. A family list obtained by fitting tabulated rows cannot know what is missing from the rows, so the earlier fitted version never ruled them out. \emph{This is the argument for deriving a spectrum even when a fit already reproduces every row on file.}
\end{remarkstar}

\paragraph{Two scope corrections, both narrowing what the lemma may be used for.} \emph{(i)} At $p=2$ the summand $S^2_0(\R^p)$ drops out of $\Weyl_n$ entirely, which is why a ten-summand count overshoots there by exactly $2$. $S^2\times S^q$ \emph{is} nondegenerate, but $c_B=3/(q-1)>0$, so \textbf{it is not a local maximum}: a live exception that a blanket statement over $p,q\ge2$ hides. \emph{(ii)} $c=0$ characterises the \emph{criterion}, not degeneracy. The direction the isolation argument uses is $\ker\mathrm{Hess}=\text{orbit tangent}\Rightarrow$ nondegenerate. The converse fails, so calling $S^3\times S^3$ ``the failure'' overstates: what fails there is the criterion, with kernel $18$ against orbit $9$.

\paragraph{Consequence: the openness half of the component argument fires.} Nondegeneracy is proved across the \emph{whole} graded classification (two-block sphere products for $p,q\ge3$ except $(3,3)$ by Lemma~27, and the $\CP^m$ family with no exception by Lemma~28), so every orbit in it is open in $\Crit(P)$. \textbf{The closedness half is a separate statement and it is where the argument stands}: \S\ref{sec:closedness} is that statement, and what it reduces to. What the component argument then leaves is \S\ref{sec:outside}.

\subsection{Closedness of the graded set: what it reduces to}\label{sec:closedness}

The component argument has two halves. Openness is Lemmas~27 and 28. Closedness was carried for some time as ``$\{S=0\}$ is closed'', and \textbf{that substitution is not available}: $S=0$ is necessary for gradedness and not sufficient, $W_{\CP^2}$ being the standing witness ($S=0$ and ungraded once padded). So the closed set is the wrong set, and the graded set's own closedness has to be argued.

\paragraph{Two facts settle the rival's orbit outright, and neither goes through $S$.} Two independent implementations, on identical briefs, agree on both.

\begin{enumerate}[leftmargin=1.8em,nosep]
\item \textbf{$W_{\CP^2}$'s orbit is a union of connected components of $\Crit(P)$}, from the Hessian. Padded $W_{\CP^2}$ is Morse--Bott nondegenerate: $\ker\mathrm{Hess}$ equals the orbit exactly at $n=6,8,10$, excess $0$ ($10=10$ at $n=6$ and $18=18$ at $n=8$, measured separately by each implementation). So $\Crit(P)$ is \emph{locally exactly the orbit} there.
\item \textbf{$S$ is identically zero on $\Crit(P)$ there, and does not jump.} Perturb-and-project gives $S\approx10^{-46}$ after projection against $20.6$--$44.9$ on the sphere, every point landing on the orbit at distance $\le7.4\times10^{-12}$ against raw distance up to $0.048$. Since $S$ is $O(n)$-invariant and $\Crit(P)$ is locally the orbit, $S$ is \emph{necessarily} constant there. \emph{The jump is a sphere phenomenon}: perturbations taken on the unit sphere with no projection back to the critical set do see one, and $\dim\mm$ jumps $2\to14$ at $n=8$ there because padding contributes $25$ zero eigenvalues. On $\Crit(P)$ there is nothing to jump.
\end{enumerate}

\paragraph{The general statement reduces to a spectral gap, and the reduction is the usable form.} Gradedness can fail in a limit only by an eigenvalue reaching zero. Hence
\[
\text{the graded critical set is closed}
\iff
\text{a lower bound on }\min|\lambda|\text{ over it.}
\]
That is checkable and needs no enumeration, which is what recommends it. \textbf{Two scope corrections are forced before it can be used, and together they change what must be proved.}

\begin{itemize}[leftmargin=1.5em,nosep]
\item \textbf{It must be the smallest \emph{nonzero} modulus.} A padded graded critical point has $\min|\lambda|=0$ exactly ($S^3\times S^3$ padded to $n=8$ and $n=12$ carries $13$ and $51$ zero eigenvalues, $S^4\times S^4$ padded to $n=12$ carries $38$), and paddings of graded points are graded, so the raw infimum is $0$ and the statement is vacuous. Padding leaves the \emph{nonzero} spectrum untouched, which is what makes the repair salvageable: the smallest nonzero modulus at $S^3\times S^3$ is $0.210818511$ at $n=6$, $8$ and $12$ alike. The quantity is the gap of $W$ on $(\ker W)^{\perp}$, padding-invariant by construction.
\item \textbf{It cannot be uniform in $n$, and the conjectured maximizer is the witness.} At the balanced product the smallest nonzero modulus is $\sqrt2(n-2)/\big(n\sqrt{(n-1)(n-2)}\big)\sim\sqrt2/n\to0$, exactly at $n=6,8,10,12,16$ (deviation $\le1.4\times10^{-17}$). Any usable statement is \emph{per $n$}. This is fence post~27.
\end{itemize}

\begin{Obs}{56}[At fixed $n$, the named graded families have a floor, and it is not the set's floor]
Among the products, the smallest gap at fixed $n$ is attained at the most unbalanced factorisation: $0.062994$ at $S^3\times S^7$ for $n=10$, $0.038925$ at $S^3\times S^9$ for $n=12$, with the balanced product (the conjectured maximizer) at the \emph{largest} gap of the family. The mechanism is elementary: the vertex entries of $S^p\times S^q$ are $-q/(p-1)$, $-p/(q-1)$ and $1$, so the smallest modulus is $p/(q-1)$, i.e.\ $3/(n-4)$ at $p=3$ before normalisation. Multi-block products interleave rather than extending the trend, $S^3\times S^3\times S^3\times S^3$ at $0.0580$ sitting between $S^3\times S^9$ and $S^4\times S^8$.
\end{Obs}

\paragraph{The candidate falsifier that stood here is dead, and it died of a sign convention.} A critical point at $n=6$ inside $\sigma=\mathrm{diag}(1,1,1,-1,-1,-1)$, at $\theta=0.418198$ with gap $2.2\times10^{-3}$, was recorded as having $S=0$ and so as sitting a hundredfold below $S^3\times S^3$'s gap. Normalized to $\theta>0$ it has $S_{\mathfrak{mmm}}=21.5213$ and $S_{\mathfrak{hhm}}=7.1738$, with $\dim\mm=10$: \textbf{it fails the first condition outright and never needed the second}. The subspace dimension $45$ and the $355$ of $400$ convergences stand, but every $S$ that hunt reported does not.

\textbf{The mechanism is general and worth stating once.} $\theta$ is cubic, so $W$ and $-W$ are both critical with opposite $\theta$, and under the flip $\hh$ and $\mm$ exchange, and every block mass with them. A hunt reporting $|\theta|$ while measuring $S$ on whichever sign its solver returned has its two columns in different conventions. One line exhibits it: at $S^3\times S^3$, sign $+1$ gives $\theta=-1.054093$ with $S_{\mathfrak{mmm}}=12$, and sign $-1$ gives $\theta=+1.054093$ with $S_{\mathfrak{mmm}}=S_{\mathfrak{hhm}}=0$.

\begin{certificate}
Adjudicated by a third implementation computing every block mass as a \emph{matrix commutator} in $\so(n)$ under $\langle A,B\rangle=\tfrac12\tr(A^{\mathsf T}B)$, sharing no machinery with the hunt, agreeing to every printed digit with an independent implementation reached separately. Two controls and a sum rule run every time: $S^3\times S^3$ must come out graded ($S_{\mathfrak{mmm}}=S_{\mathfrak{hhm}}=0$ exactly, $\dim\mm=9$); padded $W_{\CP^2}$ must come out ungraded at $\theta=\thstar$ with $S_{\mathfrak{mmm}}\approx0$, $\dim\mm=2$ and $S_{\mathfrak{hhm}}=4(n-4)$ \emph{exactly} ($0,4,8,12,16,20$ at $n=4,\dots,9$, a closed form rather than a measured figure, since without a kernel tolerance the same operator gives $\dim\mm$ anywhere in $\{3,\dots,9\}$ according to which LAPACK routine is called); and $S_{\mathfrak{mmm}}+3S_{\mathfrak{mmh}}+3S_{\mathfrak{mhh}}+S_{\mathfrak{hhh}}=2N(n-2)$, from total antisymmetry and the row marginal $2(n-2)$, residual $\le2\times10^{-13}$ at all $355$ points. \textbf{Both controls failed on the first build, and that is how both bugs were found}: a grading test passes silently in either convention, so nothing in the measured output would have shown either.
\end{certificate}

\paragraph{A second convention trap, in the $\hh/\mm$ split itself.} A padded operator has a large exact kernel and \texttt{eigh} returns those eigenvalues at $\pm10^{-17}$, so a bare $\lambda<0$ test scatters half of it into $\mm$. Padded $W_{\CP^2}$ at $n=6$ then reports $\dim\mm=5$ and $S_{\mathfrak{mmm}}=2.64$ against the true $2$ and $3\times10^{-64}$. \textbf{Zeros belong to $\hh$, and the split needs a tolerance to put them there.}

\paragraph{Across four gradings the search finds no graded point at all, and the control shows that is a fact about the search.} Sweeping $\sigma(1,5)$, $\sigma(2,4)$, $\sigma(3,3)$ and the complex structure $J$, invariant subspaces of dimension $49$, $44$, $45$, $48$, at $400$ random starts each: $1493$ converged critical points, \emph{zero} passing the full grading test.

\begin{certificate}
\textbf{The zero is forced, and the mechanism is degeneracy.} Inside $\sigma(3,3)$ the criticality map's Jacobian has $\ker J=9$ against an orbit of $0$ at $S^3\times S^3$, excess $9$, and $\ker J=6$ against an orbit of $6$ at every random-start point, excess $0$. \textbf{Every graded point at $n=6$ is a degenerate zero and every ungraded one found is nondegenerate}, so a residual-cutoff root-finder cannot reach the former: $0$ of $1205$ is structural, and no amount of compute changes it. Two audits on identical briefs, one reporting excess $0$ at all $122$ of its own ungraded points, and reproduced here independently.
\end{certificate}

\paragraph{This rescopes post~13 rather than repeating it, and it makes isotropy enumeration necessary.} Random starts reach the top of $\theta$ perfectly well: under gradient \emph{ascent} they reach $\theta_{\max}=\thstar$ from $100$ of $100$ starts in full $\Weyl_6$, and $59$--$60$ of $60$ per cone. What fails is \textbf{root-finding on the criticality equation}, which is blind to degenerate zeros by construction. And ascent does not rescue the census either: $S^3\times S^3$ is a \emph{degenerate saddle}, with zero positive Hessian directions but $\theta$ taking both signs at cubic order on the $9$ null directions transverse to its orbit. \textbf{Neither general-purpose method can reach a graded point}, so a search over isotropy subalgebras is not a convenience but the only route left.

\textbf{What the sweep does settle.} Twenty-four points across three cones have $S_{\mathfrak{mmm}}$ at $10^{-30}$ with $S_{\mathfrak{hhm}}$ nonzero ($17.0000$ at the smallest-gap one, $\theta=0.104886$), carrying \emph{no kernel and no padding}. The only prior witness for ``$S=0$ does not imply graded'' was padded $W_{\CP^2}$, whose ungradedness is manufactured by the padding, and both audits of it said so. \emph{That objection is answered.}

\textbf{So the floor question is blocked on method rather than on evidence.} A graded critical point is one with a symmetry, so a search over isotropy subalgebras reaches them by construction where a random start provably cannot. No such search has been run.

\begin{center}
\begin{tabular}{@{}p{0.40\linewidth}p{0.56\linewidth}@{}}
\toprule
\textbf{Step} & \textbf{State} \\
\midrule
Graded orbits are open in $\Crit(P)$ & \textbf{Proven} across the classification (Lemmas~27, 28) \\
$W_{\CP^2}$'s orbit is a union of components & \textbf{Proven} from the Hessian at $n=6,8,10$, excess $0$ \\
$S\equiv0$ on $\Crit(P)$ near that orbit & \textbf{Proven}. The jump is a sphere phenomenon and not one on $\Crit(P)$ \\
``$\{S=0\}$ closed'' as a substitute for closedness & \textbf{False}: $S=0$ is necessary for gradedness and not sufficient \\
Closedness $\iff$ a nonzero-gap bound, per $n$ & \textbf{Proven as a reduction}. The bound itself is \textbf{open} \\
A uniform-in-$n$ bound & \textbf{Refuted}: the balanced product's gap is $\sim\sqrt2/n$ \\
A per-$n$ floor from the named families & \textbf{Verified} for products. Whether it describes the whole set is \textbf{untouched} \\
The $n=6$, gap $2.2\times10^{-3}$ point as a falsifier & \textbf{False.} Its recorded $S=0$ was $S$ on $-W$. At $\theta>0$ it has $S_{\mathfrak{mmm}}=21.5213$ and fails the first condition \\
A census of the graded set by root-finding & \textbf{Closed structurally.} Graded points are \emph{degenerate} zeros, excess $9$ against $0$ at every ungraded point, so a residual-cutoff solver cannot reach them at any compute \\
A census by gradient ascent & \textbf{Closed too.} $S^3\times S^3$ is a degenerate saddle, so ascent passes it by. Only isotropy enumeration remains \\
\bottomrule
\end{tabular}
\end{center}

\emph{Closedness may not be on the critical path at all}: the compactness step needs only that the graded set be open, which is Lemmas~27 and 28. That is the reason the value bound of \S\ref{sec:outside} is the route this report pursues.

\paragraph{The spectral gap's floor is at $p=2$, and it has a closed form.} Two reviewers, independently and on identical briefs, place the attainment point one block lower than a $p\ge3$ reading suggests. For $p<q$ the gap is
\[
\mathrm{gap}(p,q)=\sqrt{\frac{2p(p-1)}{q(q-1)(n-1)(n-2)}},
\]
strictly increasing in $p$, so the minimum sits at $p=2$ at every $n\ge6$, by monotonicity, with no sweep. $S^2\times S^{n-2}$ is an admissible Weyl operator and is where the floor is attained, so a floor quoted at $p=3$ is reading the closed form off the wrong block.

\subsection{The residual, and what it asks for}\label{sec:outside}

With nondegeneracy proved across the graded classification, each graded orbit is open in $\Crit(P)$. Granting the closedness of \S\ref{sec:closedness}, the graded set is then a union of connected components. The question that remains is what the rest of $\{S=0\}$ looks like near $\theta_0(n)$, and the target is a \textbf{value bound} rather than a classification (a classification being both unavailable and unnecessary, for the reasons below). \emph{The value bound does not depend on the closedness}, which is why it is stated first and pursued on its own.

\paragraph{The split that makes the residual small.} $\{S=0\}$ divides by coordinate-diagonality, and \textbf{the diagonal half is granted by the problem statement} for even $n\ge12$: that is exactly what Lemma~1 secured at the threshold dimensions. So what is owed is one inequality on the other half:
\[
\sup\big\{\theta\;:\;Q(W)=\theta W,\ S=0,\ W\text{ non-diagonal}\big\}\;\le\;\thstar .
\]
\textbf{That closes every even $n\ge12$ on its own}, since $\thstar<\theta_0(n)$ there. Every known such point respects it, the closest being the rival $\mathcal S^4(\alpha)$ at exactly $\thstar=0.9909\,\theta_0(12)$, the same orbit the localization strategy localizes on. Its falsifier is a single object: a non-diagonal $S=0$ critical point above $\thstar$.

\paragraph{The statement is well posed because diagonality is decided frame-independently.} Tested in a chosen frame, ``non-diagonal'' names a set that depends on the frame, and the supremum above has no definite meaning. The criterion below removes the dependence.

\begin{Lem}{37}[Diagonality is decided by a single kernel]
For $W\in\Weyl_n$ let $\delta_A(X)=AX+XA$ for symmetric $A$, and let $\Phi_W(A)=[W,\delta_A]$. Then $W$ is diagonal in \emph{some} orthonormal frame if and only if some $A\in\ker\Phi_W$ has pairwise-distinct sums $a_i+a_j$. In particular $\dim\ker\Phi_W<n$ certifies non-diagonality outright.
\end{Lem}

\begin{proof}[sketch]
$W$ is diagonal in the frame diagonalising a symmetric $A$ with distinct eigenvalue sums exactly when $W$ commutes with $\delta_A$, since the $a_i+a_j$ are then the distinct eigenvalues of $\delta_A$ on $\Lambda^2$ and its eigenspaces are the coordinate $2$-planes. Commuting is $\Phi_W(A)=0$. Conversely a diagonalising frame produces such an $A$. The dimension certificate follows because a diagonal $W$ admits an $n$-parameter family of such $A$.
\end{proof}

\begin{remarkstar}[Two things the criterion does not do]
$\delta_{\Id}(X)=2X$, so $\Id\in\ker\Phi_W$ always and $\dim\ker\Phi_W=1$ forces $\ker\Phi_W=\R\cdot\Id$ \emph{exactly}: $\CP^3$ and $\CP^5$ have literally the same kernel. That value is generic: it holds at every $\CP^m$ and at random Weyl tensors whose $\theta$ differ by an order of magnitude, so the invariant separates diagonal from non-diagonal cleanly and separates nothing else. This is fence post~22: \textbf{use $\ker\Phi_W$ to define the target set, not to bound $\theta$ on it}. And the dimension corollary alone is not enough: implemented by itself it wrongly calls padded $W_{\CP^2}$ diagonal, contradicting that point's own rank-$4$ simple eigenvector, so the distinct-sums search is load-bearing.
\end{remarkstar}

\paragraph{It is diagonality that carries the bound, not $S$.} For any symmetric pair $(\mathfrak g,\kk)$ inside $\so(n)$ the two-eigenvalue point $a\Pi_{\kk}+b\Pi_{\mm}$ has $S=0$ identically, since $[\mm,\mm]\subseteq\kk\perp\mm$. Taking $(\so(2k),\so(k)+\so(k))$ returns \emph{the Weyl tensor of $S^k\times S^k$}, at $\theta_0(2k)^2=2-3/k+1/k^2$, which exceeds $\thstar$ from $k=6$ on and tends to $\sqrt2$. \textbf{So $\{S=0\}$ contains points running up to $\sqrt2$, and what keeps them out of the target is that they are diagonal.} An argument that uses only slack zero is bounding a set that reaches $\sqrt2$, so the diagonality hypothesis has to be carried to the end, and Lemma~37 is the tool for carrying it.

\paragraph{Why completeness was the wrong target.} Enumerating the non-diagonal $S=0$ points cannot be done and need not be: unitary isotropy is sufficient for membership but never necessary, and orthogonal direct sums manufacture an infinite unlisted family with $S=0$ (five verified at $n=12$, all at $0.65$--$0.70\,\theta_0$). A value bound does not care how many there are.

\textbf{Sums can be dropped from the search entirely, and the reason is one line.} $\circ$-orthogonality is automatic across disjoint blocks, since $[\so(n_1),\so(n_2)]=0$ kills the polarised $\#$; $Q$ maps block-supported operators into the same block, so criticality of the sum \emph{forces} criticality of each summand rather than assuming it; and $S(W)=S(W_1)+S(W_2)$. Hence $\theta^{-2}=\theta_1^{-2}+\theta_2^{-2}$ and $|\theta|<\min_i|\theta_i|$ strictly. \textbf{So the supremum over the non-diagonal $S=0$ set is attained at block-indecomposable points.}

\paragraph{Under the criterion the target was censused, and it is attained rather than merely satisfied.}

\begin{Obs}{40}[The surviving target, measured]
Over the catalogued critical points, decided by Lemma~37 rather than by frame inspection, $\sup\{\theta:Q(W)=\theta W,\ S=0,\ W\text{ non-diagonal}\}=\thstar$, \textbf{attained at $W_{\CP^2}$ and at every padding of it} (margin $-6.7\times10^{-16}$). Diagonal controls are correctly excluded, $S^5\times S^7$ at $1.228239$ among them. Three independent hunts agree: a $26$-space symmetric-space sweep, and closure sweeps at $n=6$ and $n=8$ both topping out at padded $\CP^2$.
\end{Obs}

\begin{certificate}
Double precision, with every row passed through a criticality screen that refuses non-critical input, a fix forced by an earlier table in which three published rows sat at criticality residual $0.11$--$0.22$, where the conditions being tested have no content. The diagonality verdict is integer-valued where the dimension and rank certificates decide it. Where they are silent the distinct-sums search decides it, and its failure certifies non-diagonality with probability one, the distinct-sum locus being Zariski-open. \emph{That last branch is load-bearing rather than a fallback}: $\CP^2+\CP^2$ embedded at $n=12$ has no simple eigenvalue and $\dim\ker\Phi_W=n$ exactly, so both certificates are silent, and a census folding the undecided case into ``diagonal'' would drop such points from the target set.
\end{certificate}

\paragraph{What that census does and does not license, after audit.} \textbf{Its bound survives, but its characterisation is refuted.} The claim ``the non-diagonal $S=0$ points are $\CP^m$, its paddings and $\circ$-orthogonal sums'' is false, and at least two further families lie outside it:

\begin{longtable}{@{}p{0.30\linewidth}p{0.08\linewidth}p{0.24\linewidth}p{0.30\linewidth}@{}}
\toprule
\textbf{Point} & $n$ & $\theta$ & \textbf{Why it escapes the list} \\
\midrule
\endhead
$\HP^2=Sp(3)/(Sp(2)\!\cdot\!Sp(1))$ & $8$ & $\sqrt{56/45}=1.115547$ & Triality at $n=8$: isospectral to $S^3\times S^5$, one diagonal and one not \\
\addlinespace
$(\su(m),\so(m))$, $m\ge3$ & $2m$ & $1.095445$, $1.035098$, $1.024695$ at $m=3,5,6$ & A K\"ahler family the list never contained, capped in closed form by Lemma~38 \\
\bottomrule
\end{longtable}

$\HP^2$'s certificate is analytic rather than numerical: its top eigenspace is spanned by the quaternionic structures $I,J,K$, and every non-zero $aI+bJ+cK$ is an invertible quaternion, hence of bivector rank $8$. The record's earlier rank certificate was structurally blind to this, inspecting only \emph{simple} eigenvalues where this one has multiplicity $3$. $\HP^2$'s own analytic cap is Lemma~47's, together with the closed form $\theta^2(\HP^m)=\tfrac{(m+2)^2(4m-1)}{9m(m-1)(2m+1)}$, decreasing from $\tfrac{56}{45}$ (\S\ref{sec:symmetric}).

\paragraph{The sharp case is excluded with no enumeration at all.} Writing $S(W)=\sum_{\lambda<0}c'^2_{\alpha\beta\gamma}$, total antisymmetry gives $\dim\mm\le2\Rightarrow S=0$ identically, since there are no triples. Padded $W_{\CP^2}$, the sharp point at $\thstar$, has spectrum $(2,-1,-1)/\sqrt6$ and lives in precisely that stratum. Maximising $\theta$ over $\Weyl_n$ with the third-most-negative eigenvalue driven to zero returns \textbf{exactly $\thstar$ at $n=6,8,10,12$}. Since $\thstar<\theta_0(12)=1.236033$:
\[
\textbf{at }n=12\textbf{ a maximizer cannot have fewer than three negative eigenvalues.}
\]
Padded $W_{\CP^2}$ is the only catalogued point in that stratum (every other $S=0$ point has $\dim\mm\ge6$), so the nearest competitor, at $0.9909\,\theta_0(12)$, is excluded without listing anything. In the complementary stratum the best found is $1.0560$ at $n=6$ and $1.1315$ at $n=8$, well under $\thstar$, though only one start in three or four cleared both constraints, so those are lower bounds rather than suprema.

\paragraph{Openness holds at every class that has been exhibited.} Five classes of diagonal $S=0$ point are covered: the sphere products by Lemma~27; padded points, nondegenerate at the top, with $E(\tfrac12)$ \emph{equal} to the orbit at $S^6\times S^6$ padded $12\to14$ and $S^5\times S^5$ padded $10\to12$ (principal cosines $1.000000000000$); multi-block products, which carry $S>0$ by the triangle count below; the subalgebra family, whose only $S=0$ member is the balanced product; and the $C_5$ blow-up, settled by the third-order obstruction.

\paragraph{The slack is a triangle count.} For multi-block products,
\[
S=3\sum_{a<b<c}n_an_bn_c=3\times\#\{\text{triangles of the negative block graph}\},
\]
exact on all $437$ products with $n\le20$ and on $62$ partitions to $n=34$. \textbf{So $S=0$ exactly when that graph is triangle-free.} Multi-block \emph{sphere} products have complete negative graphs, so they never achieve it, making $S>0$ for $k\ge3$ a theorem rather than a measurement. The blanket statement over all multi-block critical points is nonetheless false: of $3767$ scanned at $n=9,\dots,16$, forty-one are indecomposable with $S=0$, all at $k=5$ and all $C_5$ blow-ups. The reason none occurs at $k=3$ or $4$ is structural: the negative graph must be connected, triangle-free and non-bipartite, and $C_5$ is the smallest such graph. Being $5$-block, they sit below $\theta_0/2$.

\paragraph{One classification claim withdrawn, at no cost to the route.} ``Every diagonal $S=0$ critical point is a blow-up'' is \textbf{false}: at $n=8$ two are \emph{circulants} on the $8$-cycle, with $w$ depending only on cyclic distance, at $\theta=0.5557$ and $0.0823$, residual $8\times10^{-16}$. Blow-ups are the imprimitive case. The right invariant is vertex-transitivity, not a block partition. Both values sit far below $\theta_0(8)=1.145644$, so this costs the classification and not the route, which is the second time the same lesson has been paid for, and the reason the target is now a value bound.

\paragraph{Decomposable points, and the cap they need.} The harmonic-square law's real hypothesis is $\circ$-orthogonality, $W_A\circ W_B=0$, not disjointness of index blocks, and it \emph{implies} Frobenius orthogonality rather than assuming it. Where it applies, criticality forces $\theta=a\theta_A=b\theta_B$ with $a^2+b^2=1$, whence
\[
\frac1{\theta^2}=\sum_i\frac1{\theta_i^2},
\]
verified at twenty-five held-out block sums to $9\times10^{-16}$: three and four blocks, unequal dimensions, decomposable blocks, scattered coordinates, ambient padding. Hence $|\theta|<\min_i|\theta_i|$, the absolute values being necessary since $(+\mathrm{Fox})\oplus(-\mathrm{Fox})$ at $n=8$ gives $\theta=+0.866$ against $\min_i\theta_i=-1.2247$. Every $\circ$-summand of a critical $W$ lies in $E(1)$ of $L_e$, so $\dim E(1)=1$ \emph{proves} indecomposability: it holds at the sphere products, at $\CP^m$, at padded $W_{\CP^2}$ and at the $C_5$ blow-up, while every direct sum splits. What is owed here is an a priori $M(n)\le\sqrt3$, which would give $|\theta|\le M(n)/\sqrt2\le\thstar$ at every decomposable point.

\paragraph{A refuted constant, kept because it shows what the hypothesis costs.} The bound this was pursued for (``every decomposable critical point has $\theta\le\sqrt3/2$'') fails twice over. Its hypothesis, $\theta_i\le\thstar$ for $n_i\le11$, \emph{is the open low-dimensional half of Conjecture~A}: proved outright only at $n_i=4$, where $P=9(\det W^++\det W^-)$ on $\Weyl_4$ gives $\max P=\thstar$ exactly. And the constant is false from $n=16$: $S^6\times S^6$ on twelve coordinates plus $W_{\CP^2}$ on four gives $\theta=0.869988927484$ at residual $2.7\times10^{-16}$, with a second witness at $n=20$ reaching $0.885202$. What survives is weaker in hypothesis and stronger in reach: $\theta/\theta_0\le0.706$ for every $n\ge12$, requiring only $\theta_i<\sqrt2$.

\paragraph{$S$ is not a distance, and the trap is worth one sentence.} \textbf{The multi-block products cannot be excluded on the size of their slack}, which is a non sequitur: perturbing padded $W_{\CP^2}$, an exact $S=0$ point, by $10^{-6}$ gives $S\in[11.3,22.2]$, so openness fails through a \emph{sequence} and the value of $S$ at any one point carries no topological information. What excludes those points is $\theta$: $P$ is constant on connected components of $\Crit(P)$, and $\sup\theta/\theta_0$ over $k$-block products is $1/\sqrt{k-1}$.

\paragraph{The one uniformity that is not available.} Any padding of a seed from $n_0$ gives $\min_{\lambda\ne1/2}|\lambda-\tfrac12|=\tfrac1{2(n_0-1)}$, \emph{independent of the target dimension}, agreed to seven digits by three methods. Padding the balanced product freezes $\theta$ while the ambient bound rises, so $\theta/\theta_0\to1$ and the gap $\to0$ together. This costs less than it appears: Conjecture~A is a per-dimension statement, at fixed $n$ only the seeds $n_0\le n$ contribute, all with gap $\ge\tfrac1{2(n-1)}$, and compactness supplies a constant $c_n>0$ that may depend on $n$. \textbf{No argument may rest on a spectral gap or a $\theta$-margin uniform in $n$}. It is a constraint on the shape of the argument, not on its conclusion.



\paragraph{The $\mathrm{ad}$-form, and the sub-question it leaves behind.} Wilking's Lie-algebraic form
\[
\langle M^{\#}x,y\rangle=-\tfrac12\tr\big(\mathrm{ad}_xM\,\mathrm{ad}_yM\big),\qquad x,y\in\so(n),
\]
is the operator-level handle fence post~33 asks for. Cauchy--Schwarz on the trace closes nothing: it produces a norm-only bound, and those overshoot by $3.25\times$: separate Okumura caps on the two cubics reach $4.0129$ at $n=12$ against $\theta_0(12)=1.236$. \textbf{The reason is that the bound never uses that $M$ is a Weyl tensor}: $M$ is not a general element of $S^2(\so(n))$, and taking norms discards the first Bianchi identity along with the $\Lambda^4$ component it kills.

\emph{So what is owed is not a sharper inequality on the same data.} It is an identity: express $\langle M^{\#},M\rangle$ through the $\mathrm{ad}$-form so that the $\Lambda^4$ component \emph{separates}, making Bianchi visible as the vanishing of a named term rather than an unused hypothesis. That is a statement about the structure constants of $\so(n)$ paired against a tensor of known symmetry type, checkable by hand once found. It is the right shape of obstruction: post~33 rules out two moments with a support bound, and post~23 rules out any \emph{linear} identity coupling the two spectra's invariants at $n\ge6$ while leaving inequalities open.

\subsection{The degree-six invariant certificate, settled at $n=5$}\label{sec:sos}

The named fallback is a sum-of-squares certificate for
\[
G_{\mu,n}(W)=\Theta_n^2\|W\|^6-\langle Q(W),W\rangle^2+\mu\big(\|W\|^2\|Q(W)\|^2-\langle Q(W),W\rangle^2\big),
\]
whose second bracket is $\|W\wedge Q(W)\|^2$ and vanishes at every critical point. A certificate at one dimension settles Conjecture~A there. \textbf{At $n=5$ the certificate exists and is exhibited.}

\begin{Lem}{61}[$M(5)=\sqrt{3/2}$, the global maximum on $\Weyl_5$]
$F_5=\tfrac32\|W\|^6-9\operatorname{tr}(W^3)^2$ is a sum of squares of cubics vanishing on the $O(5)$-orbit of $W_0$, the padded $W_{\CP^2}$. Hence $\lvert\theta\rvert\le\sqrt{3/2}$ on $\Weyl_5$, with equality on that orbit. The equality set is not identified.
\end{Lem}

\begin{proof}[sketch]
Three facts compose. \emph{The target is reachable}: the map $\mu$ from Gram matrices to sextics has rank $13$ on the orbit-constrained space, so $\dim\mathrm{ANN}=1$, the single obstruction is $A(W_0^6)$, and $F_5\perp A(W_0^6)$ because $F_5(W_0)=0$. \emph{The witness exists}: an explicit rational Gram, positive definite. \emph{Agreement is identity}: every coordinate of the Gram-to-sextic map is an $O(5)$-invariant sextic, as is $F_5$, and both lie in the $14$-dimensional $(\mathrm{Sym}^6\Weyl_5)^{O(5)}$, on which evaluation at the sample points is injective: $120$ explicit invariant sextics have rank $14$ against a cap of $14$ there. So a Gram reproducing $F_5$ at those points reproduces it identically.
\end{proof}

\begin{certificate}[Exact, at every step]
The rows of the system are computed modulo four primes near $2^{40}$ and lifted to $\mathbb Q$ by the Chinese remainder theorem and rational reconstruction, unique by a margin of $2.26\times10^{-14}$. The family is an exact particular solution plus an exact $34$-dimensional kernel, each checked against all $40$ equations. The certificate is a rational combination of those, hence a member by construction. Positive definiteness is verified by two routes sharing no code (Sylvester minors, and characteristic polynomials with alternating coefficients), and holds for both packings of the off-diagonals, with $\lambda_{\min}>21/1000$ certified over $\mathbb Q$. No floating-point step enters the verification.
\end{certificate}

\paragraph{The bound is sharp and attained, not slack.} $F_5\ge0$ is exactly $|\operatorname{tr}(W^3)|\le\|W\|^3/\sqrt6$, a maximisation on the $35$-sphere. Three thousand Riemannian gradient ascents from random starts, both signs, all reach $0.4082482904639=1/\sqrt6$, at one critical value with negative semidefinite Hessian. The maximiser's $\Lambda^2$ spectrum is $(-2,1,1,0^7)$. Over $\mathbb Z$: $\|W_0\|^2=24$, $\operatorname{tr}(W_0^3)=48$, and $6\cdot48^2=13824=24^3$.

\paragraph{The block structure, and why separation is a theorem.} $\mathrm{Sym}^3(\mathcal W_5)$ carries $24$ blocks under $O(5)$, and \textbf{the \emph{pair} $(C_2,C_4)$ separates all $24$, which is forced, with no table and no case check}. By Harish-Chandra, $Z(U(\so(5)))\cong S(\hh)^{W(B_2)}=\R[p_2,p_4]$. For dominant $\lambda$ the shift $\ell=\lambda+\rho$ is \emph{strictly} dominant automatically, $\rho=(\tfrac32,\tfrac12)$ having strictly positive decreasing entries, so a regular $\ell$ meets the open chamber once and the infinitesimal character separates every dominant $B_2$ weight. Exactly two power sums suffice, of degrees $2$ and $4$. Beyond $p_4$ Newton makes the rest redundant.

\textbf{No single quartic is distinguished, so ``the quartic generator'' names nothing.} $z_4'=z_4+\alpha z_2^2+\beta z_2+\gamma$ generates for \emph{every} choice of constants, and the choices differ in what they separate on their own: at $\alpha=-\tfrac12$ the image takes only $18$ values on the $24$ blocks. \textbf{And $C_2^2$ does not separate nothing}: it separates exactly the $21$ classes $C_2$ itself does, $c_2$ taking $21$ distinct values with three collisions, at $30$, $40$ and $60$. \emph{What is forced is the pair, not a degree and not any one quartic.} The trivial multiplicity is $1$ at $n=5$ and $2$ from $n\ge6$, which is the parity-sensitive Lovelock collapse of Lemma~51 seen in the block count.

\begin{Lem}{59}[The certificate's necessary condition is one dimension count]
Let $K$ be the span cut out by the orbit constraint and let $\mathrm{ANN}$ be the invariant sextics $F$ with $\mathrm{Cat}_F(K,K)=0$. Apolarity gives $\langle F,\mu(B)\rangle=\langle\mathrm{Cat}_F,B\rangle$, so, for $\mu_K$ the restriction of $\mu$ to the invariants supported on $K$,
\[
\mathrm{rank}(\mu_K)=14-\dim\mathrm{ANN},\qquad F_5\in\mathrm{image}(\mu_K)\iff F_5\perp\mathrm{ANN}.
\]
$\mathrm{Cat}$ is injective on the invariant sextics (the unrestricted $\mu$ is surjective, so its adjoint is injective), placing a copy of the $14$ inside the $65$ invariant symmetric forms on $\mathrm{Sym}^3$. The forms vanishing on $K\times K$ number $29$, counted block by block: an odd-$\lambda_2$ block dies entirely, $K$ containing all of $M_\lambda$; an even one keeps $u\otimes v+v\otimes u$, $K$ having codimension $1$ there. So $65=36+29$ splits the invariant forms into those supported on $K$ (the unknowns) and those killing $K\times K$.
\end{Lem}

\begin{Lem}{62}[$\dim\mathrm{ANN}=1$, so the rank is $13$]
Order-$1$ vanishing at $W_0$ caps the restricted rank at $13$: a product of two cubics of $K$ vanishes at $W_0$, and evaluation there is a nonzero functional on the $14$. A sampled rank is bounded above by the true one, and reducing modulo $p$ can only lower it further (the $40$ evaluation points being exact rationals fixed before any prime is chosen, so what is reduced is the reduction of a rational matrix), so a sampled rank \emph{attaining} $13$ forces equality. It attains it, at four primes, and exactly over $\mathbb Q$. Hence $\dim\mathrm{ANN}=1$, spanned by $A(W_0^6)$, and $F_5\perp\mathrm{ANN}$.
\end{Lem}

\noindent The cap is conditional on the copies being identified equivariantly, and the multi-probe Schur test is what supplies that. A family identified non-equivariantly can sample \emph{above} its own rank.

\noindent\textbf{Every dimension of $\mathrm{ANN}$ is a coincidence}: a generic $14$ meets a generic $29$ inside a $65$ in $14+29-65=-22$ dimensions. \emph{The one obstruction guaranteed to exist is exactly the one $F_5$ passes.} $A(W_0^6)$'s catalecticant is the orbit \emph{average} of the rank-one forms $(f,g)\mapsto f(gW_0)g(gW_0)$, of rank $3570$. It kills $K\times K$ because every $f\in K$ vanishes on the whole orbit, not merely at $W_0$. The rank-one form itself is not $O(5)$-invariant and is the catalecticant of no invariant sextic.

\paragraph{The obstruction criterion without catalecticants.} $\mathrm{Sym}^6$ is spanned by sixth powers, so the invariant sextics are spanned by the projections $A(v^6)$. Apolarity makes $\mathrm{Cat}_{v^6}$ the rank-one form $(f,g)\mapsto f(v)g(v)$, and averaging over the group gives $\mathrm{Cat}_{A(v^6)}(f,g)=\int_{O(5)}f(gv)g(gv)$. Hence
\[
A(v^6)\in\mathrm{ANN}\iff v^3\in C,\qquad C=\mathrm{span}\{(gW_0)^3\},
\]
and the cone over $W_0$'s orbit contributes the one direction. $W_0$ is a smooth point of $V(K)$ ($\dim\operatorname{span}\{\nabla f(W_0):f\in K\}=28$, so the tangent space has dimension $7$, the orbit cone's), and $160$ ascents, forty of them over $\mathbb C$, find no $v$ with $v^3\in C$ off that cone.

\paragraph{Two scope points the machinery carries.} The unknowns are \emph{real symmetric}:
\[
\dim\mathrm{Sym}^2(\mathrm{Sym}^3\Weyl_5)^{O(5)}=65,\qquad \dim\Lambda^2(\mathrm{Sym}^3\Weyl_5)^{O(5)}=24,
\]
summing to the $89$ a Hermitian ansatz over $\mathbb Q(i)$ offers, so $24$ directions are gauge, though which $24$ is not the imaginary coordinates of that basis, the identification holding only in a basis real for the real structure, where $h$ would be real symmetric and is measured otherwise. And positive definiteness is read in the $G$ coordinates because the ansatz matrix meets them by a \emph{congruence} by the multiplicity Gram, which is positive definite. A congruence preserves definiteness in both directions.

\subsection{The certificate at $n\ge12$: the block structure is built}\label{sec:sos12}

Conjecture~A at $n\ge12$ needs the same construction one dimension up, and the block data is computed rather than quoted.

\begin{Lem}{63}[The $n=12$ block table]
$\dim\Weyl_{12}=1638$, the irreducible of highest weight $(2,2,0,0,0,0)$ for $D_6$. $\mathrm{Sym}^3(\Weyl_{12})$ carries $77$ $SO(12)$ constituents with $191$ multiplicities. Twelve are chiral, in six pairs with $\lambda_6\ne0$ at multiplicities $2,1,1,1,1,1$, and an orientation-reversing element fuses each pair into one $O(12)$ block. So $O(12)$ sees $\mathbf{71}$ blocks and $\mathbf{184}$ multiplicities, with $\sum m_\lambda d_\lambda=\binom{1640}3=733\,813\,080$.
\end{Lem}

\noindent Four independent routes agree on $1638$ (the Weyl dimension formula, Freudenthal, Littlewood restriction with Jacobi--Trudi, and Sage), and three on the decomposition, the third being Racah's alternating sum over $W(D_6)$, which reproduces all $191$ multiplicities and proves completeness. Six invariants separate the $71$ blocks, the sixth being $p_{12}$ or equivalently $\mathrm{Pf}^2$: the two are different polynomials, $\mathrm{Pf}^2$ vanishing where $p_{12}=261\,453\,379$, but generate the same extension of $\mathbb Q[p_2,\dots,p_{10}]$.

\paragraph{The table is stable from $n=12$, and $12$ is the threshold.} $\Weyl_n$ is the partition $(2,2)$, so a Newell--Littlewood product of three has at most six rows; a partition labels an $O(n)$ irreducible exactly when $\lambda'_1+\lambda'_2\le n$; and across the constituents $\max(\lambda'_1+\lambda'_2)=12$, at $(2,2,2,2,2,2)$. Koike--Terada carries this from labels to multiplicities. So no modification rule fires for $n\ge12$ and the decomposition cannot move.

\begin{certificate}[Four independent routes]
Littlewood's inverse restriction formula in the ring of symmetric functions gives $71/184/477$ and reproduces $64/182/493$ at $n=10$; the Weyl character formula block by block gives $70/183$ at $n=11$, $77/191$ at $n=12$ and $71/184$ at $n=13$ with zero mismatches; $\sum m_\lambda\dim_{O(n)}=\binom{\dim\Weyl_n+2}{3}$ holds at $n=12,13,14,16,18$ and \emph{fires} at $n=10$; and the same identity confirms the $n=4k$ boundary for $\mathrm{Sym}^2$ at $n=8$ and $\mathrm{Sym}^4$ at $n=16$. The boundary is pinned by the Weyl tensor itself: $\dim\Weyl_3=0$ where $\lambda'_1+\lambda'_2=n+1$, and $\dim\Weyl_4=10=n$.
\end{certificate}

\paragraph{What stability does \emph{not} give, and this is the report's own correction.} It fixes the \textbf{shape} of the certificate problem at every $n\ge12$: which blocks, which multiplicities, how many unknowns. It does \textbf{not} make the constraint data $n$-independent: the $89$ invariant sextics, the covariants and the orbit constraint are computed at a given $n$. \emph{So a certificate built at $n=12$ settles Conjecture~A at $n=12$, and transferring it upward is a separate question.}

\begin{Lem}{64}[The transfer mechanism]
Fix $d\ge1$ and let $F$ be an $O(n_0)$-invariant form of degree $d$ on $\Weyl_{n_0}$. For every $n\ge n_0$, $F$ has an $O(n)$-invariant extension to $\Weyl_n$. When $n\ge n_0\ge k_B(d)$, the presentation degree of $\{\mathrm{Sym}^d\Weyl_n^*\}$ over the signed permutation groups $\{B_n\}$, it is the unique $B_n$-invariant extension. The first fundamental theorem writes $F$ as a combination of complete contractions, each of which restricts correctly under padding because every slot sits on a Weyl factor, and Proposition~2.9 gives uniqueness. With $SO$ in place of $O$, uniqueness holds as stated and existence holds when $n_0$ is odd or $n_0\ge4d$.
\end{Lem}

\noindent\textbf{Equivariance was therefore never the obstruction.} What the mechanism needs is a value of $k_B(d)$, which is finite with none known, and $\{\Weyl_n\}$ is generated in degree exactly $4$ over $B_n$ with $\mathrm{Sym}^\ell$ in degree $4\ell$, so $k_B\ge24$ for the degree-six certificate. \emph{The dimension-free machinery of Levin--Chandrasekaran does not apply to $\{O_n\}$ at all}: it is defined only for finite-index subgroups, and $O(d)\times O(n-d)$ has infinite index, so no module over $\{O_n\}$ has a presentation degree while the theorems require one.

\paragraph{The certificate's own target is not freely described.} With $A=\lVert W\rVert^6$, $B=\langle Q(W),W\rangle^2$ and $C=\lVert W\rVert^2\lVert Q(W)\rVert^2$ (each the same contraction at every $n$), the target is $\Theta_n^2A-B+\mu(C-B)$, and its only $n$-dependence is that scalar. A restriction-compatible extension carries $\Theta_{n_0}^2$ unchanged, and $\Theta_n^2$ is strictly increasing, so the extension asserts $\theta\le\Theta_{n_0}$. \textbf{It is worse than false}: $C-B$ vanishes \emph{exactly} at the balanced product, where $Q(W)=\theta W$, so the extension is negative there for \emph{every} $\mu$ and the description set is \textbf{empty}. Writing $n$ as a contraction does not repair it: $I_n$ is freely described but its trace is not, $\mathrm{tr}(I_{n+1})=n+1$ while $\mathrm{tr}(P(I_{n+1}))=n$. \emph{Scope: $2-\Theta_n^2>0$ always, so the non-sharp target $2A-B+\mu(C-B)$ \emph{is} freely described. This bites the sharp constant alone.}

\paragraph{Owed, and it is the binding constraint.} The $184$ covariants are $7$ of $71$ blocks, and \textbf{none is built here}: the extraction produces vectors rather than diagrams, and three steps rest on their independence as functions. \emph{The cost is not what blocks it}: extraction by composition runs at roughly three units per dimension and grows linearly, so a sampling route is affordable. What is missing is the construction.

\paragraph{A certificate at $n=12$, built outside this project.} An outside project built the covariants and the certificate, and it replays here end to end.

\begin{Lem}{254}[Conjecture~A at $n=12$, on an outside certificate]
On $\Weyl_{12}$, with the full four-index norm,
\[
f_{12}(R)=\tfrac{55}{36}\lvert R\rvert^6-\langle Q(R),R\rangle^2
\]
is a sum of $170$ positive rational multiples of squared norms of $O(12)$-equivariant cubic maps, over $65$ of the $71$ blocks of Lemma~63. Hence $\langle Q(R),R\rangle\le\Theta_{12}\lvert R\rvert^3$ on $\Weyl_{12}$, attained at $S^6\times S^6$, so $M(12)=\Theta_{12}$.
\end{Lem}

\begin{certificate}
The identity is verified exactly, over $\mathbb Q$ and at two primes near $10^6$, against a $22{,}915\times477$ table of sextic invariants built by the outside project. It was replayed here from the parent archives with every validator passing, and on a second backend. \textbf{The qualification travels with every use of the lemma}: the translation from Gram coordinates to polynomials (the contraction seeds, the seminormal bases and the block scales $\kappa_\lambda$) is independently checked only for the target and for blocks of output rank at most $4$, and the outside project recorded and corrected two errors in exactly that layer. An independent computation of the $71$ scales $\kappa_\lambda$ is the check that would lift it. No reviewer verdict is on file.
\end{certificate}

The equality case is not claimed: that the balanced orbit is the only maximiser at $n=12$ does not follow, as it does not for Lemma~61. And the certificate is one a reader cannot check by hand, so the construction this subsection owes stands: a certificate at $n=12$ settles $n=12$, and the transfer upward is the question of Lemma~64.

\begin{Lem}{65}[The seed is the inclusion padded by cups]
Let $n\ge1$ and let $\lambda$ be an $O(n)$-admissible label with $|\lambda|\in\{0,2,\dots,12\}$. Write $k=(12-|\lambda|)/2$, let $\iota_\lambda\colon V_\lambda\hookrightarrow V^{\otimes|\lambda|}$ be the inclusion and $\mathrm{cup}=\sum_ae_a\otimes e_a$. Then $\iota_\lambda\otimes\mathrm{cup}^{\otimes k}$ is a nonzero element of $\mathrm{Hom}_{O(n)}(V_\lambda,V^{\otimes12})$, so the seed need not be searched for. Admissibility makes $V_\lambda$ nonzero and holds whenever $|\lambda|\le n$.
\end{Lem}

\begin{certificate}
Two reviewers, independently and on identical briefs, built $\iota_\lambda$ at non-trivial $\lambda$ rather than at the two smallest labels only, and the lemma holds at every block. At $16$ of the $71$ blocks $|\lambda|=12$, so $k=0$ and the seed is the inclusion $\iota_\lambda$ with no cups. The rank identity behind the count is $\dim\nu(B_f(n))=\sum_\lambda m_\lambda(n)^2$ over \emph{restricted} up--down tableaux, matching at every pair tested, and it is unconditional in $n$: no threshold enters.
\end{certificate}

\paragraph{The seed does not close the obligation, and the gap is a search.} The chain runs seed $\to$ a spanning set of $23!!=316{,}234{,}143{,}225$ elements $\to$ the \emph{existence} of a combination. Only the first link is supplied. \textbf{A spanning set of that size is not a construction}, and the deliverable is a readable proof rather than a certificate nobody can check.

\paragraph{Restitution accounts for none of the difficulty.} The classical name for the diagonal restriction is restitution, and Weyl settles its kernel, but polarization accounts for $0$ of the $10{,}393$ kernel dimensions, which are the non-symmetric complement. The three-slot rank is $2$ at $n=12$ as well, so restitution costs nothing and buys nothing. \emph{A gate that shipped alongside the earlier reading tested a scale-sensitive tolerance rather than the claim.}

\subsection{The stratum with at most two negative eigenvalues}\label{sec:stratum}

The sharp case of the residual lives where $S=0$ holds vacuously. Total antisymmetry makes $S$ a count of triples drawn from $\mm$, so $\dim\mm\le2$ forces $S=0$ with nothing to check, and the rival, with spectrum $(2,-1,-1)/\sqrt6$, sits exactly there. Bounding $\theta$ on that stratum therefore excludes the nearest competitor with no enumeration at all. It is settled at one positive eigenvalue and open beyond.

\begin{Lem}{29}[One positive and two negative eigenvalues cap $\theta$ at $\thstar$]
Dimension-free, and attained at $(2,-1,-1)/\sqrt6$, i.e.\ at padded $W_{\CP^2}$.
\end{Lem}

\begin{proof}[sketch]
Three of the four sign-blocks of the cubic vanish structurally. The pair cap bounds the fourth, and trace-freeness collapses the estimate to a single closed form with one maximum.
\end{proof}

Lemma~29 covers one positive eigenvalue, and the descent that would reduce the general case to dimension four is what Lemma~34 refutes. The stratum closes directly instead, at any number of positives and any dimension.

\begin{Lem}{35}[At most two negative eigenvalues caps $\theta$ at $\thstar$]
Let $W\in\Weyl_m$ have unit norm and at most two negative eigenvalues, at any number of positive ones and any $m$. Then $\theta\le\thstar$. Content and sharpness both begin at $m=4$.
\end{Lem}

\begin{proof}[sketch]
In an orthonormal eigenbasis, writing $T_{\alpha\beta\gamma}=\langle[v_\alpha,v_\beta],v_\gamma\rangle^2$,
\[
\theta=\sum_\alpha\lambda_\alpha^3+3\!\!\sum_{\alpha<\beta<\gamma}\!\!T_{\alpha\beta\gamma}\lambda_\alpha\lambda_\beta\lambda_\gamma,
\]
an identity that needs no Weyl hypothesis, though it does need $\lVert W\rVert=1$, the display being $P(W)$. Relax $T$ to any nonnegative totally symmetric tensor with row marginals $m-2$ and pair caps $\sum_\gamma T_{\alpha\beta\gamma}\le2$: the problem becomes a transportation relaxation. Dropping the row marginals leaves a \emph{pair-capped triangle packing} on the support alone (free of $m$, which is what makes the bound dimension-free), and the sharp harmonic-mean dual (Lemma~42) bounds that by $\thstar$. The closing one-variable inequality is machine-checked in Lean: in closed form on $(0,\tfrac13]$, by an exact factorisation with a double root at the endpoint $\sqrt{2/3}$, and by rational interval arithmetic in the kernel between; validated interval arithmetic at $200$ bits is an independent check. The lemma as a whole, from the eigenframe to the bound and for at most two negative eigenvalues including one and none, is machine-checked in Lean.
\end{proof}

\paragraph{The dual step of that chain is an argument, not a computation.} It is linear-programming weak duality, and the whole of it is four lines. The first interval and the endpoint band close symbolically, the endpoint through the factorisation, and every endpoint constant reproduces by hand; the middle interval closes by rational interval arithmetic that the Lean kernel carries out, which is mechanical rather than readable. The argument as written covers exactly two negative eigenvalues. The one-negative case is the boundary of that family.

\begin{Lem}{42}[The sharp dual weight]
Let $\lambda_P$ be at least three positive eigenvalues and let
\[
M(\lambda_P)\;=\;\max\Big\{\textstyle\sum_{\alpha<\beta<\gamma}t_{\alpha\beta\gamma}\lambda_\alpha\lambda_\beta\lambda_\gamma\;:\;t\ge0,\ \sum_{\gamma\notin\{\alpha,\beta\}}t_{\alpha\beta\gamma}\le2\ \text{for every pair}\Big\}.
\]
Then $3M(\lambda_P)\le h\,(s^2-u)$, with $h$ the harmonic mean of the three largest positives, $s=\sum\lambda$ and $u=\sum\lambda^2$. \textbf{The constant cannot be improved.}
\end{Lem}

\begin{proof}
The dual is $\min\,2\sum_{\alpha<\beta}z_{\alpha\beta}$ subject to $z_{\alpha\beta}+z_{\beta\gamma}+z_{\alpha\gamma}\ge\lambda_\alpha\lambda_\beta\lambda_\gamma$ on every triple, $z\ge0$. Try the product weight $z_{\alpha\beta}=c\,\lambda_\alpha\lambda_\beta$. Dividing a triple constraint by $\lambda_\alpha\lambda_\beta\lambda_\gamma>0$ turns it into $c\big(\lambda_\alpha^{-1}+\lambda_\beta^{-1}+\lambda_\gamma^{-1}\big)\ge1$, that is $c\ge\tfrac13H$ with $H$ the harmonic mean of that triple, since $\sum\lambda^{-1}=3/H$. The harmonic mean is increasing in each argument, so the binding triple is the three largest positives and $c=h/3$ is feasible. Weak duality then gives $M\le2\sum_{\alpha<\beta}z_{\alpha\beta}=2\cdot\tfrac h3\cdot\tfrac{s^2-u}2$.

For sharpness, take $k$ equal positives $\lambda$: then $h=\lambda$ and $s^2-u=k(k-1)\lambda^2$, so the bound reads $3M\le k(k-1)\lambda^3$. It is attained: $t_{\alpha\beta\gamma}=2/(k-2)$ on every triple has pair marginal $(k-2)\cdot\tfrac2{k-2}=2$, hence is feasible, and $\sum t=\binom k3\tfrac2{k-2}=k(k-1)/3$.
\end{proof}

\begin{remarkstar}[Two constants, each easy to state wrongly]
The feasibility threshold is $\mu\ge2/(3\sum\lambda^{-1})$, not $\mu\ge2/\sum\lambda^{-1}$: under the normalisation $z=\mu\cdot\tfrac32\lambda_\alpha\lambda_\beta$ the multiplier is $c=\tfrac32\mu$, so the shorter form is off by a factor $3$. \textbf{The prose gloss that replaced it was worse}: ``at least the harmonic mean of the triple'' is $\tfrac92$ times the true threshold, and taking $\mu=h$ gives $3M\le\tfrac92h(s^2-u)$, four and a half times weaker than the lemma. The lemma itself is unaffected by either: what establishes it is the argument above. The exhibited maximiser is also a repair: the first version of the sharpness paragraph derived the same \emph{upper} bound twice and never exhibited a feasible $t$, which both reviewers caught independently.
\end{remarkstar}

\paragraph{What this settles.} \textbf{The sharp stratum of the residual is excluded at every even $n\ge12$}, without enumeration and without the descent Lemma~34 refuted, since $\thstar<\theta_0(n)$ from $n=12$. Two independent audits cover it: one on the reduction and the packing bound, one on the residual inequality, where interval arithmetic replaces the grid certificate with a proof.

The remaining machinery is what turns a small eigenspace into a small support, and it is Lie-theoretic rather than analytic.

\begin{Lem}{30}[A small negative eigenspace forces a small support]
Ricci-flatness is the positive-semidefinite identity $\sum_\alpha\lambda_\alpha v_\alpha^{\mathsf T}v_\alpha=0$, so every eigenvector \textbf{with nonzero eigenvalue} has its image in $U=\operatorname{im}v_{N-1}+\operatorname{im}v_N$. When $\dim U<n$ the tensor is a padding of a $\Weyl_{\dim U}$ one, in a single step.
\end{Lem}

\begin{Lem}{31}[The bracket cap, characterised]
For orthonormal $X,Y\in\so(n)$, $\|[X,Y]\|^2=2$ exactly when they anticommute, both have rank $4$ with every singular value $1/\sqrt2$, and their images agree. The root system gives $\max_Y\|[X,Y]\|^2=(s_1+s_2)^2$, and Cauchy--Schwarz forces the rest.
\end{Lem}

\begin{Lem}{32}[A bivector's centraliser collapses to a line]
For $m\ne3$ and skew $x$ with $\mathfrak c=\mathfrak z(x)\cap x^\perp$, one has $\mathfrak z(\mathfrak c)\cap\mathfrak c^\perp=\R x$, unless $x$ is a multiple of a complex structure on a $4$-plane, where it is an $\su(2)$. Three root-system lines, uniform in the singular values.
\end{Lem}

\begin{Lem}{33}[One eigenvalue of one sign forces the tensor to vanish]
If $W\in\Weyl_m$ has a \textbf{one-dimensional} eigenspace on one side (exactly one eigenvalue counted \emph{with multiplicity}), and that eigenvector is not a complex-structure multiple on a $4$-plane, then $W=0$.
\end{Lem}

The lever is a Weitzenb\"ock-type identity: $\sum_\alpha\lambda_\alpha N_{v_\alpha}=0$ with $N_v=2\,v\otimes v-\operatorname{ad}(v)^2\succeq0$, whose per-eigenvector splitting turns the statement into the centraliser containment of Lemma~32. \textbf{This answers the question item B0 had been reduced to}: no rank-$3$ element at even $m\ge6$ has a full-rank minority eigenvector, so the object that $520$ earlier ascent starts had hunted does not exist. The padded $\Weyl_4$ family is the only boundary case.

\paragraph{And the reduction those lemmas were serving is refuted.}

\begin{Lem}{34}[Full support with two negative eigenvalues, at $m=6$ and $m=8$]
Two padded signature-$(2,1)$ $\Weyl_4$ tensors on $4$-planes sum to a $\Weyl_m$ element with $\dim U=m$. In general position they reach $\theta=+0.485$.
\end{Lem}

So ``no element with $\dim\mm\le2$ has $\dim U=m$'' is \textbf{false}, and the descent through Lemma~30 does not close this stratum. What the refutation cost was the \emph{route}, not the statement: Lemma~35 proves the same bound directly, and the measured $\sup\theta=\thstar$ at $m=6,8,10,12$ is now a theorem rather than an observation.

\paragraph{Nothing is owed here any more.} What was owed (the bound with two or more positive eigenvalues alongside at most two negative) is Lemma~35.


\begin{Lem}{72}[The rank-three stratum, closed in three lines for every $n$]
Let $W$ be a critical point of $\Weyl_n$ with exactly one positive and two negative eigenvalues. Then its spectrum on the support is $(2,-1,-1)/\sqrt6$ and $\theta(W)=\sqrt{3/2}$.
\end{Lem}

\begin{proof}
Write the spectrum $(\lambda_1,-a,-b,0,\dots)$ with $a,b>0$. Since $T$ vanishes on repeated indices the only live entry is $T:=T_{1,N-1,N}$, and the criticality rows $(W^{\#})_{\beta\beta}=\lambda_\beta(\theta-\lambda_\beta)$ at $\beta=1,N-1,N$ read
\[
Tab=\lambda_1(\theta-\lambda_1),\qquad T\lambda_1b=a(\theta+a),\qquad T\lambda_1a=b(\theta+b),
\]
three expressions for one number. Equating the last two gives $a^2(\theta+a)=b^2(\theta+b)$, so either $a=b$ or $\theta=-(a^2+ab+b^2)/(a+b)$. On the second branch $\theta+a=-b^2/(a+b)$ and the middle equation returns $T=-ab/\lambda_1^{2}<0$, impossible for $T=c^2$, while $T=0$ fails it outright. Hence $a=b$, and $\tr W=0$ gives $\lambda_1=2a$, the first two equations give $\theta=3a$ and $T=2$, and $\|W\|=1$ fixes $6a^2=1$.
\end{proof}

\paragraph{Remark.}
The stratum is \emph{overdetermined}: the rows give $r$ equations against $\binom r3$ live entries, so $T$ is overdetermined at rank $3$, determined at rank $4$, and free from rank $5$ on. That is why this stratum closes and no other does. The same rows also settle what was recorded as a conjectured bridge across $0<\delta<0.005$: the per-row Bianchi balance $\sum_{\alpha\ne\beta}\lambda_\alpha\epsilon_{\alpha\beta}=0$ forces all three pair caps to saturate, and rows $N-1$ and $N$ added give the bridge as an identity.

\subsection{What $S=0$ costs the negative eigenspace, and what it does not buy}\label{sec:horns}

Lemma~35 closes one end of the residual. The other end turns out to be closed too, and by an unexpected route: the slack condition is a statement in extremal graph theory in disguise.

\begin{Lem}{36}[$S=0$ caps the negative eigenspace at $\lfloor n^2/4\rfloor$]
Slack zero says the negative eigenspace $\mm$ is \emph{isotropic} for the Cartan $3$-form $\langle[x,y],z\rangle$ on $\so(n)$, and such a subspace has dimension at most $\lfloor n^2/4\rfloor$, attained by the balanced symmetric pair.
\end{Lem}

This is a continuous analogue of \textbf{Mantel's theorem} (a triangle-free graph on $n$ vertices has at most $\lfloor n^2/4\rfloor$ edges, with the balanced complete bipartite graph extremal), and the correspondence is exact rather than decorative: $S=0$ reads as triangle-freeness of the negative graph, and the extremal configuration is the balanced split in both settings. Two reviewers prove the cap independently, by different routes. Rigidity at the cap holds for even $n$ and is open for odd $n$.

\paragraph{Both ends of the stratum are settled, and the middle was read as a slope between them. That reading is refuted.} The residual is stratified by $\dim\mm$:

\begin{center}
\begin{tabular}{@{}p{0.26\linewidth}p{0.32\linewidth}p{0.34\linewidth}@{}}
\toprule
\textbf{Stratum} & \textbf{Value there} & \textbf{Status} \\
\midrule
$\dim\mm\le2$ & $\thstar=1.2247$, attained at padded $W_{\CP^2}$ & \textbf{Closed} by Lemma~35, at every $m$ \\
\addlinespace
$3\le\dim\mm<\lfloor n^2/4\rfloor$ & reaches $1.228239$ at $S^5\times S^7$ and $1.324617$ at $S^{11}\times S^{13}$ & \textbf{$\theta\le\thstar$ here is FALSE} \\
\addlinespace
$\dim\mm=\lfloor n^2/4\rfloor$ & $\theta_0(n)$, the conjectured maximizer itself & \textbf{Closed} by Lemma~36, with rigidity at even $n$ \\
\bottomrule
\end{tabular}
\end{center}

\paragraph{Why the two closed ends do not make the middle low.} The natural reading of a stratum capped at both ends (\emph{two horns and a low middle}, with $\theta\le\thstar$ on $\{S=0,\ 3\le\dim\mm<\lfloor n^2/4\rfloor\}$) is false, and $S^5\times S^7$ at $n=12$ settles it alone: $\theta=1.228239158>\thstar$, $\dim\mm=35<36$, $S_{\mm\mm\mm}=0$ exactly, criticality residual $4\times10^{-16}$. \textbf{No floating point is load-bearing in that witness}: built over $\mathbb Q$, $Q(W)$ is exactly diagonal with $Q(W)_{aa}/W_{aa}=\tfrac{11}{35}$ at all $66$ entries, so $\theta^2=\tfrac{264}{175}$ and the excess over $\tfrac32$ is exactly $\tfrac3{350}$. Between $25$ and $28$ catalogued band points exceed $\thstar$ at $n\le24$.

\textbf{And the band is asymptotically extremal rather than low.} Along $S^{n/2-1}\times S^{n/2+1}$ the ratio $\theta/\theta_0(n)$ runs $0.9937$, $0.9976$, $0.9988$, $0.9997$, $0.99996$ at $n=12,16,20,30,60$, with $\theta^2\to2$: there is no gap between the middle of the band and the cap in either coordinate. The stratum is inhabited by named points, not only by search output: $\CP^m$ lies in it for every $m\ge3$, provably, since $\dim\mm=m(m-1)$ satisfies $3\le m(m-1)<m^2$.

\textbf{The missing hypothesis is diagonality.} Every witness above is a \emph{diagonal} point. Slack zero plus a dimension bracket does not constrain $\theta$, whereas slack zero plus non-diagonality is the statement of \S\ref{sec:outside}. Lemmas 35 and 36 are unaffected: each caps its own end of the stratum, and neither says anything about the middle.

\paragraph{The middle cannot inherit either end's argument either, and that much is proved rather than suspected.} Two obstructions are on record.

\emph{The cap's rigidity does not propagate downward.} Maximal isotropic subspaces of the Cartan $3$-form occur at \emph{many} dimensions: at $n=6$ all $211$ edge-maximal triangle-free graphs give maximal isotropics, at dimensions $5,7,8,9$, and generic isotropics at $d=5,6,7$ are maximal too. So the $3$-form is unlike a quadratic form, and the middle stratum is not made of proper subspaces of a symmetric pair.

\emph{And Lemma~35's machinery provably cannot be reused.} It would be reusable if $S=0$ were inert in that relaxation: $S=0$ is exactly $T=0$ on $(-\,-\,-)$ triples, whose products are negative, so a maximiser looks as though it empties them anyway. \textbf{That ignores the row marginals.} At the relaxation's own maximising spectrum at $n=6$ the $(-\,-\,-)$ mass is $12.000000$, and imposing $S=0$ renders that spectrum \emph{infeasible}: with most indices negative, the marginal $\sum_{\beta\gamma}T_{\alpha\beta\gamma}=2(n-2)$ forces mass into $(-\,-\,-)$ because a negative index has nowhere else to put it. The relaxation is blind to $S=0$ only where $S=0$ is slack. What survives is the operational content and it is unaffected: with the vector identity \emph{and} $S=0$ both imposed the optimum still reaches $1.306498$ against $\thstar$, down from $1.5305$ on marginals alone but over. \textbf{So the next argument must use $S=0$ through something other than the bracket tensor's marginals}, and Lemma~36 is the first channel found that does.

\subsection{Capping the non-diagonal families analytically}\label{sec:families}

Enumeration was the wrong target and a value bound is the right one, but a value bound over an uncountable set still has to be assembled from somewhere. The route that has moved is to \textbf{cap each family by a closed form in its own parameter}, which is a proof on that family rather than a data point.

\begin{Lem}{38}[Two families capped, in closed form and strictly decreasing]
For the $\CP^m$ family,
\[
\theta^2\;=\;\frac{(m+1)(2m-1)}{3m(m-1)},
\]
strictly decreasing in $m$ with maximum $\tfrac32$ at $m=2$. For the K\"ahler family of the symmetric pair $(\su(m),\so(m))$ inside $\so(2m)$, with $p=\dim\kk=\binom m2$ and $q=\dim\mm=\binom{m+1}2-1$,
\[
\theta^2\;=\;\frac{m(m+1)}{(m-1)(m+2)}\;=\;\frac X{X-2},\qquad X=m(m+1),
\]
again strictly decreasing, with maximum $\tfrac32$ at $m=2$. \textbf{The $m=2$ member of the second family is $W_{\CP^2}$ itself.}
\end{Lem}

\begin{proof}[sketch, second family]
Take the two-eigenvalue graded ansatz $W=a\Pi_{\kk}+b\Pi_{\mm}$ with $ap+bq=0$ and $a^2p+b^2q=1$. Two of the four block masses vanish identically from the symmetric-pair relations $[\kk,\kk]\subseteq\kk\perp\mm$ and $[\mm,\mm]\subseteq\kk\perp\mm$: $S_{\kk\kk\mm}=0$ and $S_{\mm\mm\mm}=0$. The second is \emph{why} slack zero holds along the whole family. The survivors have closed forms, $S_{\kk\kk\kk}=(m-2)p$, the Killing form of $\so(m)$, and $S_{\kk\mm\mm}=mq$. Then $\theta=\tr(W^3)+\tfrac12\big[a^3S_{\kk\kk\kk}+3ab^2S_{\kk\mm\mm}\big]$, and eliminating $b=-ap/q$ and $a^2=q/(p(m^2-1))$ gives the stated form identically in $m$. Monotonicity is immediate from $X/(X-2)$.
\end{proof}

\paragraph{The equality case stops being a coincidence.} $W_{\CP^2}$ is the smallest member of an infinite family along which $\theta$ decreases monotonically to $1$. \textbf{That is a structural statement about the conjectured boundary point, and it is what this attack buys.} What it does not buy is a cap: the family's own supremum is $\thstar$, attained, so the closed form can never bound the class \emph{strictly} below the target. \emph{The cap is $\thstar$ and not $\sqrt{6/5}$}, which is what the family gives from $m=3$ up. The $m=2$ member is the point that sets the threshold, and a reading that starts the family at $m=3$ omits it.

\begin{Lem}{39}[The two-eigenvalue formula for an arbitrary symmetric pair]
Nothing in the derivation above used $\su(m)$. For \emph{any} symmetric pair $(\mathfrak g,\kk)$ inside $\so(n)$ with $p=\dim\kk$, $q=\dim\mm$, the two-eigenvalue point has $S_{\kk\kk\mm}=S_{\mm\mm\mm}=0$ identically and
\[
\theta^2\;=\;\frac{\big(S_{\kk\kk\kk}q^2+3S_{\kk\mm\mm}p^2-2p^3+2pq^2\big)^2}{4p^3q(p+q)^3}.
\]
When $\mathfrak g$ is simple with Killing constant $c$, $S_{\kk\kk\kk}+S_{\kk\mm\mm}=cp$, so $\theta$ depends on four invariants only: $p$, $q$, $S_{\kk\kk\kk}$, $c$. The constant is classical table data, $c=2h^\vee(\mathfrak g)/I_\rho$, dual Coxeter number over Dynkin index.
\end{Lem}

\begin{certificate}
Confirmed on $16$ symmetric pairs across five Cartan series. The Killing constant was checked intrinsically as $-B/\langle\cdot,\cdot\rangle$ at four independent points: $\mathfrak g_2\subset\so(7)$ returns $8$; $\so(3)$ on $\mathrm{Sym}^2_0(\R^3)$, where the index is far from $1$, returns exactly $2/5$; every adjoint representation returns $2$. \emph{A normalisation trap, recorded because it is easy to fall into}: the indices here are \emph{half} the standard Dynkin index of the complexified representation, so applying the rule with indices taken from the usual tables gives every $c$ wrong by a factor $2$. And $S_{\kk\kk\kk}=c_{\kk}p$ fails whenever $\kk$ is not simple (a centre contributes nothing, and every Hermitian symmetric pair has one), the correct form being $S_{\kk\kk\kk}=\sum_ic_i\dim\kk_i$ over the simple ideals.
\end{certificate}

\paragraph{The reduction this invited is dead in principle, and that is worth more than a failed attempt.} Four invariants and Cartan's finite list read as a finite arithmetic problem. It is not one, for three separate reasons, all established at Gate~2.

\begin{enumerate}[leftmargin=1.8em,nosep]
\item \textbf{Admissibility is provably not a function of the four invariants.} $(\so(6),\so(2)+\so(4))$ and $(\su(4),\mathfrak s(\uu(2)+\uu(2)))$ share $(p,q,S_{\kk\kk\kk},S_{\kk\mm\mm})=(7,8,24,32)$ and share $\theta^2=9245/8232$, yet fail differently: one on Ricci-flatness, one on Bianchi. So no function of those invariants decides whether the ansatz is a Weyl tensor at all.
\item \textbf{The representation is a free infinite parameter.} $(\su(2),\uu(1))$ on one, two and three copies of $\mathbb C^2$ keeps $(p,q)=(1,2)$ while $\theta^2$ runs $3/2$, $2/3$, $25/54$.
\item \textbf{Semisimple $\mathfrak g$ with non-isomorphic factors has no single $c$.} $\so(3)+\so(5)$ inside $\so(8)$ gives $c\in\{2,6\}$ and the Killing relation fails outright.
\end{enumerate}

Of $16$ pairs built, only \emph{three} pass Ricci-flatness, Bianchi and criticality together. \textbf{So the real content of ``which pairs threaten $\thstar$'' is which $(\mathfrak g,\kk,\rho)$ give a Weyl operator at all}, a question the reduction never posed, and the unrestricted maximum over the class is $\sqrt2$, not $\thstar$. \emph{Scope, stated plainly}: all of this is the \textbf{two-eigenvalue} ansatz. $\CP^m$ and $\HP^2$ have three distinct eigenvalues and are not of this form. What covers them, and every eigenvalue count at once, is the quadratic reformulation of Lemma~46.

\subsection{The $\lambda$-weighted bound, and the one ingredient it lacks}\label{sec:lambdabound}

The second attack is an inequality of the right shape (an upper bound on $\theta$ that \emph{decreases} in the ungradedness), and it is the one such bound the programme has. Two reviewers derived it independently, and it appears in none of the literature searched.

\begin{Lem}{41}[A $\lambda$-weighted identity, and the bound it gives]
Write $g_\gamma=\lVert[\mathrm{ad}(v_\gamma),W]\rVert^2$ in an eigenbasis of $W$. Then
\[
\sum_\gamma\lambda_\gamma g_\gamma\;=\;-4\,\theta(W),
\]
in two lines from $A(W)=2W$ and the definition of $\#$. With $g_\gamma\ge0$ and $\sum_\gamma g_\gamma=4(n-1)$ this gives
\[
\theta\;\le\;|\lambda_{\min}|\cdot\frac{4(n-1)-\Delta-G_0}{4},\qquad G_0:=\sum_{\gamma\in\ker W}g_\gamma ,
\]
with equality exactly when $g$ is supported on the $\lambda_{\min}$ eigenspace.
\end{Lem}


\begin{certificate}
The identity holds to $12$ digits at $n=5,6,7,8,10$ and at residual $\le4\times10^{-15}$ over $14$ catalogued points. The bound is never violated, worst margin $7\times10^{-16}$. \textbf{Equality holds at all nine symmetric-space points} (the four sphere products, $\CP^{2,3,4}$ and both padded points) and \textbf{at the balanced product the bound returns $1.236033=\theta_0(12)$ exactly, so it is sharp at the conjectured maximizer}. Slack appears only where $\Delta>0$: $0.671$ against $0.894$ at the $C_5$ blow-up, and $1.2$--$1.4$ at random Weyl tensors.
\end{certificate}

\begin{Lem}{66}[No maximizer is $\circ$-decomposable]
Call $W$ \emph{$\circ$-decomposable} if $W=\sum_{i=1}^ka_iA_i$ with $k\ge2$, each $A_i$ a unit element of $\mathrm{Weyl}_n$, $\sum_ia_i^2=1$, every $a_i\ne0$, and $A_i\circ A_j=0$ for $i\ne j$. Then for $n\ge4$ no maximizer of $P$ on the unit sphere of $\mathrm{Weyl}_n$ is $\circ$-decomposable.
\end{Lem}

\begin{proof}
$A_i\circ A_j=0$ makes $P$ additive: $P(W)=\sum_ia_i^3P(A_i)$. Each $A_i$ is a unit element of $\mathrm{Weyl}_n$ and $P$ is odd, so $\lvert P(A_i)\rvert\le M(n)$, and $M(n)>0$ at $n\ge4$ because $S^k\times S^l$ with $k,l\ge2$ has $P=\theta_{k,l}>0$ (Nienhaus, Proposition~2.25). Hence
\[
P(W)\;=\;\sum_ia_i^3P(A_i)\;\le\;M(n)\sum_i\lvert a_i\rvert^3\;\le\;M(n)\max_i\lvert a_i\rvert\;<\;M(n),
\]
using $\sum_i\lvert a_i\rvert^3\le\max_i\lvert a_i\rvert\sum_ia_i^2=\max_i\lvert a_i\rvert$, and $\max_i\lvert a_i\rvert<1$ because $\sum a_i^2=1$ with $k\ge2$ and every $a_i\ne0$.
\end{proof}

The absolute values are load-bearing: with $a_i^3$ in place of $\lvert a_i\rvert^3$ the chain fails at any negative coefficient. The pieces' support dimensions never enter.


\paragraph{What blocks it, and what has been ruled out as a source.} $|\lambda_{\min}|$ is uncontrolled, so the question is two-variable in $(\Delta,\lambda_{\min})$, and at random points the bound reaches $1.367$ at $n=8$, above $\thstar$. Three sources were assessed for the missing ingredient and none supplies it as an eigenvalue estimate: sharp Weitzenb\"ock/PIC2 estimates (fence post~21: blind to $\lambda_{\min}$, and $2.5$--$7.8\times$ short quantitatively), the curvature operator of the second kind, whose sphere theorems bound the \emph{smallest} eigenvalues from \emph{below} and so run in the wrong direction, and the $\uu(m)$-isotypic decomposition of Lemma~28 (exact but $\CP^m$-specific). \textbf{The second of those is not closed as a source of structure, only as a source of this estimate}: the operator carries half of $\theta$ exactly (Lemma~50), which is \S\ref{sec:bispectral}. \textbf{The untried lever that is named rather than hoped for is whether $|\lambda_{\min}|$ can be traded against $\Delta$}, since the bound needs only the \emph{product} controlled.

\begin{remarkstar}[$\theta$ is a named invariant, and its constant is convention-sensitive]
$\Theta_3(W)=-2\,\theta(W)$, the stable cubic of Weingart's graphical calculus, proportional at every $n$ with a universal constant and no Ricci correction, verified independently twice. \textbf{The constant is $-2$ under Definition 6.1's normalisation}, and it moves to $-8$, $-32$ or $+32$ under a dropped factor $\tfrac14$ or a flipped bracket sign. The magnitude is convention-dependent, but the sign is not, since $\Theta_3$ is cubic. That matters beyond bookkeeping: fence post~19's conclusion turns on this sign, so the convention in use must always be named. The identification opened Ge's sharp-Weitzenb\"ock literature to the search, where $q_2(W)=-2W$ on Weyl tensors was verified at $10^{-16}$ (their operator is ours), and it is fence post~21 that closes the route regardless.
\end{remarkstar}

\paragraph{A per-direction cap sharpens the bound pointwise and cannot close it.} Substituting Petersen--Wink's Proposition 1.7 into a \emph{single} $g_\gamma$ rather than into the $\lambda$-weighted sum, and using $T\ge0$ with the row marginal $\rho_\alpha=2(n-2)$, gives
\[
g_\gamma\;\le\;2(n-2)\,(\lambda_{\max}-\lambda_{\min})^2 .
\]
This sits \emph{outside} $\{g\ge0,\ \sum g=4(n-1)\}$, which is the only information Lemma~41 uses. Since $\sum_{\mm}g=4(n-1)-\Delta-G_0$ exactly and the cap is smaller than $4(n-1)$, mass is forced off the bottom eigenvalue and water-filling turns it into a valid sharpening: $15$--$25\%$ against the a priori bound, $2$--$12\%$ against the $\Delta$-aware one, and enough to put every random-point maximum tested below $\thstar$.

\textbf{It costs nothing at the maximizer, and that is the obstruction rather than the recommendation.} The cap is saturated \emph{exactly} at $S^k\times S^k$ for $k=3,\dots,10$ (slack $0.0000$), so the capped bound returns $\theta_0(2k)$ to machine precision and no sharpening of it can gain anything where the problem is hard. All the slack sits at random Weyl tensors, which are not the live band. \emph{And the class is already closed}: the cap is built from $T\ge0$ and one marginal, so it lives inside fence post~10's transportation class, ceilinged at $1.60$/$2.29$/$2.49$ at $n=6,8,12$.

\textbf{Two-moment refinements are strictly stronger and change nothing where it matters.} Expanding $g_\gamma=2\sum_\alpha\lambda_\alpha^2R_{\gamma\alpha}-2Q_\gamma$ with $Q_\gamma=\sum_{\alpha\beta}\lambda_\alpha\lambda_\beta T_{\gamma\alpha\beta}$ and using the per-row identity $\sum_\beta\lambda_\beta R_{\gamma\beta}=-2\lambda_\gamma$ as a constraint gives a genuinely $\gamma$-dependent cap. Bounding the two terms separately cancels the gain exactly, $\mathrm{Rf}_\gamma-U=(a-b)(\lambda_{\max}^2-\lambda_{\min}^2)$, whereas spending both moments jointly in one LP does not, beating both caps by $17$--$33\%$ and dominating the uniform cap everywhere: ratios $0.857$ at $S^3\times S^5$, $0.917$ at $S^4\times S^6$, $0.945$ at $S^5\times S^7$. \emph{The gain is exactly $0$ at every balanced product}, and the LP is inside post~10's class again.

\textbf{So the term to attack is $Q_\gamma$, and its equalities are definitional.} The two exact constraints $\sum_\gamma Q_\gamma=-2\|W\|^2$ and $\sum_\gamma\lambda_\gamma Q_\gamma=2\langle W^{\#},W\rangle$ hold at \emph{random} Weyl tensors too, so criticality contributes nothing and the second is the definition of $\#$ in the eigenbasis. Contracted against weight $1$ and weight $\lambda$ they return $\sum g=4(n-1)$ and $\sum\lambda g=-4\theta$, both already held. The one Weyl-specific input is the per-row identity, and it is weaker than it looks: $\Pi_{\mathcal W}(\mathrm{ad}(v)^2)=-2\,\Pi_{\mathcal W}(vv^{\mathsf T})$ holds for \emph{every} unit $v\in\so(n)$, eigenvector or not, so it is one linear identity restating trace-freeness. \emph{And the natural inequality on $Q_\gamma$ is fenced}: $Q_\gamma=\lambda^{\mathsf T}A\lambda$ with $A=T_{\gamma\cdot\cdot}\ge0$ symmetric with prescribed marginals, and minimizing that is a transportation problem, post~10 exactly. \textbf{This exhausts the per-direction cap from inside that structure. Whether it exhausts the $\lambda$-weighted attack as a whole is a route-status claim and is not made here.}

\subsection{The three relations on the bracket data, and what they cannot decide}\label{sec:pinned}

Write $T_{abc}=\langle[v_a,v_b],v_c\rangle^2$ in an eigenframe of $W$: nonnegative, totally symmetric, vanishing when two indices coincide. Three relations hold per index, and they are the whole of what the Killing form, Ricci-flatness with first Bianchi, and criticality give:
\[
\text{(A)}\ \sum_{b,c}T_{abc}=2(n-2),\qquad
\text{(B)}\ \sum_{b,c}\lambda_bT_{abc}=-2\lambda_a,\qquad
\text{(C)}\ \sum_{b,c}\lambda_b\lambda_cT_{abc}=2\lambda_a(\theta-\lambda_a).
\]
(A) and (B) are B\"ohm--Wilking's Lemma 2.1 read as $I\#I=(n-2)I$ and $W\#I=-W$. (C) is the diagonal of $W^{\#}=\theta W-W^2$. \textbf{The list is complete}: no fourth per-index relation exists, the four candidate sources (first Bianchi, the square-producing contractions of Jacobi, the off-diagonal criticality block, and the quartic Casimir) being empty by termwise vanishing, by averaging back to (C), and by circularity through $Q$ preserving $\Weyl_n$.

Row $a$ of (A) makes $\mu_a:=T_{a\bullet\bullet}/2(n-2)$ a probability measure on pairs, and (B), (C) become expectations against it. Since an expectation lies in the range of what it averages, each row yields a necessary condition, and the two extreme rows recover both of this programme's standing spectral bounds:

\begin{Lem}{73}[A spectral necessary condition on all of $\Weyl_n$]
For every $W\in\Weyl_n$, critical or not, the spectrum satisfies $\lambda_{\max}\le(n-2)\lvert\lambda_{\min}\rvert$ and $\lvert\lambda_{\min}\rvert\le(n-2)\lambda_{\max}$. The constant $n-2$ is not improvable: it is attained at $\CP^m$ for every $m$, where $\lambda_{\max}/\lvert\lambda_{\min}\rvert=2(m-1)=n-2$, and at every maximal torus.
\end{Lem}

\begin{proof}
Row $a$ of (B) reads $\mathbb E_{\mu_a}\big[\tfrac{\lambda_b+\lambda_c}2\big]=-\lambda_a/(n-2)$, and the averaged quantity lies in $[\lambda_{\min},\lambda_{\max}]$, so $-\lambda_a/(n-2)$ does too, for every $a$. Only (A) and (B) are used, and neither involves $\theta$. The saturating cases have $E_{\max}$ abelian, so $[E_{\max},E_{\max}]=0$ and $[E_{\max},\cdot]\subseteq E_{\min}$.
\end{proof}

\paragraph{Remark.}
Tracelessness alone gives only $N-1$ in place of $n-2$, for every $n\ge3$. The condition is slack at every sphere product, so it cuts away from the extremum rather than toward it.

Applying the same calculus to (C) at the $\lambda_{\min}$ row gives $\theta\le\lambda_{\min}+(n-2)\lambda_{\max}$, and at the $\lambda_{\max}$ row, using $\lambda_b\lambda_c\le\big(\tfrac{\lambda_b+\lambda_c}2\big)^2$ and $\mathbb E[X^2]\le(\lambda_{\min}+\lambda_{\max})\mathbb E[X]-\lambda_{\min}\lambda_{\max}$, gives $\theta\le(n-1)\lvert\lambda_{\min}\rvert$. \textbf{Both standing bounds are one system read at opposite ends}, neither an independent result.

\begin{Lem}{70}[The deficiency identity]
At a critical point of unit norm, for every unit $v\in E_{\min}(W)$,
$\theta=\lambda_{\min}+(n-2)\lambda_{\max}-D(v)/\lvert\lambda_{\min}\rvert$ with $D(v)=\tfrac12\sum_{a,b}(\lambda_a\lambda_b-\lambda_{\min}\lambda_{\max})c_{vab}^2\ge0$, the nonnegativity needing $\tr W=0$, and the identity itself needing $\lVert W\rVert=1$, because its left side is scale-invariant and its right homogeneous of degree one. Equality holds exactly when $\ad(v)$ is supported on $E_{\min}\otimes E_{\max}$, and it forces $\dim E_{\min}\ge n/2$ and $\dim E_{\max}\ge(n-2)/2$, both attained at $W_{\CP^2}$ at $n=4$.
\end{Lem}

\begin{Lem}{71}[Saturation forces the balanced product]
Let $n\ge5$ and let $W$ be a critical point saturating Lemma~70. Then $\so(n)=\hh\oplus\mm$ is a symmetric pair with $\hh=[\mm,\mm]$, $W$ has exactly two eigenvalues, and $W$ is the Weyl tensor of $S^{n/2}\times S^{n/2}$. Consequently $\theta=\theta_0(n)$, and no critical point saturates at odd $n\ge5$.
\end{Lem}

\begin{proof}
Saturation gives $[\mm,\mm]\subseteq\hh$, $[\mm,\hh]\subseteq\mm$ and $[\mm,\rr]=0$ for $\rr$ the remaining eigenspaces. Then $\mathfrak{g}'=\mm+[\mm,\mm]$ is a subalgebra, $\rr$ centralises it, and $\hh$ normalises it since $[\hh,[\mm,\mm]]\subseteq[[\hh,\mm],\mm]+[\mm,[\hh,\mm]]\subseteq[\mm,\mm]$. So $\mathfrak{g}'$ is an ideal, and $\so(n)$ is simple for $n\ge5$ and $\mm\ne0$, whence $\mathfrak{g}'=\so(n)$, $[\mm,\mm]=\hh$ and $\rr=0$. Writing $p=\tfrac n2+t$, $\dim(\so(p)\oplus\so(q))=\tfrac{n^2}4+t^2-\tfrac n2$, so the required $\dim\hh=\tfrac{n(n-2)}4$ forces $t=0$. $\mathfrak u(n/2)$ has the wrong dimension, and $\dim\mathfrak{g}^\theta$ is an $\Aut$-invariant so triality adds none at $n=8$. Finally $\dim\mm=n^2/4$ must be an integer.
\end{proof}

\paragraph{Remark.}
$n=4$ is the sole exception, $\so(4)$ being the only non-simple $\so(n)$, and $W_{\CP^2}$ realises it: $\mathfrak{g}'$ is one $\so(3)$ factor and $\rr$ the other. Lemma~71 \emph{classifies} rather than bounds: its hypothesis is the scalar equation $\theta=\lambda_{\min}+(n-2)\lambda_{\max}$, so verifying it presupposes $\theta$, and at a non-saturating point its only content is the already-weak Lemma~70.

\paragraph{What the three relations cannot decide.} Maximising $\theta$ over $T\ge0$ subject to (A), (B), (C) with the spectrum free returns $1.340542$ at $n=6$, against a true maximum of $\sqrt{3/2}=1.224745$. Relation (B) is what closes most of the gap (the same programme without it returns $1.590$, and the marginal relaxation of post~10 returns $1.606$), \textbf{and the class still fails}. The failure has a published shape: in the analogous nilsoliton problem the linear system in the squared structure constants admits a positive-dimensional family of nonnegative solutions on which the Jacobi identity selects a single point, and Jacobi is quadratic and sign-dependent in those constants, hence not a function of $T$ at all. Two non-isometric algebras in dimension~$6$ share their $T$. \textbf{The binding condition is orthonormal-\emph{frame} realizability, not the Jacobi identity}: in $\so(n)$ Jacobi is automatic, the structure constants being $\so(n)$'s own rather than unknowns. What $T$ fails to determine is the algebra itself, two non-isometric algebras in dimension~$6$ sharing their $T$, and the first free contraction $\sum\lambda_a^2\lambda_b^2T_{abc}$ is no way around it: at a two-eigenvalue spectrum $\lambda^2$ is affine in $\lambda$, and the balanced product has two eigenvalues, so that contraction is spectrum-determined exactly where the conjecture is decided.

\subsection{The diagonal half has left curvature entirely}\label{sec:diagonal}

The other half of the residual (diagonal $S=0$ points other than the balanced product) has turned into a problem with no curvature in it. On the diagonal locus $\Weyl_n$ is $\{w$ symmetric, hollow, with vanishing row sums$\}$, and
\[
\theta(w)=\sum_{a<b}w_{ab}^3+3\!\!\sum_{\text{triangles}}\!\!w_{ab}w_{bc}w_{ca},
\]
under zero row sums and unit norm, with $S=0$ reading as \emph{triangle-freeness of the negative graph} and the balanced product as the complete-bipartite weight. Equivalently $\theta=\tfrac12\big(\sum_{a,b}M_{ab}^3+\tr M^3\big)$ for symmetric $M$ with $M\mathbf1=0$, $\tr M=0$, $\tr M^2=2$.

\begin{Lem}{68}[The symmetry defect on the diagonal locus, in closed form]
On the diagonal locus, with $M$ the curvature matrix of $W$,
\[
g_{ab}=2\|M_a-M_b\|^2-4M_{ab}^2,
\]
so $\Delta=0$ exactly when the positive graph's components are twin cliques with non-positive off-block entries. At a critical point the entries obey $g_{ab}=2(s_a+s_b)-4\theta M_{ab}$, an equality at every twin pair, and $\theta\|M\|_1=\sum_ad_a^+s_a-\tfrac12\Delta$: defect converted into a $\theta$ deficit, exactly there.
\end{Lem}

The deficit does not close the half on its own: the entry cap is an equality at every balanced product, so no entry-free relaxation of it can separate them. A falsifier hunt over triangle-free seed graphs reproduces the recorded ceilings exactly at $n=10$ and $n=12$ and raises the $n=12$ ceiling to $0.716134$, a factor $1.71$ below the target.

\paragraph{Why the two halves cannot be separated.} The spectral half alone caps $\tfrac12\tr M^3$ at $1.2309$ against $\theta_0(12)=1.2360$ (\textbf{$99.6\%$ of the target}), and the balanced product attains that spectral cap \emph{exactly}, from $M\mathbf1=0$ and the two trace conditions alone, with no graph condition and no $S=0$. So what is owed is one trade-off, $E-E_0\le T_{\max}-T$, between the entrywise and spectral shares. Both carry the same constant $c_n=(n-3)/\sqrt{(n-1)(n-2)}$, and bounding them separately cannot work.

\paragraph{Criticality compresses to an $n\times n$ equation.} The multiplier vanishes and
\[
M^{\circ2}+M^2=\theta M+\operatorname{diag}(s_i),\qquad s_i=\sum_jM_{ij}^2 ,
\]
a complete characterisation of diagonal critical points, verified at $10^{-16}$. The vertex system of \S\ref{sec:gradedhalf} is its block-constant restriction. Criticality is load-bearing rather than decorative: over $15{,}000$ random admissible $M$ the residual quantity reaches $3.50$ against a target of $2$, so a proof must feed in this equation and not merely the linear constraints.

\paragraph{The residual as a Poincar\'e inequality.} Contracting at the top eigenvector reduces the statement on the vertex-transitive locus to the single inequality
\[
\sum_{i,j}M_{ij}^2(u_i-u_j)^2\;\ge\;\frac{2\mu_1^2}{n},
\]
the Dirichlet energy of $M$'s top eigenvector against the weights $M_{ij}^2$, with \textbf{equality at both sphere products}, to nine digits, and strict at the $C_5$ blow-up. Its tersest form is $\mu_1(M)^2+2\nu_2(\tfrac n2M^{\circ2})\le\tr M^2$, again with equality exactly at the sphere products.

\paragraph{The rigidity is an association scheme, and on its two-class case the half is settled.} Feeding criticality back in gives $M^{\circ2}=\theta M-M^2+\tfrac2nI$ on the vertex-transitive locus: \emph{the Hadamard square is a quadratic polynomial in $M$}, which is exactly the rigidity behind strongly regular graphs. The algebra $\mathrm{span}\{I,M,\dots\}$ contains $J$ and is Hadamard-closed at $10^{-15}$ at both balanced products (two-class, strongly regular) and at the $C_5$ blow-up (three-class), while the non-vertex-transitive $S^4\times S^8$ fails closure at $0.18$.

\textbf{On the two-class locus the diagonal half is settled}: $\theta$ is an explicit function of the parameters $(n,k,\lambda,\mu)$, and maximising over every feasible parameter set returns exactly $\theta_0(n)$ at $K_{n/2,n/2}$, at every even $n$ from $6$ to $24$. The enumeration runs over a \emph{superset} of realisable strongly regular graphs, so it holds a fortiori. And the two halves meet: the spectral cap is attained only by three distinct eigenvalues, which under the rigidity \emph{is} the two-class case, so cap-attaining implies closed, and otherwise there is a strict spectral deficit to spend against $E-E_0$.

\begin{Obs}{144}[No two-class point lies in the band except the balanced product]
Every strongly regular graph gives a critical point of the diagonal locus, with $\theta(\Gamma)=\lvert1+\lambda-\mu\rvert\sqrt{2(n-1)/(nk(n-1-k))}$. The imprimitive graphs give $\theta=\sqrt{2(n-1)(g-1)/(n(n-g))}$, \emph{maximised uniquely at $g=n/2$}, which is the balanced product. The primitive ones satisfy $\theta<\tfrac12$ over the range examined.
\end{Obs}

\begin{certificate}
The imprimitive half is proved. The primitive half is a census: every parameter set to $n\le200$, maximum $0.4758$. The non-sphere supremum is $0.7059\,\Theta_v$ over $v\le300$. \emph{The label is Observation because a uniform bound over all primitive parameter sets is owed.}
\end{certificate}

\emph{Granting the screened half, the falsifier a two-class point could have supplied cannot exist.}

\paragraph{Owed.} The branch with three or more classes (the $C_5$ blow-up sits there, at $0.6708$) and the non-vertex-transitive points, where the rigidity fails outright.

\subsection{$P$ as a Rayleigh quotient}\label{sec:rayleigh}

The newest reformulation is one line, and it may be the whole conjecture in disguise. In an eigenbasis of $W$, with the structure constants transported there,
\[
P=\sum_\alpha\lambda_\alpha^3+\tfrac12\sum c'^2_{\alpha\beta\gamma}\lambda_\alpha\lambda_\beta\lambda_\gamma=\big\langle\lambda,M(\lambda)\lambda\big\rangle,
\qquad
M(\lambda)_{\alpha\beta}=\delta_{\alpha\beta}\lambda_\alpha+\tfrac12\sum_\gamma c'^2_{\alpha\beta\gamma}\lambda_\gamma ,
\]
with $M$ symmetric because $c'^2$ is symmetric in its first two indices. Since $\|\lambda\|=1$ this gives $\theta\le\lambda_{\max}(M(\lambda))$, and \textbf{at a critical point $\theta$ is an eigenvalue of $M$ rather than merely bounded by one}, since $Q(W)=\theta W$ is diagonal in the eigenbasis, so $M\lambda=\theta\lambda$ exactly.

The cubic has become a quadratic form whose matrix depends on the point, and the identity is exact. \textbf{But it is a restatement rather than a route, and the reason is worth recording as a rule.} Writing $\langle x,Mx\rangle=\langle L_WX,X\rangle$ for $X$ diagonal in $W$'s eigenframe shows $\lambda_{\max}(M)$ is the top eigenvalue of $L_W$ on the eigenframe-diagonal algebra $\mathcal D$.

\textbf{Lemma~24 does not reach that algebra, and the equality needs a separate argument.} Its cap lives on $\Weyl_n\cap W^\perp$, while $\mathcal D$'s elements are generically neither trace-free nor Bianchi ($\tr(v_\alpha\otimes v_\alpha)=1$ where every Weyl operator is trace-free), and at a generic $W$ in $n=6$ the intersection $\mathcal D\cap\Weyl_6$ is $\R W$ alone, so everything beyond the free direction sits outside Lemma~24's reach. \textbf{On the diagonal locus the equality is proved.} Coordinate-diagonal operators are hollow symmetric $S$, Weyl exactly when $S\mathbf 1=0$, so the non-Weyl complement is $\{S_{ij}=u_i+u_j\}$, with no cross term between the two parts, and there the Rayleigh quotient is at most $\tfrac12\lambda_{\max}(\Lambda)<1/\sqrt2$, $\Lambda$ the hollow matrix of $W$'s entries. Hence, with $M$ formed in the coordinate frame, $\lambda_{\max}(M)=\theta$ at every coordinate-diagonal global maximizer, where Lemma~24 caps the Weyl part and $\theta\ge\thstar>1/\sqrt2$, and at every $S^m\times S^m$ with $m\ge2$, where a Cauchy--Schwarz bound caps the Weyl part with no maximality assumed. \emph{In general the equality is owed a cap at level $\theta$ on the part of $\mathcal D$ orthogonal to $\Weyl_n$, which is nowhere on record, and at repeated eigenvalues $M$ depends on the eigenframe.}

\textbf{And $\sup_W\lambda_{\max}(M)\le\theta_0(n)$ is not stronger than Conjecture~A.} The record measures $\sup\lambda_{\max}(M)=M(n)$ at $n=6,8,10$, so the two suprema agree rather than one dominating; at odd $n\ge13$ the target is $\Theta_n<\theta_0(n)$, which makes the $\theta_0$ form \emph{weaker}; and below $n=12$ it is simply false, $M(n)\ge\sqrt{3/2}>\theta_0(n)$ there. \emph{The pointwise statement at maximizers is the conjecture itself, which is the honest reading and the one the remark below states as a rule.}

\begin{remarkstar}[A bound tight at the extremum is the original problem]
Replacing a hard maximisation by a pointwise upper bound looks like progress the moment the bound is easier to write down, and the question of how much was given away goes unasked. \textbf{The cheap test is one evaluation}: if the bound is tight at the known or conjectured extremiser, its global supremum equals the original one, and proving it is at least as hard. A reformulation can still earn its place by connecting the problem to machinery that already exists, but that is a different claim, and it has to be made separately.
\end{remarkstar}

Retired as a shortcut. Its identities stand, and the ratio $\lambda_{\max}(M)/\theta=1.206011$ at the $C_5$ blow-up matches the recorded $\lambda_{\max}(L_e|_{W^\perp})$ there.

% ==============================================================================

% ==============================================================================
